% Created 2026-09-04 Fri 12:28
% Intended LaTeX compiler: pdflatex
\documentclass[preprint,12pt]{elsarticle}
\usepackage{graphicx}
\usepackage{rotating}
\usepackage[normalem]{ulem}
\usepackage{amsmath,amssymb,amsthm}
\usepackage{booktabs}
\usepackage[final]{microtype}
\usepackage{xurl}
\usepackage{hyperref}
\usepackage{orcidlink}
\usepackage{listings}
\usepackage{xcolor}
\usepackage{array}
\newcolumntype{L}[1]{>{\raggedright\arraybackslash}p{#1}}
\usepackage{placeins}
\usepackage{algorithm}
\usepackage{algpseudocode}
\usepackage{tikz}
\usetikzlibrary{shapes.geometric,arrows.meta,positioning,calc}
\usepackage{pgfplots}
\usepgfplotslibrary{groupplots}
\pgfplotsset{compat=1.17}
% The seams-core release that the paper cites and that the program files hold: the main-branch commit it is built on, its source tree and program-files archive, and the Zenodo concept DOI of the seams-core GitHub releases with its latest version.
% No commit holds the tree of the release; it is the base commit with the ten patches of the program files.
% Set the version DOI in \seamscoredoi once a release with that tree is archived.
% The reproducibility package is cited at its commit (\reprocommit, tag v2.7.0).
% The 31 October records used the first revision of the package patch, sha256 0e53bc8e, tree 5dde84a5.
% The revision cited here, sha256 f1721ac7, tree 5722048b, adds the PTM cutoff 0.1 and scripts/ptm_cutoff_sweep.py, fits q6 at the mW coexistence point, reads the dump column header, and writes the freud means with scripts/sota_freud.py.
% It rewrites the Lechner-Dellago curves of Figure sota with scripts/regen_ld_figure.py.
% That script takes its sigmas and seeds from the command line.
% Set \reprodoi once that tree is archived.
% snapper:off
\RequirePackage{seqsplit}
\newcommand*{\seamscoretag}{v2.12.0}
\newcommand*{\seamscoretree}{598e964e}
\newcommand*{\seamscorebase}{305fb3c67c395303df6d9171a0c1af3520061625}
\newcommand*{\seamscoresha}{1f5b88b4956c9a7c46bceab0150650f5b9ca6b469b5641e774a52a06b93b7352}
\newcommand*{\seamscoredoi}{}
\newcommand*{\seamscoreconceptdoi}{10.5281/zenodo.3874291}
\newcommand*{\seamscoreconceptlatest}{v2.11.0}
\newcommand*{\reprocommit}{f625883238bf0fb3b5b56b61a48994e4cb12e43e}
\newcommand*{\reprobase}{f625883}
\newcommand*{\reprotree}{5722048b}
\newcommand*{\repropatch}{dseams2\_repro-v2.7.0-fixes.patch}
\newcommand*{\repropatchsha}{f1721ac779fe171705156c8845ed8d971bf89812eec64db053deead80b76d57a}
\newcommand*{\reprodoi}{}
\makeatletter
\edef\repro@doiurl{https://doi.org/\reprodoi}
\newcommand{\reproarchive}{%
  \ifx\reprodoi\@empty
    commit \texttt{\expandafter\seqsplit\expandafter{\reprocommit}} (tag v2.7.0)
    with the package patch \texttt{\repropatch} of the supplementary data,
    sha256 \texttt{\expandafter\seqsplit\expandafter{\repropatchsha}}, which gives git tree \texttt{\reprotree}%
  \else Zenodo, \expandafter\url\expandafter{\repro@doiurl}\fi}
\newcommand{\reproid}{%
  \ifx\reprodoi\@empty commit \texttt{\reprobase} (tag v2.7.0) with the package patch of the supplementary data, git tree \texttt{\reprotree}%
  \else \expandafter\url\expandafter{\repro@doiurl}\fi}
\edef\core@doiurl{https://doi.org/\seamscoredoi}
\edef\core@concepturl{https://doi.org/\seamscoreconceptdoi}
\newcommand{\coretag}{\seamscoretag}
\newcommand{\coretree}{\texttt{\seamscoretree}}
\newcommand{\corehowpublished}{%
  \url{https://github.com/d-SEAMS/seams-core}, commit
  \texttt{\expandafter\seqsplit\expandafter{\seamscorebase}} with the ten
  release patches of the program files, tree \coretree}
\newcommand{\corearchive}{%
  \ifx\seamscoredoi\@empty
    Program files archive \texttt{seams-core-\seamscoretag.tar.gz},
    sha256 \texttt{\expandafter\seqsplit\expandafter{\seamscoresha}}.
    Zenodo concept DOI \expandafter\url\expandafter{\core@concepturl},
    latest version \seamscoreconceptlatest%
  \else Zenodo, \expandafter\url\expandafter{\core@doiurl}\fi}
\makeatother
% snapper:on
% The tag of the seams-core release that contains the measured fix to the cage labels.
% Set it once that release is made; until then both documents print the placeholder.
% snapper:off
\newcommand*{\releasetag}{[to be added]}
\newcommand{\DOIPENDING}{pending}
% snapper:on

% References between the manuscript and the supplement.
% \sref{label} in rg_main.tex prints the number of a label of rg_supplement.tex, and \mref{label} in rg_supplement.tex the number of a label of rg_main.tex.
% scripts/mkxref.py writes supp_xref.tex and main_xref.tex from the .aux of each document, so each compiles alone; scripts/build_pdfs.sh fails when either file is stale.
% A missing label prints ?? and warns like \ref.
% snapper:off
\makeatletter
\newcommand{\xr@def}[3]{\expandafter\def\csname xr@#1@#2\endcsname{#3}}
\newcommand{\xr@get}[2]{%
  \ifcsname xr@#1@#2\endcsname\csname xr@#1@#2\endcsname
  \else\textbf{??}\PackageWarning{xrefs}{Reference `#2' on page \thepage\space undefined}\fi}
\newcommand{\sref}[1]{\xr@get{s}{#1}}
\newcommand{\mref}[1]{\xr@get{m}{#1}}
\makeatother
\InputIfFileExists{supp_xref}{}{}
\InputIfFileExists{main_xref}{}{}
% snapper:on

\definecolor{coral}{RGB}{255,127,80}
\definecolor{teal}{RGB}{0,128,128}
\lstset{basicstyle=\ttfamily\small,breaklines=true}
\lstdefinelanguage{lua}{morekeywords={local,function,for,do,end,require,in,if,then,else,return},sensitive=true,morecomment=[l]{--},morestring=[b]"}
\renewcommand{\thealgorithm}{S\arabic{algorithm}}
\renewcommand{\thefigure}{S\arabic{figure}}
\renewcommand{\thetable}{S\arabic{table}}
\newtheorem{theorem}{Theorem}
\newtheorem{lemma}[theorem]{Lemma}
\newtheorem{proposition}[theorem]{Proposition}
\newtheorem{corollary}[theorem]{Corollary}
\renewcommand{\thetheorem}{S\arabic{theorem}}
\newtheorem{paperprop}{Proposition}
\renewcommand{\thepaperprop}{P\arabic{paperprop}}
\newtheorem{papercor}{Corollary}
\renewcommand{\thepapercor}{C\arabic{papercor}}
\newcommand{\leanname}[2]{\par\noindent{\normalfont\small Lean: \texttt{DSeams.\detokenize{#1}}, file \texttt{#2.lean}}}
\journal{Computer Physics Communications}
\biboptions{sort&compress}
\makeatletter\def\ps@pprintTitle{\let\@oddhead\@empty\let\@evenhead\@empty\let\@oddfoot\@empty\let\@evenfoot\@oddfoot}\makeatother
\renewcommand{\thepage}{S\arabic{page}}
\renewcommand{\theequation}{S\arabic{equation}}
\renewcommand{\thesection}{S\arabic{section}}
\date{\today}
\hypersetup{
 pdfauthor={Rohit Goswami, Amrita Goswami, Ruhila Goswami, Jayant K. Singh},
 pdftitle={Supplementary material for d-SEAMS 2.0},
 pdfkeywords={},
 pdfsubject={},
 pdfcreator={},
 pdflang={English}}
\setlength{\emergencystretch}{2em}
\begin{document}

\begin{frontmatter}
\title{Supplementary material for ``d-SEAMS 2.0: Incremental ring and cage algorithms for ice structure analysis''}
\author[iitk]{Rohit Goswami\fnref{rg-now}\orcidlink{0000-0002-2393-8056}}
\ead{rgoswami@ieee.org}
\author[iitk]{Amrita Goswami\corref{cor1}\fnref{ag-now}\orcidlink{0000-0001-8706-2383}}
\ead{amrita@hi.is}
\author[hi-life]{Ruhila Goswami\orcidlink{0000-0002-5443-9356}}
\ead{rug17@hi.is}
\author[iitk]{Jayant K. Singh\corref{cor2}\orcidlink{0000-0001-8056-2115}}
\ead{jayantks@iitk.ac.in}
\cortext[cor1]{Corresponding author}
\cortext[cor2]{Co-corresponding author}
\fntext[rg-now]{Present address: SURF B.V., Amsterdam, The Netherlands, and TurtleTech ehf., Reykjavik, Iceland.}
\fntext[ag-now]{Present address: Universidad Complutense de Madrid, Madrid, Spain.}
\address[iitk]{Department of Chemical Engineering, Indian Institute of Technology Kanpur, Kanpur 208016, Uttar Pradesh, India}
\address[hi-life]{Faculty of Life and Environmental Sciences, University of Iceland, Reykjavik, Iceland}

\begin{abstract}
This supplement holds the material that the main text points to.
Sections~\ref{sec:s-background}--\ref{sec:s-defs} give the background,
the frame loops of other codes, the neighbour graphs, the bond-orientational order parameters and notes
on the ring and cage definitions.
Sections~\ref{sec:s-seeded}--\ref{sec:s-invariance} describe the seeded
label, ring completion, cage signatures, topology keys, the numerical
techniques, the device batch, the costs in full and the cage-object
order.
We prove Propositions~P1--P5 on paper in
Sections~\ref{sec:org0f685be}--\ref{sec:org5328df8}.
Algorithm~\mref{alg:rings} emits every primitive ring once, a source of
Algorithm~\mref{alg:rings} reads only a bounded ball,
Algorithm~\mref{alg:hopbound} equals Algorithm~\mref{alg:rings} on the
next frame, and the completion loop of Algorithm~\ref{alg:seeded}
returns the least fixed point.
Theorems~S1--S34 state 34 of the 53 public theorems of a Lean~4 development
\citep{moura2021lean4}, and each names its declaration and states what
that declaration proves.
The Lean statements are about abstract lists, predicates and natural
numbers (Section~\ref{sec:lean}) and are not a proof about the C++
implementation.
Sections~\ref{sec:s-accuracy}--\ref{sec:s-repro} hold the
label-accuracy tests, shape matching, the verification suite with the
corrections that change results, and the code overview with the carried
command line.
They then give the further benchmarks with the machines, the records and
the head-to-head calls, the applications and the commands that rerun
every timing record.
Primitive rings are those of Franzblau and Yuan
\citep{franzblau1993computation,yuan2002rings}, and the cage
predicates are those of d-SEAMS \citep{goswami2020dseams}.
\end{abstract}
\end{frontmatter}
\section{Background}
\label{sec:s-background}
Molecular-dynamics trajectories resolve ice nucleation from the first freezing event \citep{matsumoto2002molecular} through the homogeneous rates obtained by seeding \citep{sanz2013seeding,espinosa2014rates,espinosa2016seeding,goswami2020seeding} and by the shear-nucleation programme of \citet{goswami2021shearprl,goswami2021shearedreview}.
Those trajectories, however, record particle positions and do not track crystal structures or defects \citep{goswami2020dseams}.
Nucleation of soft matter is further complicated by competing ice polymorphs with small free-energy differences \citep{lupi2017stacking,russo2014anomalies}, while thermal fluctuations distort the ice-like structures that form and locally disrupt long-range order.
Surface interactions and confinement add further structure-determination cost.
Automated classification is therefore a prerequisite for interpreting those trajectories.

Water exhibits rich phase behaviour in bulk and in confinement.
At least seventeen bulk ice polymorphs have been observed experimentally \citep{bartels2012ice,salzmann2019phase}, and water confined within nanometer length scales forms ordered hydrogen-bond networks of ice nanotubes, monolayers, bilayers and trilayers \citep{koga2001nanotube,goswami2020nanoribbon,goswami2020hpice,suraj2021salt}.
Rigid non-polarizable models remain the standard tools of those simulations \citep{vega2011models,abascal2005tip4p2005,abascal2005tip4pice}, and the metastable liquid of supercooled water is itself a long-standing structural problem \citep{poole1992phase,debenedetti2001supercooled,palmer2014llt}.
Aqueous ions add a second length scale.
The solvation shell of Fe\(^{3+}\) switches between octahedral and capped trigonal prism states \citep{goswami2024fe3}, and the viscosity of those solutions tracks complex formation \citep{goswami2026viscosity}.
Those polymorphs, and the ordered structures that form during a nucleation event, must be identified before a trajectory can be analysed \citep{goswami2020dseams}.

Classification involves scoring local neighborhoods of molecules.
The Steinhardt bond-orientational parameters \citep{steinhardt1983bond} and their neighbour-averaged forms \citep{lechner2008accurate} are cheap to evaluate \citep{winczewski2016bop}, yet they are known to be deficient on disordered particulate matter \citep{mickel2013shortcomings}.
Smoothed Overlap of Atomic Positions (SOAP) environments \citep{bartok2013representing,caro2019optimizing}, the descriptor libraries that include them \citep{himanen2020dscribe}, and rotationally invariant hydrate order parameters \citep{zeron2024hydrate} extend that neighbourhood description.
Analysis suites \citep{ramasubramani2020freud,menon2019pyscal,michaud2011mdanalysis,gowers2016mdanalysis} expose those scores, as do collective-variable engines written for the same trajectories \citep{tribello2014plumed,tribello2019plumed}.
None of them name the cages of a growing nucleus.
CHILL and CHILL+ reduce the four-neighbour star to staggered and eclipsed bonds \citep{moore2010chill,nguyen2015identification}.
Learned identifiers treat ice versus liquid \citep{fulford2019deepice} or agreement with CHILL+ \citep{maity2025icecoder}; they do not return a hexagonal or double-diamond cage.

A nucleation study assigns one number per molecule.
CHILL+ is available in OVITO \citep{stukowski2010visualization} and in in-house code; freud and pyscal evaluate the averaged Steinhardt parameters \citep{ramasubramani2020freud,menon2019pyscal}.
A PLUMED bias requires a derivative, so the collective variables are Steinhardt scalars or the environment similarity of \citet{piaggi2022ice}.
The nucleus is then the largest cluster of molecules over a threshold.
That coordinate carries history dependence \citep{montero2026nonmarkovian}.
Those tools do not return a cage, a ring shared between two cages, or a prism.

Hexagonal cages and double-diamond cages are the units of ice Ih and ice Ic.
These form order parameters in a homogeneous nucleation-rate calculation \citep{haji2015nucleation}.
d-SEAMS \citep{goswami2020dseams} is the engine with which we turned cage membership into a published per-molecule ice assignment.
A molecule that belongs to a hexagonal cage (HC) is ice Ih, and a molecule that belongs to a double-diamond cage (DDC) is ice Ic.
A molecule that belongs to neither is liquid water.
We keep that assignment as the ice score, together with the CHILL and CHILL+ labels and a count of clathrate cages following \citet{jacobson2010amorphous}.
Confined ice is a separate problem.
Visual inspection of the hydrogen-bond network does not scale to a trajectory.
The topological network criteria of \citet{goswami2020nanoribbon} classify quasi-one- and quasi-two-dimensional ices without that inspection, and the \(n\) -gonal prisms of prismatic ice are among them.
The same programme describes high-pressure ices and salt water in diamond confinement \citep{goswami2020hpice,suraj2021salt}.
The largest connected component of cage-labelled molecules is the cluster size of the seeding study of ice nucleation under shear \citep{goswami2020seeding,goswami2021shearprl,goswami2021shearedreview}.
A hybrid topological and shape-matching formalism can map a non-convex cage onto its template \citep{goswami2021tum}.
We applied that mapping to heterogeneous nucleation on silver-exposed \(\beta\)-AgI \citep{vonnegut1947agi,zielke2015agi,zielke2016plates,prerna2019ice,goswami2020dseams}, to homogeneous bulk nucleation \citep{matsumoto2002molecular,haji2015nucleation,goswami2020dseams}, to monolayer square ice \citep{algarasiller2015squareice,goswami2020nanoribbon}, and to the freezing of an ice nanotube \citep{koga2001nanotube,goswami2020nanoribbon}.

Those local scores do not return a cage.
For the cage searches of d-SEAMS 1 we enumerated closed walks and applied Franzblau's shortest-path criterion afterwards \citep{franzblau1993computation}.
On disordered networks most generated walks are then discarded.
Hexagonal cages were then found by scanning every pair of surviving six-rings, and that pair scan dominates the cost on nucleation-scale cells.
Independently, the same searches read a cutoff neighbour graph.
On an ordered ice lattice the first coordination shell lies inside the \(3.5\,\mathring{\mathrm{A}}\) clustering cutoff of \citet{nguyen2015identification} and the second shell lies outside it, so the cutoff graph coincides with the four nearest neighbours of each molecule.
Under thermal motion first-shell neighbours can cross that cutoff, and the same cage tests then lose edges.

Herein, we describe d-SEAMS 2.0, which keeps the HC and DDC predicates of that assignment and replaces the searches that evaluate them on nucleation-scale cells.
Primitive rings follow the shortest-path construction of \citet{yuan2002rings}; an incremental updater returns the same output as a full rebuild.
An inverted atom-to-rings index makes both cage searches linear in the ring count at fixed degree.
A new labelling rule on the four nearest neighbours restores first-shell edges that a cutoff loses, and it labels nothing on a frame whose mutual four-nearest graph closes no cage.
The engine is a C++20 library with a command-line interface.
It replaces the single YAML--Lua executable of the original release.
We release a Python package and a Lua library separately as front ends to it.
We measure lattice-label retention against CHILL+, the averaged Steinhardt parameters and polyhedral template matching.
The test frames are jittered lattices, a random packing, a supercooled liquid and held-out mW trajectories of ice Ic, ice Ih and liquid water.
We also time the computations the engine shares with OVITO, freud, GenIce2 and RINGS on the same frames (Section~\mref{sec:comparison}).
Section~\mref{sec:introduction} lists what is new against d-SEAMS 1 and against those codes.
We follow the largest ice cluster through every public nucleation trajectory of \citet{niu2019temperature} and score CHILL+ on the same frames.
The remaining new applications are a seeded-growth campaign in mW, a Madrid-2019 ice/brine front and the occupancy of ideal GenIce methane hydrates.
We rerun the five deposits of the original release through the same library.
The rerun includes the confined-ice criteria that no per-molecule label provides.

d-SEAMS 2.0 keeps the HC and DDC predicates of d-SEAMS 1 \citep{goswami2020dseams} on a given graph and changes how they are computed; the seeded and completed labels are a new rule built on them.
The name d-SEAMS 2.0 denotes the program generation; the release described here is seams-core \coretag.
The shortest-path ring construction \citep{yuan2002rings} and inverted indices are established techniques.
What we add is their use in d-SEAMS, the exact source update and the seeded completion.
Table~\ref{tab:complexity} gives the cost of each stage in both versions, and Table~\ref{tab:features} places the program beside the codes a study would otherwise combine.
We list the new parts against d-SEAMS 1 below.
\begin{enumerate}
\item Primitive rings from pairs of shortest paths (Algorithm~\mref{alg:rings}). Each source sweeps its ball of radius \(\ell=\lfloor D/2\rfloor\) instead of growing every walk of length below \(D\), decides primitivity from distance rows of radius \(\ell'=\max(\ell-1,1)\), and emits every primitive ring once (Proposition~P1).
\item An exact incremental ring update between frames (Algorithm~\mref{alg:hopbound}). Only the rings that meet the set \(T\) of vertices within \(r\) hops of a changed bond on \(G\cup G'\) are enumerated again, with \(r=0\) up to seven members, and the result equals a rebuild (Proposition~P4).
\item An inverted atom-to-rings index for the hexagonal and double-diamond cage tests (Algorithm~\mref{alg:cages}). The HC test was an \(O(R_6^2)\) pair scan and drops to \(O\bigl(R_6\,\Delta\,d_I(\Delta+\log(\Delta d_I))+B\,d_I\bigr)\). The DDC test drops to \(O(R_6\,d_I^6\log(d_I+1))\). Both are linear in the ring count at fixed degree, and both cage searches visit rings in an order keyed to ring content.
\item A seeded hysteresis on the mutual and union four-nearest graphs, completed over the six-rings to a least fixed point (Algorithm~\ref{alg:seeded}, Proposition~P5). A frame whose mutual graph closes no cage receives no label.
\item A device-resident batch of \(F=\min(F_{\mathrm{req}},\lfloor A/W\rfloor)\) frames on one GPU stream (Algorithm~\ref{alg:resident}).
\item Ring-size cage signatures with guest occupancy, ion-shell classes and topology keys on the same ring network (Sections~\ref{sec:signatures} and \ref{sec:topokeys}).
\item A C++20 library with a command line, a Python package, a Lua library, a PLUMED action and a LAMMPS compute, in place of the single YAML--Lua executable (Section~\ref{sec:code}).
\item A Lean~4 development of 55 theorems about the abstract definitions (list lengths, ball agreement, the completion closure, batch arithmetic). This supplement states 34 of them as Theorems~S1--S34, and \texttt{lake build} checks them on Lean 4.15.0 with no other library.
      It does not model the C++ code (Section~\ref{sec:lean}).
\end{enumerate}

The other codes overlap with parts of that list.
OVITO \citep{stukowski2010visualization} ships CHILL+ \citep{nguyen2015identification} and polyhedral template matching \citep{larsen2016robust} as modifiers.
freud and pyscal evaluate the Steinhardt and averaged Steinhardt parameters \citep{ramasubramani2020freud,menon2019pyscal}.
RINGS counts the rings of a network under the King \citep{king1967ring}, Guttman \citep{guttman1990ring}, primitive and strong \citep{goetzke1991rings} definitions \citep{leroux2010rings}.
GRADE determines the clathrate hydrate cages of a trajectory from its water rings \citep{mahmoudinobar2019grade}.
GenIce2 builds ice and clathrate lattices, and its \texttt{analice2} front end reports the rings and polyhedral cages of a given configuration \citep{matsumoto2018genice,matsumoto2024genicecore}.
RINGS, GRADE and GenIce2 compute ring and cage topology; none of these codes returns d-SEAMS HC or DDC membership per molecule, an incremental ring update or a seeded assignment.
We time every computation they share with d-SEAMS 2.0 on the same frames in Section~\mref{sec:comparison}.

\subsection{Frame loops of other codes}
\label{sec:s-relwork}
Ring and cage codes compute the rings of each frame from that frame alone.
RINGS \citep{leroux2010rings} and ISAACS \citep{leroux2010isaacs} clear the ring table and rebuild the neighbour table at every step of a trajectory.
Their successor atomes \citep{leroux2025atomes} does the same and distributes the steps over OpenMP threads.
ChemNetworks \citep{ozkanlar2014chemnetworks} processes one snapshot per run, and its LAMMPS fix passes it one configuration at a time.
cycless, which replaced CountRings \citep{matsumoto2007topological,matsumoto2026cycless}, takes one graph per call, and RingsStatisticsMatter.jl \citep{morrow2023ringsstatisticsmatter} takes one structure.
fastatomstruct \citep{holle2025fastatomstruct,holle2026fastatomstruct}, Trajan \citep{maksimov2026trajan} and RingDetect-HPC \citep{ringdetecthpc2026} loop over the frames of a trajectory and search each frame on a new graph.
None of these papers, manuals or sources describes reusing the rings of an earlier frame.
GenIce2 \citep{matsumoto2018genice,matsumoto2024genicecore}, which uses cycless, and GRADE \citep{mahmoudinobar2019grade} find the rings and cages of each frame from that frame.
OVITO \citep{stukowski2010visualization} and freud \citep{ramasubramani2020freud} evaluate CHILL+ and the Steinhardt parameters frame by frame.
Between two saved frames of a crystal or of a growing nucleus most bonds persist, and the rings away from the changed bonds are the rings of the previous frame.
d-SEAMS 1 \citep{goswami2020dseams} also rebuilt every frame.
Its ring search grew every closed walk shorter than the largest ring size and kept the walks that pass Franzblau's shortest-path test \citep{franzblau1993computation}.
Its HC search tested every pair of six-rings, so its cost grew with the square of the ring count.
Table~\ref{tab:features} lists the computations built into each code.

\begin{table}[htbp]
\caption{\label{tab:features}Computations built into each code.
A dash means the code does not provide the computation.
The versions are those of Table~\mref{tab:compare}.}
\centering
\footnotesize
\setlength{\tabcolsep}{2.6pt}
\begin{tabular}{L{2.2cm}L{1.75cm}L{2.45cm}L{1.3cm}L{1.25cm}L{1.4cm}L{1.9cm}}
 & d-SEAMS 1 & d-SEAMS 2.0 & OVITO & freud & GenIce2 & RINGS\\
\hline
primitive rings & walks, then filter & shortest-path pairs & -- & -- & yes & yes\\
ring update between frames & rebuild & exact, changed-atom set & -- & -- & -- & --\\
HC and DDC membership per molecule & pair scan & inverted index & -- & -- & -- & --\\
cage update between frames & rebuild & exact, ring adjacency & -- & -- & -- & --\\
seeded four-nearest assignment & -- & yes & -- & -- & -- & --\\
prisms and confined-ice criteria & yes & yes & -- & -- & -- & --\\
polyhedral cages by ring signature & -- & yes, with guests & -- & -- & yes & --\\
CHILL+ & yes & yes & yes & -- & -- & --\\
averaged Steinhardt \(\bar q_6\) & yes & yes & -- & yes & -- & --\\
template matching (PTM) & -- & -- & yes & -- & -- & --\\
common-neighbour analysis, identify diamond & -- & -- & yes & -- & -- & --\\
GPU batch & -- & yes & -- & -- & -- & --\\
interfaces & Lua, YAML & CLI, Python, R, Julia, Lua, PLUMED, LAMMPS & GUI, Python & Python & CLI, Python & CLI (Fortran, MPI)\\
\end{tabular}
\end{table}

\section{Neighbour graphs}
\label{sec:s-graphs}
A hexagonal cage identifies ice Ih and a double-diamond cage identifies ice Ic, as in d-SEAMS \citep{goswami2020dseams}.
CHILL and CHILL+ remain the four-neighbour staggered and eclipsed labels of \citet{moore2010chill} and \citet{nguyen2015identification}. \citet{jacobson2010amorphous} count clathrate nuclei as \(5^{12}\) and \(5^{12}6^{2}\) polyhedra, which are not those ice cages.
We restate here the published tests on a neighbour graph \(G\), the input of the searches in Section~\mref{sec:algorithms}.
Voronoi-weighted \(q_l\), SOAP \citep{bartok2013representing,caro2019optimizing} and IRA templates \citep{gunde2021ira} are available in the library and are not the ice score.


A frame is a finite set of oxygen positions \(V\) in a periodic cell together with a neighbour relation \(E\subset V\times V\).
We write \(G=(V,E)\) for the undirected graph and \(N(v)\) for the neighbours of \(v\).
For tetrahedral water the published CHILL and CHILL+ tests take \(|N(v)|=4\) when those four neighbours exist \citep{moore2010chill,nguyen2015identification}.
We use two constructions of \(E\) in this paper.

A cutoff graph bonds \(u\) to \(v\) when the minimum-image separation satisfies \(\|x_u-x_v\|\le r_c\).
For water, \(r_c=3.5\,\mathring{\mathrm{A}}\) is the clustering cutoff of \citet{nguyen2015identification}.
CHILL itself takes the four closest neighbours \citep{moore2010chill}.
On an ordered ice lattice the first shell lies inside \(r_c\) and the second shell lies outside it, so the cutoff graph coincides with the four nearest neighbours of each molecule.
On a thermally disordered or confined frame that identity fails.
Some molecules carry fewer than four neighbours inside \(r_c\), and some pairs inside \(r_c\) are not first-shell neighbours of each other.

A \(k\)-nearest graph, with \(k=4\), nominates for each \(v\) the \(k\) molecules of smallest minimum-image distance.
An exact tie at the \(k\)-th distance is resolved by the traversal order of the neighbour backend; ideal ice lattices have none, since their fifth neighbour lies beyond the first shell.
Two undirected graphs are then available.
The mutual graph \(G_{\mathrm{mut}}\) retains an edge when both endpoints nominate each other.
The union graph \(G_{\cup}\) retains an edge when either nominates the other.
On an undistorted crystal the first shell is mutual, so \(G_{\mathrm{mut}}=G_{\cup}\) coincides with the cutoff graph.
On a disordered liquid the union graph admits one-sided edges that the mutual graph refuses.
In Section~\mref{sec:algorithms} we use that pair as the two thresholds of a hysteresis construction.

Write \(B_1(v)=\{v\}\cup N(v)\) for the closed star of \(v\) and \(B_2(v)\) for the two-hop ball.
A score that is a function of \((G[B_1(v)],\{x_u:u\in B_1(v)\})\) alone is a star function; Eq.~\ref{eq:ql} at \(N_b=4\) is one.
The scores of the next subsection are not all star functions, and Section~\ref{sec:cagestar} separates the two cases.

\section{Bond-orientational order and CHILL+}
\label{sec:s-boo}
\subsection{Bond-orientational order}
\label{sec:orgca7518c}

The Steinhardt bond-order parameters \citep{steinhardt1983bond} characterize the star of molecule \(i\) through the spherical harmonics \(Y_{lm}\) of the unit vectors \(\hat r_{ij}\) toward its neighbours.
The complex vector
\begin{equation}
q_{lm}(i)
  = \frac{1}{N_b(i)}\sum_{j=1}^{N_b(i)} Y_{lm}(\hat r_{ij}),
\qquad
N_b(i)=4,
\label{eq:qlm}
\end{equation}
is the mean of those harmonics over the four nearest neighbours \citep{moore2010chill,nguyen2015identification}.
The rotationally invariant
\begin{equation}
q_l(i)
  = \sqrt{\frac{4\pi}{2l+1}\sum_{m=-l}^{l}\bigl|q_{lm}(i)\bigr|^2}
\label{eq:ql}
\end{equation}
is the length of that vector.
The neighbour-averaged \(\bar q_l\) of \citet{lechner2008accurate} replaces \(q_{lm}(i)\) by the mean of \(q_{lm}\) over \(B_1(i)\), so it reads a two-hop ball as averaged spherical harmonics, not as a ring complex. d-SEAMS reports both \(q_l\) and \(\bar q_l\).
The CHILL and CHILL+ algorithms use \(l=3\).

Averages of \(q_{lm}\) over every neighbour inside a distance cutoff diverge from the published definition, which specifies exactly four neighbours.
On ordered ice the two averages coincide.
On the disordered example trajectory they do not (Section~\ref{sec:verification}).
\subsection{CHILL and CHILL+}
\label{sec:org1f78dce}

CHILL \citep{moore2010chill} and CHILL+ \citep{nguyen2015identification} reduce the star of \(i\) to four scalars, one per neighbour.
The correlation
\begin{equation}
c_{ij}
  = \frac{\sum_{m=-l}^{l} q_{lm}(i)\,q_{lm}^*(j)}
         {\sqrt{\sum_{m}\lvert q_{lm}(i)\rvert^2}\,
          \sqrt{\sum_{m}\lvert q_{lm}(j)\rvert^2}}
\label{eq:cij}
\end{equation}
is the cosine of the angle between the two \(q_l\) vectors.

In ice Ic every water--water bond is staggered; in ice Ih the bond along the \(c\) direction is eclipsed and the other three are staggered \citep{moore2010chill}.

The algorithms encode those geometries as windows on \(c_{ij}\).
A bond is staggered when
\begin{equation}
c_{ij}\le -0.8
\label{eq:staggered}
\end{equation}
in both CHILL and CHILL+.
A bond is eclipsed when
\begin{equation}
-0.20\le c_{ij}\le -0.05
\label{eq:eclipsed-chill}
\end{equation}
in CHILL, and when
\begin{equation}
-0.35\le c_{ij}\le 0.25
\label{eq:eclipsed-plus}
\end{equation}
in CHILL+.
Every other value is out of range.
Write \(S(i)\) and \(E(i)\) for the number of staggered and eclipsed bonds among the four neighbours of \(i\).
The published tables then assign
\begin{align}
\mathrm{Ic}(i)
  &\iff S(i)=4,\; E(i)=0,
\label{eq:chill-ic}\\
\mathrm{Ih}(i)
  &\iff S(i)=3,\; E(i)=1,
\label{eq:chill-ih}\\
\mathrm{clathrate}(i)
  &\iff E(i)=4,\; S(i)=0,
\label{eq:chill-clath}\\
\mathrm{interClathrate}(i)
  &\iff E(i)=3.
\label{eq:chill-interclath}
\end{align}
The last two rows exist only in CHILL+. Interfacial ice is the remaining four-coordinated case that still has a staggered neighbour. CHILL treats \(i\) as interfacial when \(S(i)=2\) and some \(j\in N(i)\) has \(S(j)>2\), or when \(S(i)=3\), \(E(i)=0\) and some \(j\in N(i)\) has \(S(j)=2\) \citep{moore2010chill}. CHILL+ keeps the first rule and relaxes the second.
In the relaxed rule the neighbour of a \(3+0\) molecule needs \(S(j)>1\) \citep{nguyen2015identification}.
Every other molecule is liquid, and so is every molecule with fewer than four listed neighbours.

The four neighbours of Eqs.~\ref{eq:qlm} and \ref{eq:cij} are the four nearest entries of a neighbour list, and we use two such lists.
CHILL+\((r_c)\) takes them from the cutoff graph at \(r_c=3.5\,\mathring{\mathrm{A}}\).
It is the variant we report in Tables~\ref{tab:sota}, \ref{tab:atlas}, \ref{tab:ions} and \ref{tab:icescores} and in the label walks of Section~\ref{sec:applications}.
It is also the variant of \texttt{seams chill-plus}, of the Python and Lua \texttt{chill\_plus} calls, and of the PLUMED action in the seeding and brine campaigns.
CHILL+(4NN) takes them from the mutual four-nearest graph \(G_{\mathrm{mut}}\), in which a molecule can have fewer than four neighbours and is then liquid.
It is therefore a mutual-graph variant and not the four-neighbour list of \citet{nguyen2015identification}.
The \texttt{seams chill-plus -{}-graph knn} option, the LAMMPS compute and the device kernel of Section~\ref{sec:resident} use it.

The Nature study of the crystallization rate \citep{moore2011structural} uses CHILL.
The algorithm itself is defined in the nanopore paper.
In \citet{li2011homogeneous} ice is identified with a distance-shell \(q_6\) (mean neighbour correlation at \(l=6\), threshold \(0.5\)), and the core is later split with \(q_3\).
That identifier is not CHILL+ and is not the ice score of this engine.

\section{Notes on the ring and cage definitions}
\label{sec:s-defs}
d-SEAMS identifies hexagonal cages and double-diamond cages from the six-rings of \(G\) \citep{goswami2020dseams}. \citet{jacobson2010amorphous} count clathrate nuclei by \(5^{12}\) and \(5^{12}6^{2}\) cages, not those ice units.
Franzblau's SP ring sits at the centre of his \(m\)-irreducibility hierarchy; Yuan's primitive ring is the same object.
Strong rings are a stricter subclass \citep{yuan2002rings}.
King's shortest-path analysis is a different object.
It neither contains nor is contained in the primitive set \citep{franzblau1993computation,yuan2002rings}.
Ice Ih and ice Ic are built from six-rings.
In the benchmarks of Section~\mref{sec:benchmarks} we take \(D=6\) as the largest ring size sought.
The accuracy sweeps and the cage benchmark take \(D=7\) and keep the six-rings, which are the same set at both values because \(\lfloor D/2\rfloor=3\) for both.
Eq.~\mref{eq:franzblau} is the test.
The candidates of \citet{yuan2002rings} come from prime mid-nodes.
Every primitive ring has that form, and not every such candidate is primitive.
Candidates therefore arise only from shortest paths, and each ring is enumerated once from its lowest-indexed member.
In \citet{goswami2020dseams} we enumerated closed walks and applied Eq.~\mref{eq:franzblau} afterwards.
The method of \citet{yuan2002rings} still tests primitivity after the mid-node construction and does not stop at the criterion.
The choice of \(p\), \(q\) and the matching covers every starting atom and direction of both rings, so the pair test does not depend on how a ring is stored.
The cage is the two basals together with those prismatic rings, twelve vertices in all.
With \(j=2\) these are the alternating-intersection conditions of d-SEAMS \citep{goswami2020dseams}.
The choice of \(j\) and of each \(R'_k\) makes the test independent of the starting atom and direction of \(R\) and of the order of the ring list.
The ice score of this engine is cage membership.
Only a ring that is not HC-affiliated is tested as an equator, as in the baseline commit of Section~\mref{sec:benchmarks}; the 2020 release ran the DDC search first and tested every ring.
A peripheral can itself be HC-affiliated, and a molecule that belongs to both cage types (a stacking fault) carries both labels.
Eq.~\mref{eq:ice} is invariant under any motion that preserves \(G\), and it changes when a bond is made or broken even if the coordinates move only slightly.
The engine evaluates \(\alpha_{\mathrm{HC}}\) for every ring and then \(\alpha_{\mathrm{DDC}}\), from the ring list and the neighbour list alone.
The d-SEAMS 1 search assembled cage objects by claiming rings as it accepted cages, in the order it visited them.
The flags of Eqs.~\mref{eq:ddc-aff} and \mref{eq:ice} test every ring without claiming, so that an incremental update equals a full evaluation.
Where cages overlap they can flag a ring that the claiming search leaves to an earlier cage.
We check that on mW cubic and on an isolated hexagonal cage the two coincide ring for ring (Section~\mref{sec:benchmarks}).

\section{Seeded hysteresis and ring completion}
\label{sec:s-seeded}
\label{sec:org393f5ae}

At a Gaussian displacement of \(0.3\,\mathring{\mathrm{A}}\) the cutoff graph loses edges as first-shell neighbours cross \(3.5\,\mathring{\mathrm{A}}\), and Eq.~\mref{eq:ice} on that graph degrades (Table~\ref{tab:sota}).
The four nearest neighbours of each molecule keep their identity.
We therefore evaluate the same cage criteria on the mutual and union \(k\)-nearest graphs of Section~\mref{sec:definitions}, with \(k=4\).

Whenever the \(k\)-th neighbour of every molecule lies inside the cutoff and the \((k+1)\)-th lies outside it, the cutoff graph and both \(k\)-nearest graphs coincide edge for edge.
On a crystal the first shell is mutual, so \(G_{\mathrm{mut}}=G_{\cup}\) on an undistorted lattice.
A union graph admits one-sided edges, which close rings that neither partner's own first shell supports.
A mutual graph removes those edges, and on the liquid of Table~\ref{tab:sota} it labels \(0.023\) of the molecules.

The residual loss on a jittered crystal is a boundary effect.
A mutual bond drops when one partner's nomination fails, and a single missing edge dissolves every cage through it.
The union graph keeps those edges, together with the one-sided edges of a disordered region.
When we replace \(G_{\cup}\) by the mutual five-nearest graph, the seeded values at \(0.4\,\mathring{\mathrm{A}}\) drop to \(0.574\) and \(0.505\), because the extra edges are chords that break the primitivity of the six-rings (Eq.~\mref{eq:franzblau}).

The remaining construction is a hysteresis, the double-threshold logic of edge tracking applied to bonds.
We treat each cage label \(X\in\{\mathrm{HC},\mathrm{DDC}\}\) separately.
Write \(\alpha^X_{\mathrm{mut}}(v)\) and \(\alpha^X_{\cup}(v)\) for the label \(X\) of Eq.~\mref{eq:ice} on \(G_{\mathrm{mut}}\) and on \(G_{\cup}\), and \(A^X=\{v:\alpha^X_{\mathrm{mut}}(v)\lor\alpha^X_{\cup}(v)\}\).
Affiliation on \(G_{\mathrm{mut}}\) supplies seeds.
A molecule of \(A^X\) is accepted only when its component in \(G_{\cup}[A^X]\) contains a seed,
\begin{equation}
\mathrm{seeded}^X(v)
  \iff
  \exists\, u\in C^X(v)
  \text{ with }
  \alpha^X_{\mathrm{mut}}(u),
\label{eq:seeded}
\end{equation}
where \(C^X(v)\) is the connected component of \(G_{\cup}[A^X]\) that contains \(v\in A^X\).
A molecule with \(\alpha^X_{\mathrm{mut}}\) and not \(\alpha^X_{\cup}\) is its own seed and is accepted; an edge of \(G_{\cup}\) absent from \(G_{\mathrm{mut}}\) can make a mutual-graph six-ring non-primitive, so \(\alpha^X_{\mathrm{mut}}\) does not imply \(\alpha^X_{\cup}\).
If the mutual pass affiliates no molecule, nothing is accepted, whatever rings \(G_{\cup}\) builds (Theorem~S17).
On the liquid of Table~\ref{tab:sota} the mutual pass affiliates \(0.023\) of the molecules, and those components are what the seeded rule labels.
One linked-cell walk writes both graphs.

The hysteresis leaves unlabelled a molecule whose own cage lost one edge while the cages around it did not.
We count those molecules with \texttt{scripts/seeded\_residual.py} on seeds 0 to 2.
At \(0.3\,\mathring{\mathrm{A}}\) the hysteresis misses 0 to 2 of the 512 Ic sites and 0 to 5 of the 360 Ih sites in the five realizations of Table~\ref{tab:sota}.
A six-ring of \(G_{\cup}\) with every vertex labelled but one is a cage ring with a vacancy in the label, and its last vertex takes the label,
\begin{equation}
\mathrm{completed}(v)
  \iff
  \mathrm{seeded}(v)
  \ \lor\ 
  \exists\, R\ni v:\ \forall u\in R\setminus\{v\}:\ \mathrm{completed}(u),
\label{eq:completion}
\end{equation}
taken as the least fixed point over the six-rings \(R\) of \(G_{\cup}\), with \(\mathrm{seeded}=\mathrm{seeded}^X\), separately for each label \(X\).
The rule fills a ring only when five of its six vertices already carry the label, so a ring at the edge of a labelled region with one or two labelled vertices is not filled.
With no seeded molecule no ring has five labelled vertices, and nothing is completed (Theorem~S21).
At \(0.4\,\mathring{\mathrm{A}}\), where the seeded pass alone keeps \(0.698\) and \(0.589\) of the two lattices, the completed rule keeps \(0.942\) and \(0.728\) (Table~\ref{tab:sota}).
The implementation is \texttt{ringAdjacentCompletion}, a stack of rings that starts with every six-ring of \(G_{\cup}\) and pushes again only the rings through a newly labelled vertex.
We prove in Proposition~P5 of the supplement that it returns the least fixed point for every pop order.
We state in Lean, for the inductive definition and not for the C++ loop, that a completed molecule is seeded or sits on a ring whose other molecules are completed (S19).
Lean also states that every set closed under the rule contains the completion (S20) and that an empty seed completes nothing (S21).
On a five-molecule instance a rule that labels every vertex of a ring sharing an edge with a fully labelled ring labels a molecule the completion does not (S22).

\begin{algorithm}[htbp]
\caption{Seeded hysteresis with ring completion}
\label{alg:seeded}
\begin{algorithmic}[1]
\Require positions \(X\); cell \(H\); \(k=4\); one cage label \(X\) (HC or DDC)
\Ensure the completed labels of Eq.~\ref{eq:completion} (Proposition~P5)
\State \(N_k(v)\gets\) the \(k\) nearest neighbours of each \(v\) \Comment{linked-cell walk on \(H\)}
\State \(E_{\mathrm{mut}}\gets\{uv:u\in N_k(v)\wedge v\in N_k(u)\}\); \(E_{\cup}\gets\{uv:u\in N_k(v)\vee v\in N_k(u)\}\)
\State \(\alpha^X_{\mathrm{mut}},\alpha^X_{\cup}\gets\) Eq.~\mref{eq:ice} on \(G_{\mathrm{mut}}\) and on \(G_{\cup}\) \Comment{Algorithms~\mref{alg:rings}, \mref{alg:cages}}
\State \(A\gets\) molecules of \(A^X\) reached from a seed (\(\alpha^X_{\mathrm{mut}}\)) by a sweep of \(G_{\cup}[A^X]\) \Comment{Eq.~\ref{eq:seeded}}
\State \(F\gets\) every six-ring of \(G_{\cup}\) \Comment{stack}
\While{\(F\ne\emptyset\)}
  \State pop \(R\) from \(F\)
  \If{\(R\setminus A=\{u\}\)}
    \State \(A\gets A\cup\{u\}\); push every six-ring through \(u\) not on \(F\)
  \EndIf
\EndWhile
\State \Return \(A\)
\end{algorithmic}
\end{algorithm}

The \(k\)-nearest walk examines \(m\) candidate points at \(O(\log k)\) each, \(O(m\log k)\), and both graphs have \(O(nk)\) edges.
The walk stores the binned coordinates, \(O(n+C)\) for \(C\) cells, and the \(nk\) nominations; each worker adds a heap of \(k\) entries and at most a copy of the molecules in the cells it scans.
Each source examines at most the \(n-1\) other points, so \(m=O(n^2)\) in the worst case.
When every molecule has a bounded number of molecules within the cell shells that reach its \(k\)-th neighbour, \(m=O(n)\).
That holds at the bounded density of a liquid or a solid.
The two ice scores are Algorithms~\mref{alg:rings} and \mref{alg:cages} on graphs of degree at most \(k\) (mutual) and \(k+\iota\) (union), where \(\iota\) is the largest number of molecules that nominate one molecule.
The seed sweep is \(O(n+|E_{\cup}|)\).
The stack starts with \(R_6\) rings, and each of at most \(n\) labelling events pushes at most \(d_I\) rings of six vertices, so the completion is \(O(R_6+n\,d_I)\).
A ring with one unlabelled vertex is on the stack after its last labelled vertex was labelled.
The loop therefore stops at the least set that contains the seeds and is closed under the all-but-one rule, whatever the pop order (Proposition~P5).
With no mutual seed \(A\) is empty and stays empty.

\section{Cage signatures, occupancy and ion shells}
\label{sec:signatures}

A signature is a ring-size census \(\{s:n_s\}\).
The enumerator grows a connected set of primitive rings on the ring-adjacency graph, in which two rings are adjacent when they share an edge.
It stays inside that census and keeps a set whose edges form a closed polyhedron.
In such a set every edge of every face belongs to exactly two faces of the set,
\begin{equation}
\mathrm{closed}(F)
  \iff
  F\neq\emptyset
  \ \land\
  \forall e\in E(F):\ \bigl|\{f\in F:e\subset f\}\bigr|=2.
\label{eq:closed-poly}
\end{equation}
Named \texttt{hc} and \texttt{ddc} call the TUM finders so the vertex sets match Eq.~\mref{eq:ice}.
Distinct cages are distinct sorted vertex sets.

A guest (methane, THF, an ion) is not a vertex of the water graph.
We define occupancy as two predicates on the same cages and the same guests.
Write \(c_F\) for the periodic centroid of the vertices of \(F\) and \(d(g,F)\) for the minimum-image distance between guest \(g\) and \(c_F\).
The radius assignment is
\begin{equation}
\mathrm{radius}(g)
  =
  \arg\min_{F:\,d(g,F)\le r} d(g,F),
\label{eq:occ-radius}
\end{equation}
with \(r=4\,\mathring{\mathrm{A}}\), half a \(5^{12}\), by default in \texttt{seams cages -{}-guest-radius} and \texttt{guest\_occupancy}; a guest with no centre inside \(r\) is free.
The inside assignment unwraps vertices and the guest about \(c_F\), fan-triangulates each face from the face centroid, and takes
\begin{equation}
\begin{aligned}
\mathrm{inside}(g,F)
  \iff{}& g\text{ lies on a triangle of }F\\
  &\lor\ \#\{T\subset F: \text{a ray from }g\text{ meets }T\}
  \text{ is odd}.
\end{aligned}
\label{eq:occ-inside}
\end{equation}
A hit on a shared edge retries a perturbed direction.
A guest inside two cages goes to the nearer centroid.
\texttt{seams cages -{}-guest-types} applies the centroid predicate, and adding \texttt{-{}-inside} applies the ray-parity predicate; \(r\) is not part of Eq.~\ref{eq:occ-inside}.

An ion is classed by its shell, the water molecules within a fixed cutoff of it.
Every shell molecule is labelled ice, none is, or the shell is mixed (ice, liquid, front).
A water-only primitive ring whose vertices all lie in the shell is capped, and the ion then sits above a complete water ring.
We add to the water graph a virtual ion bonded to every shell water.
A primitive ring of that graph that contains the ion is broken, because the water-only network lost that ring.
The Lean statements for this subsection concern the predicates as definitions, not the geometry.
An empty face list is not closed, and a radius assignment lies within \(r\).
The parity predicate holds for a point on a face or an odd crossing count and fails for an even count off the faces.
An empty ring is not capped.

\section{Topology keys}
\label{sec:topokeys}
The ring and cage graphs carry more than the two ice I labels.
A local key names the isomorphism class of the rooted bonded graph within \(h\) bonds of a molecule.
Two molecules share a key exactly when their neighbourhoods are the same graph with the same centre, whatever the atom indices.
With nauty \citep{mckay2014nauty} linked the key is the canonical certificate of the rooted graph.
The centre sits in a colour cell of its own, and each species has one cell when species are given.
Without nauty the key is a Weisfeiler--Leman colour refinement hash seeded by degree, species and the root flag.
A frame key hashes the sorted multiset of per-molecule refinement hashes together with the primitive ring census.
Two configurations share a frame key when every molecule has a counterpart with the same environment, and a single broken bond changes it.
The k-ART method uses that classification to index its event catalogue \citep{trochet2015kart}.
We compute the same classification on the graphs the engine already builds, and a configuration database can recognise a state up to relabelling.
Every molecule of a periodic diamond lattice carries the same key, which the Catch2 suite checks.
The bonded graph of the stored mW cubic frames is that lattice, so their \(4096\) molecules form one class.
A zincblende colouring of the diamond lattice splits it into two equal classes and survives a joint permutation of molecules and species.
A key library collects the keys of reference structures under labels and names any molecule whose key it contains.
Two hops close no ring in a four-connected network, so every tetrahedral ice shares one two-hop key.
A library therefore needs three hops before ice polymorphs separate by their ring structure.
We build that library over the GenIce lattices of Table~\ref{tab:atlas} in the reproducibility package.
The keys reach the command line as \texttt{seams fingerprint}, with \texttt{-{}-{}emit-library} and \texttt{-{}-{}library}, and the Lua and Python front ends as \texttt{topologyFingerprint}, \texttt{topology\_library} and \texttt{classify\_topology}.
In Table~\ref{tab:library} we read the three-hop library of twenty-four GenIce lattices back against the jittered ice I lattices of Table~\ref{tab:sota} and against the lattices themselves at rest and displaced.
At rest every network is named.
The only keys shared between references are the \(5^{12}\) cage environment common to the sI, sII, sT and sIV frameworks and the environment ice III and IX share with the hypothetical ice A. Under displacement the exact match is only as stable as the graph.
At \(0.16\,\mathring{\mathrm{A}}\) ice Ih, Ic, XI, 0 and XVII and the FAU, RHO and SOD networks keep all their keys, and sI and sT keep their fractions at rest.
At that displacement iceB keeps \(0.820\), and iceA keeps \(0.642\) under the key it shares with ice III and IX.
The remaining dense polymorphs, whose mutual four-nearest graph is not the bonded graph, keep at most \(0.023\) of their own labels.
At \(0.31\) no network keeps more than \(0.194\).
The seeded cage assignment remains the ice I classifier under displacement.
The key library is the exact match, and it is the one that names the other polymorphs.

A second library at two hops sits behind the three-hop one.
A molecule the deep library does not name is looked up in the shallow one, and the match records the depth that named it.
The fallback changes nothing where the exact match already names the molecule.
Among the open networks at rest and at \(0.16\,\mathring{\mathrm{A}}\) it acts only on iceB at \(0.16\,\mathring{\mathrm{A}}\).
There it names the \(69\) of \(384\) molecules that three hops leave unnamed.
Where the exact match fails it names the residue, and the depth says how far to trust the name.
On the GenIce Ih lattice displaced by \(0.31\,\mathring{\mathrm{A}}\) three hops name \(0.09\) of the molecules with their own label.
The fallback adds \(0.41\) under the label the tetrahedral ices share and \(0.35\) under another polymorph's label, and it leaves \(0.16\) unnamed.
On the dense polymorphs at \(0.16\,\mathring{\mathrm{A}}\) the mutual four-nearest graph is not the bonded graph.
There the fallback names \(0.79\) of ice IV and \(0.55\) of ice XII with other polymorphs' labels and none with their own.
Ice VII and VIII recover \(0.64\) and \(0.48\) of their own labels at two hops.
A two-hop name is a family-level name, correct as far as the graph is, and the depth field is what separates it from the exact three-hop match.

\begin{table}[htbp]
\caption{\label{tab:library}Fraction of molecules named with their own label by the three-hop key library on the mutual four-nearest graph.
Left: the jittered ice I lattices of Table~\ref{tab:sota}, means over three realizations.
Right: GenIce lattices at rest and at \(\sigma = 0.16\,\mathring{\mathrm{A}}\); structures that share one oxygen network carry one label (Ih/XI, III/IX, XII/XIV, VII/VIII, XVI/sII).
A key shared by several references names all of them, so sI, sII, sT and sIV at rest are named as a group where their \(5^{12}\) environments coincide.}
\centering
\scriptsize
\resizebox{\textwidth}{!}{%
\begin{tabular}{lcccc|lcc}
lattice & \(0.0\) & \(0.1\) & \(0.2\) & \(0.3\) & structure & rest & \(0.16\)\\
\hline
Ic & 1.000 & 1.000 & 0.926 & 0.174 & Ih, Ic, XI, 0 & 1.000 & 1.000\\
Ih & 1.000 & 1.000 & 0.976 & 0.182 & XVII, FAU, RHO, SOD & 1.000 & 1.000\\
 &  &  &  &  & iceB & 1.000 & 0.820\\
 &  &  &  &  & IV, VI, XII & 1.000 & 0.000\\
 &  &  &  &  & V, VII, VIII, XIV & 1.000 & 0.004--0.023\\
 &  &  &  &  & III, IX (with iceA) & shared & 0.000\\
 &  &  &  &  & sI & 0.478 & 0.478\\
 &  &  &  &  & sT & 0.287 & 0.287\\
 &  &  &  &  & sII, sIV & shared & shared\\
\end{tabular}
}
\end{table}

\section{Numerical techniques}
\label{sec:org387ec10}

Cutoff neighbour lists are the periodic pair search of linkcell \citep{goswami2026linkcell}; vesin \citep{vesin} takes over when linkcell is not built or fails, and a brute-force loop after that.
The engine commits of the timing records predate the linkcell pair search and took their cutoff pairs from vesin.
The four nearest neighbours of each molecule are not that pair list.
They come from a linked-cell \(k\)-nearest search \citep{goswami2026linkcell}.
It bins molecules in fractional coordinates, keeps a \(k\)-heap per source, and grows Chebyshev shells until the \(k\)-th neighbour cannot sit outside the visited cube.
The search does not take a cutoff; a cell edge is only the bin width.
The LAMMPS dump box (bound spans plus tilt \(xy,xz,yz\)) is recovered as the three lattice vectors \(H\) of the minimum-image convention.
The independent wrap on each box edge is not the sheared minimum image; the linkcell pair search, vesin and the linked-cell walk take that \(H\).
Without the wrap, an in-tree linked-cell walk takes its place on orthorhombic boxes and an all-pairs \(k\)-heap on tilted ones.
The same walk is the first kernel of the device-resident batch below.
The brute-force cutoff loop and the skin list of the trajectory update evaluate minimum-image squared distances in batches when the box is orthorhombic.
For a displacement \((\Delta x,\Delta y,\Delta z)\) and box \((L_x,L_y,L_z)\),
\begin{equation}
\tilde{\Delta x}
  = \lvert\Delta x\rvert
    - L_x\,\mathrm{round}(\lvert\Delta x\rvert/L_x),
\qquad
r^2
  = \tilde{\Delta x}^2+\tilde{\Delta y}^2+\tilde{\Delta z}^2.
\label{eq:simd_dist}
\end{equation}
A tilted dump does not take that wrap, and its pair distances go through \(H\).
Unlike-type neighbour lists and the partial \(g_{IJ}(r)\) of \texttt{seams rdf} use the same \(H\).
The \(q_{lm}\) of Eq.~\ref{eq:qlm} are closed forms at \(l=3\) and \(l=6\).
They share powers of \(\sin\theta\) and \(\cos\theta\) across orders and generate \(e^{im\phi}\) by recurrence.
The Condon--Shortley relation \(Y_{lm}=(-1)^m Y_{l,-m}^{*}\) is then exact rather than rounded from two independently evaluated exponentials.

\section{Device-resident ice-score batch}
\label{sec:resident}

The searches above classify one frame on the host.
The cage assignment on the mutual four-nearest graph is also a device-resident batch of frames when that working set fits on the GPU.
A runtime probe reports free device memory \(M\), and the budget is \(A=\lfloor sM\rfloor\) with \(s=0.8\).
We write \(W>0\) for the working set of one frame, with each array rounded up to 256 bytes.
The working set holds the coordinates, the neighbour columns, the mutual-graph CSR, the six-ring CSR, the basal-pair buffer, affiliation flags and the per-atom HC/DDC labels.
For a request of \(F_{\mathrm{req}}\) frames the batch is
\begin{equation}
F=\min\bigl(F_{\mathrm{req}},\lfloor A/W\rfloor\bigr).
\label{eq:resident}
\end{equation}
Then \(FW\le A\), and when \(F<F_{\mathrm{req}}\) one more frame of size \(W\) would exceed \(A\) (Theorems~S30 and S31).
That \(F\) is the offload unit.
The suite compares a specialized CHILL+(4NN) kernel, on the device when one is present, with the host labels on the 250 oxygens of the example trajectory.

Coordinates for the \(F\) frames upload once.
Every subsequent kernel is enqueued on one stream \(Q\) (Figure~\ref{fig:gpuresident}).
The kernels fold and bin the coordinates, run an exclusive scan of the cell counts and scatter into cell-major order.
They then run the Chebyshev-shell \(k\)-nearest walk of \citet{goswami2026linkcell} and build the mutual four-nearest graph and the primitive hexagons.
The last kernels build the inverted index of Eq.~\mref{eq:index} and the basal pairs, then evaluate HC and DDC and write the per-atom labels.
The host waits on \(Q\) once and downloads only those labels.
The kernels are compiled at run time by NVRTC.

In seams-core \coretag{} the \(k\)-nearest walk is one launch over all \(F\) frames.
Per-frame orthorhombic lengths are passed to that launch, so an NPT dump does not share a single cell.
Occupant stride \(\tau\) and thread-block size \(B\) are chosen at launch to maximize the number of sources per block \(B/\tau\) under the compiled kernel's thread limit and a shared-memory cap.
The number of thread blocks is
\begin{equation}
\bigl\lceil nF\tau/B\bigr\rceil,
\label{eq:blocks}
\end{equation}
and a grid with fewer blocks than streaming multiprocessors (SMs) leaves some of them without work.
The A100 has 108 SMs \citep{nvidia2020a100}, and every later kernel launches blocks of 128 threads.
The per-atom kernels (mutual graph, hexagon enumeration and the sort of each atom's ring row) take one thread per atom, a grid of \(\lceil nF/128\rceil\) blocks.
That grid is 32 blocks at \(n=4096\) and \(F=1\), below the SM count, and 352 for the stored eleven-frame batch (Figure~\ref{fig:gpuoccupancy}).
The ring kernels (index inversion, basal emission, DDC and the atom labels) take one thread per ring slot, with \(16\) slots per atom by default.
The HC kernel takes one thread per basal-pair slot, eight per ring slot.
Those grids are therefore \(16\) and \(128\) times larger and exceed the SM count already at \(F=1\).
A source that has already filled its \(k\)-heap still takes every warp barrier of the walk, so \(\tau\) can vary at launch without leaving a lane off a synchronizing instruction.

We record the device numbers in one process of \texttt{bench\_gpu\_batch} under \texttt{nsys} on the eleven stored frames of the 4096-molecule mW cubic dump.
The job script is \texttt{repro/elja\_gpu\_job.sh} of seams-core, which holds the record from commit \texttt{e1003cc1}.
The script \texttt{scripts/elja\_gpu\_job.sh} of the reproducibility package differs only in its output directory.
The script wrote neither the engine commit nor the NVIDIA driver version into the record, and the rerun of Section~\ref{sec:s-rerun} has no GPU node to supply them.
Its kernel table has the \texttt{emit\_basal} and \texttt{apply\_hc} kernels and no \texttt{zero\_i32}.
The walk ran once per frame with a synchronization after every launch (66 \texttt{knearest\_shells} launches and 312 synchronizations for 312 launches over the six calls).
That matches seams-core \texttt{b197e3be} and no later engine.
The job held both NVIDIA A100-PCIE-40\,GB of an Elja \texttt{gpu-2xA100} node and used one, with CUDA 12.2.
The probe gave \(W=2867712\) bytes and \(\lfloor A/W\rfloor=11672\) frames, so the request of 11 frames fits.
We time upload, compute and download on the host with \texttt{std::chrono::steady\_clock}, and each phase ends at a device synchronization.
One cold call, which includes the NVRTC compile, takes \(6394\,\mathrm{ms}\) of compute.
The best of five warm calls takes \(21.986\,\mathrm{ms}\) of compute for the eleven frames, with \(0.120\,\mathrm{ms}\) of upload and \(0.176\,\mathrm{ms}\) of download.
The batch labels \(45056\) of the \(45056\) molecules Ic over \(90112\) six-rings, the \(8192\) per frame of the host search.
The \texttt{nsys} kernel summary of the same process gives \(69.9\%\) of kernel time to the \(k\)-nearest walk, \(17.9\%\) to primitive hexagons and \(7.2\%\) to basal emission (Figure~\ref{fig:gpukernels}).

\begin{figure}[htbp]
\centering
\resizebox{\textwidth}{!}{\input{figures/gpu_resident}}
\caption{Device-resident ice-score batch. Coordinates upload once.
The $k$-nearest walk, the mutual graph, the hexagons and the
HC/DDC predicates enqueue on one stream. Labels come back.
CHILL stays on the host. A100-PCIE-40\,GB: best warm compute
$21.986\,\mathrm{ms}$ for the eleven stored 4096-molecule
cubic frames.}
\label{fig:gpuresident}
\end{figure}

\begin{figure}[htbp]
\centering
\resizebox{\textwidth}{!}{\input{figures/gpu_kernels}}
\caption{Share of device kernel time in the \texttt{nsys} summary of the
eleven-frame 4096-molecule cubic batch, one cold and five warm calls.
The Chebyshev-shell $k$-nearest walk is $69.9\%$; primitive
six-rings and basal emission are $17.9\%$ and $7.2\%$.}
\label{fig:gpukernels}
\end{figure}

\begin{figure}[htbp]
\centering
\resizebox{\textwidth}{!}{\input{figures/gpu_occupancy}}
\caption{Grid $\lceil nF/128\rceil$ of the per-atom kernels (mutual
graph, hexagon enumeration, row sort) against the 108 streaming
multiprocessors of the A100. Left: one frame against system size.
Right: the 4096-molecule cubic dump against batch width. For these
kernels $F=1$ launches fewer blocks than SMs and $F=11$ more; the ring
and pair kernels launch 16 and 128 times as many blocks.}
\label{fig:gpuoccupancy}
\end{figure}

\FloatBarrier

\begin{algorithm}[htbp]
\caption{Device-resident ice-score batch}
\label{alg:resident}
\begin{algorithmic}[1]
\Require frames $X_f$, boxes $L_f$, $f=1,\ldots,F$, with $F$ from Eq.~\ref{eq:resident}
\State upload $\{X_f\}$
\State enqueue on $Q$: bin, scan, scatter, $k$-NN walk, mutual $G$,
       hexagons, $I(v)$, $\beta$, HC, DDC, per-atom labels
\State wait on $Q$
\State download the per-atom HC and DDC labels
\end{algorithmic}
\end{algorithm}

A batch is one upload of \(F\) frames and one download of their labels.
The kernels are the \(k\)-nearest walk and Algorithms~\mref{alg:rings} and \mref{alg:cages} restricted to six-rings on the mutual graph, one thread per source or ring, over \(\lceil nF\tau/B\rceil\) blocks (Eq.~\ref{eq:blocks}).
The device numbers six-rings by atomic slot allocation; the predicates of Algorithm~\mref{alg:cages} do not read the ring order, so its labels do not depend on that numbering.
Both device paths allocate from a plan.
The plan holds \(16n\) six-rings and \(8\) basal pairs per ring in each frame, and \(32\) ring identifiers in each row of \(I\) on the offload path and \(16\) in the CUDA batch.
The kernels count every ring, row entry and basal pair, and the count covers those that do not fit.
A frame whose counts exceed a buffer is rerun with that buffer grown to the count.
The local rows of the DDC test hold \(32\) entries, so a frame with an atom in more than \(32\) six-rings is labelled by the host path on the same mutual graph.
No ring, row entry or pair is dropped.
\texttt{tests/check\_device\_host.cpp} compares the offload path with the host and with a brute-force evaluation of \(\beta\) and \(\varepsilon\).
The cases are frames 1, 74 and 148 of the nucleation deposit, the mW crystal and the bundled example, each on the cutoff and mutual four-nearest graphs.
They also include the Levi graphs of the projective planes of orders 2, 3 and 5, in which every vertex lies in 12, 54 and 375 six-rings.
We run the offload path on each case with the planned buffers, with every buffer at one entry, and with its OpenMP target regions executed on the host.
In all thirteen cases every run gives the per-atom HC and DDC labels of the host and of the brute-force evaluation.
The one-entry runs grow their buffers in up to three reruns, and the planes of orders 3 and 5 go to the host path.
Under two to eight random relabellings of each case the labels map back unchanged.
One atom of frame 1 lies in 18 six-rings, more than a row of the CUDA plan holds.
\texttt{tests/gpu\_kernels\_host.cpp} compiles the CUDA kernels for the host and runs them under the same rerun and host decisions.
On the planes of orders 2 and 3 and on ice Ih, with planned and with one-entry buffers, every frame matches the host labels.
Both programs exit nonzero on any difference, and \texttt{meson test} runs them, \texttt{check\_device\_host} on the mW crystal and the example on both graphs and on the three planes.
Eq.~\ref{eq:resident} keeps \(FW\le A\), is the largest such \(F\) at fixed \(W\) when the request exceeds it (Theorems~S30 and S31), and the grid covers every source (Theorem~S32).

\section{Costs in full}
\label{sec:s-costs}
\begin{table}[htbp]
\caption{\label{tab:complexity-full}\label{tab:complexity}Cost per frame.
Here \(n\) counts the molecules and \(\Delta\) is the largest degree of \(G\) (of \(G\cup G'\) for the update).
The largest ring size is \(D\), with \(\ell=\lfloor D/2\rfloor\) and \(\ell'=\max(\ell-1,1)\), and \(b_r\le 1+\Delta\sum_{i<r}(\Delta-1)^i\) is the largest \(r\)-hop ball, and in the ring-update row \(r=\max(\lfloor(\ell-2)/2\rfloor,0)\) is the radius of Eq.~\mref{eq:hopbound}.
The stored shortest paths number \(P_s\) and the path tuples inspected from source \(s\) number \(K_s\), rejected ones included, with \(P=\max_s P_s\) and \(K=\max_s K_s\).
The six-rings number \(R_6\), \(d_I=\max_v|I(v)|\), and \(B\) counts the ordered basal pairs that pass \(\beta\). The bonds that change between frames form \(\Delta E\), \(H\le R_6^2\) counts the comparisons within equal-hash groups of the ring diff, and \(C\) holds the dirty rings of Algorithm~\mref{alg:affupdate}. Each molecule makes \(k\) nominations, \(\iota\) is the largest number of molecules that nominate one molecule, and the \(k\)-nearest search examines \(m\) points. At fixed \(\Delta\) and \(D\) every quantity but \(n\), \(R_6\), \(B\), \(H\), \(|C|\), \(m\) and \(|\Delta E|\) is bounded, and \(R_6,B=O(n)\).}
\centering
\footnotesize
\setlength{\tabcolsep}{3pt}
\begin{tabular}{L{2.2cm}L{3.3cm}L{4.6cm}L{2.7cm}}
stage & d-SEAMS 1 & d-SEAMS 2.0 & guarantee\\
\hline
primitive rings & \(O\bigl(n\bigl[1+\Delta\sum_{j=0}^{D-2}(\Delta-1)^j\bigr]\bigr)\) walk prefixes, then a test per walk & \(O\bigl(n\,\Delta\,b_{\ell'}+\sum_s(\Delta b_\ell+\ell P_s+D^2K_s)\bigr)\) & primitive rings, each once (Prop.~P1)\\
next frame, rings & rebuild & \(O\bigl(n[1+\Delta\log(\Delta+1)]+nDK+|\Delta E|\,\Delta\,b_{\ell'}^2+|\Delta E|\,b_{r}(\Delta b_\ell+\ell P+D^2K)\bigr)\) & equals rebuild (Prop.~P4)\\
HC search & \(O(R_6^2)\) & \(O\bigl(R_6\,\Delta\,d_I(\Delta+\log(\Delta d_I))+B\,d_I\bigr)\) & flags of Eq.~\mref{eq:ice}\\
DDC search & \(O(R_6^2)\) & \(O(R_6\,d_I^6\log(d_I+1))\) & flags of Eq.~\mref{eq:ice}\\
next frame, cage flags & rebuild & \(O\bigl(n[1+\Delta\log(\Delta+1)]+R_6\log(R_6+1)+H\bigr)\) and the HC and DDC tests of the rings near \(C\) & equals Algorithm~\mref{alg:cages}\\
cage objects & visiting order & position-keyed order, \(O(n\log n+R_6\log R_6)\) & function of the ring set, neighbour rows and positions\\
seeded and completed labels & -- & \(O(m\log k+n\,d_I)\) after two ice scores on degrees \(\le k\) and \(\le k+\iota\) & least fixed point (Prop.~P5)\\
device batch & -- & \(F=\min(F_{\mathrm{req}},\lfloor A/W\rfloor)\) frames per upload & fits; maximal at fixed \(W\)\\
\end{tabular}
\end{table}
\subsection{Ring construction}
The rows cost \(O(n\,\Delta\,b_{\ell'})\) time and \(O(n\,b_{\ell'})\) memory.
Source \(s\) costs \(O(\Delta\,b_\ell)\) for its sweep, \(O(\ell P_s)\) to store its \(P_s\) paths and \(O(D^2)\) lookups for each of the \(K_s\) path tuples it inspects, rejected ones included.
The pass is \(O\bigl(n\,\Delta\,b_{\ell'}+\sum_s(\Delta b_\ell+\ell P_s+D^2K_s)\bigr)\), linear in \(n\) at fixed \(\Delta\) and \(D\).
The walk enumeration of d-SEAMS 1 grows a search tree rooted at each source over the walks of length below \(D\).
The tree has up to \(1+\Delta\sum_{j=0}^{D-2}(\Delta-1)^j\) nodes, which is \(485\) at \(\Delta=4\) and \(D=6\), where the sweep visits at most \(b_3=53\) vertices.
\subsection{Ring update}
Here \(\Delta\) is the largest degree of \(G\cup G'\).
The edge diff sorts and compares every neighbour row in \(O(n[1+\Delta\log(\Delta+1)])\), and a frame with \(\Delta E\ne\emptyset\) builds every row of \(G\cup G'\) in \(O(n[1+\Delta\log(\Delta+1)])\).
The row update recomputes \(|U|\le 2|\Delta E|\,b_{\ell'}\) rows at \(O(\Delta\,b_{\ell'})\) each.
The enumeration runs from \(|T|\le 2|\Delta E|\,b_{r}\) sources, with \(b_0=1\), at the cost of one source of Algorithm~\mref{alg:rings} each, and the stored rings that meet \(T\) are dropped by one pass over the stored list.
The local recomputation is \(O\bigl(|\Delta E|\,\Delta\,b_{\ell'}^2+|\Delta E|\,b_{r}(\Delta b_\ell+\ell P+D^2K)\bigr)\), linear in \(|\Delta E|\) at fixed \(\Delta\) and \(D\).
The returned list is rebuilt on every frame, \(\Delta E=\emptyset\) included, by copying every stored ring, \(O\bigl(D\sum_s|\mathcal{R}[s]|\bigr)\) with \(|\mathcal{R}[s]|\le K\), so a frame costs \(O(n[1+\Delta\log(\Delta+1)]+nDK)\) for the diff and the output, plus the local recomputation.
We prove that a source outside \(S\) reads the same edges on both frames (Proposition~P2), that a row outside \(U\) is unchanged (Proposition~P3), and that the output equals a rebuild (Proposition~P4).
The Lean statements S14--S16 check the ball-agreement step and the carry-or-rebuild case split for an abstract enumerator that is assumed local; they do not model Algorithm~\mref{alg:rings}.
\subsection{Cage search}
The index holds \(6R_6\) ring identifiers in \(n\) rows and costs \(O(n+R_6)\) time and memory to build.
A ring draws at most \(2\Delta d_I\) basal candidates, sorts them, and tests each against the neighbour rows of its atoms in \(O(\Delta)\) at fixed ring size.
Each of the \(B\) passing ordered pairs then reads the rows of its twelve atoms for its prismatic rings, so the HC pass is \(O\bigl(R_6\,\Delta\,d_I(\Delta+\log(\Delta d_I))+B\,d_I\bigr)\).
An equatorial test forms the six triplet intersections from ascending rows of length at most \(d_I\) and enumerates at most \(d_I^3\) chains of alternate peripherals for each parity.
It tests at most \(d_I^6\) pairs of chains and sorts the peripherals of the passing pairs, so the DDC pass is \(O(R_6\,d_I^6\log(d_I+1))\).
Beyond the index and the returned lists, the HC flag sweep uses \(O(R_6+\Delta d_I+Bd_I)\) working memory and the DDC flag sweep \(O(R_6d_I+d_I^6)\).
In d-SEAMS 1 we tested \(\beta\) on every pair of six-rings and scanned the whole ring list for every atom of a candidate equator, \(O(R_6^2)\) for both.
On ice Ic \(R_6=2n\) and every molecule lies on twelve six-rings (\(8192\) six-rings on the 4096-molecule frame), so both passes are linear in \(n\) where the scans were quadratic.
Theorem~S29 states that \(R\in I(v)\) exactly when \(R\) is stored and \(v\in R\), which is the property of the stored index that the two passes use.

\subsection{Cage update}
\label{sec:s-costs-cage}
Every update of Algorithm~\mref{alg:affupdate} sorts the neighbour rows in \(O(n[1+\Delta\log(\Delta+1)])\).
An incremental update adds the ring diff, \(O(R_6\log(R_6+1)+H)\), the index, \(O(n+R_6)\), and the HC and DDC tests of the rings near \(C\).
The code also computes the closure of the changed rows before matching rings and returns the batch pass at once when that closure is already above \(\lfloor\phi R_6\rfloor\).
The Catch2 case ``AffiliationUpdater tracks synthetic topology churn exactly'' and the trajectory benchmarks of Section~\ref{sec:s-incremental} compare it with Algorithm~\mref{alg:cages} on every frame.

\section{Cage-object order and relabelling}
\label{sec:s-invariance}
The d-SEAMS 1 search claimed rings as it accepted cage objects, in the order in which it visited them.
d-SEAMS 2.0 fixes both orders.
Each atom is ranked by its position wrapped into the box, compared by \(x\), then \(y\), then \(z\), and both searches visit rings in the order of their sorted atom ranks.
\texttt{ring::positionRank} rejects a frame in which two atoms share a wrapped position, so every rank is fixed by the positions alone.
Hexagonal cages are claimed over basal pairs in that order, and the search skips a ring that an earlier cage claimed as prismatic.
Equatorial candidates skip HC rings and the peripherals of earlier double-diamond cages, and a cage takes the passing choice of \(R'_k\) whose ranks in that order are least.
A rank moves with its atom when the atoms are relabelled.
The cage objects are therefore a function of the ring set, the neighbour rows and the positions.
They do not depend on the atom labels, the order of the ring list or where each ring starts and which way it runs.
The bulk topological analysis, the topological unit matching and \texttt{seams cages -{}-signature hc} or \texttt{ddc} rank the atoms of every frame they analyse.
The reference HC and DDC of the unit matching come from the fixed files \texttt{templates/hc.xyz} and \texttt{templates/ddc.xyz}, whose atom order sets the order of the reference points.
The per-ring flags of Eqs.~\mref{eq:ddc-aff} and \mref{eq:ice} claim no rings and read neither the positions nor the labels.
\texttt{tests/check\_invariance.cpp} reverses the ring list, rotates every ring by one position, reverses every ring, rotates every ring by two positions and reverses it, and relabels the atoms of the frame at random and enumerates its rings afresh.
We run it on frames 1, 74 and 148 of the nucleation deposit, the mW crystal and the bundled example.
Every variant leaves the six-ring set, the HC and DDC ring sets, the per-atom affiliation and every cage, as a set of rings, unchanged once the labels are mapped back.
Frame 148 keeps its 224 cages and the mW crystal its 2040.
The program exits nonzero on any difference, and \texttt{meson test} runs it on the mW crystal and the example.
The Catch2 case ``Cage objects are invariant under atom relabelling'' relabels the mW crystal at random three times, rebuilds its neighbour list and rings, and requires every cage back.

\section{Primitive rings}
\label{sec:org0f685be}

d-SEAMS identifies ice Ih and ice Ic from hexagonal and
double-diamond cages on the primitive six-rings of the neighbour
graph \citep{goswami2020dseams}. Those rings are not the $5^{12}$
and $5^{12}6^{2}$ clathrate cages of Jacobson et al.\ \citep{jacobson2010amorphous}.
Let $G=(V,E)$ be undirected and simple, $V=\{0,\ldots,n-1\}$. A
\emph{walk} of length $k$ is a sequence $w=(u_0,\ldots,u_k)$ with
$u_{t-1}u_t\in E$. Set
\begin{equation}
d_G(u,v)
  \;=\;
  \min\bigl\{\,k: \text{there is a walk of length $k$ from $u$ to $v$}\,\bigr\},
\label{eq:si-dg}
\end{equation}
or $d_G(u,v)=\infty$ if no such walk exists. Let
$R=(v_0,\ldots,v_{n_R-1})$ be a cycle of length $n_R\ge 3$ with
distinct vertices. We write
\begin{equation}
d_R(v_i,v_j)
  \;=\;
  \min\bigl(|i-j|,\,n_R-|i-j|\bigr)
\label{eq:si-dr}
\end{equation}
for the shorter arc. The cycle is a primitive (shortest-path) ring
in the sense of Franzblau \citep{franzblau1993computation} when
\begin{equation}
d_G(v_i,v_j)=d_R(v_i,v_j)
\qquad\text{for all }i,j.
\label{eq:si-franzblau}
\end{equation}
Yuan's primitive ring is the same object; strong rings are a
stricter subclass \citep{yuan2002rings}. Two cycles are the same
ring when one is a rotation or a reversal of the other. Ice Ih and
ice Ic are built from six-rings, so the cage searches take $n_R=6$
and $D=6$.

\paragraph{No shortcut.}
We take any walk $w$ that starts at $v_i$ and ends at $v_j$. By~\eqref{eq:si-dg},
$d_G(v_i,v_j)\le\mathrm{len}(w)$. Substituting~\eqref{eq:si-franzblau}
gives $d_R(v_i,v_j)=d_G(v_i,v_j)\le\mathrm{len}(w)$.
Since the arc of $R$ is itself a walk, $d_G\le d_R$ always holds, so
Eq.~\eqref{eq:si-franzblau} fails exactly when some pair has
$d_G(v_i,v_j)<d_R(v_i,v_j)$, a strict shortcut.

\begin{theorem}[Franzblau ring has no shortcut]\label{thm:franzblau-noshort}
If $R=(v_0,\ldots,v_{n-1})$ satisfies Eq.~\eqref{eq:si-franzblau}, then
every walk in $G$ between $v_i$ and $v_j$ has length at least $d_R(v_i,v_j)$.
\leanname{franzblau_no_shortcut}{Ring}
\end{theorem}

\begin{theorem}[Ring-size bound]\label{thm:ring-size}
A list of distinct vertices of $G$, each less than $n$, has length at
most $n$.
\leanname{ring_size_bound}{Ring}
\end{theorem}

\section{Shortest-path ring construction}
\label{sec:org0890993}

Algorithm~\mref{alg:rings} of the main text fixes $\ell=\lfloor D/2\rfloor$ and
$\ell'=\max(\ell-1,1)$. For each source $s$ it runs a breadth-first
sweep of $G$ to depth $\ell$, records the level $L(v)=d_G(s,v)$ and drops
every vertex $v<s$. At each kept $v$ it stores the set $\Pi(v)$ of
shortest $s$--$v$ paths of $G$ whose vertices are all at least $s$, each
as an explicit vertex list. We write a shortest path of length $t$ between $s$ and
$p$ as $P=(s,a_1,\ldots,a_{t-1},p)$, and its interior is
$(a_1,\ldots,a_{t-1})$. Two candidate types are assembled as ordered
cycles.
\begin{itemize}
\item Even: for $p$ with $L(p)=t$ and $3\le 2t\le D$, and each
  unordered pair $\{P,P'\}$ of distinct paths of $\Pi(p)$ whose
  interiors are disjoint, the cycle
  $C=(s,P\text{ interior},p,\text{reverse of }P'\text{ interior})$ of
  length $2t$.
\item Odd: for each edge $pq$ with $p<q$, $L(p)=L(q)=t$ and
  $2t+1\le D$, and each $P\in\Pi(p)$ and $Q\in\Pi(q)$ whose interiors
  are disjoint and avoid $p$ and $q$, the
  cycle $C=(s,P\text{ interior},p,q,\text{reverse of }Q\text{ interior})$ of
  length $2t+1$.
\end{itemize}
The Keep step emits $C$ when no pair of its vertices has a strict
shortcut. It reads shortcuts in the rows $d_{\ell'}(u,\cdot)$ of breadth-first
distances out to $\ell'$ hops, computed once for the frame. An absent entry
counts as greater than $\ell'$.
A candidate with a repeated vertex is discarded as it is assembled, and
every other candidate is kept if and only if it satisfies
Eq.~\mref{eq:franzblau} of the main text. No closed walk that is not a
shortest-path pair is generated.

\paragraph{Lengths.}
In the even case the two interiors have $t-1$ entries each. The
cycle lists $s$, $p$ and both interiors, so it has $2+2(t-1)=2t$ entries. In
the odd case it lists $s$, $p$, $q$ and two interiors, so it has
$3+2(t-1)=2t+1$ entries. Theorems~\ref{thm:yuan-even}
and~\ref{thm:yuan-odd} state that arithmetic for the Lean lists; they
do not state that the entries are distinct, that they form a path,
or that the interiors are disjoint. Distinctness is the disjointness
test of the construction above.

\begin{theorem}[Even candidate list length]\label{thm:yuan-even}
If $\ell\ge 2$, $\mathtt{int1.length}=\ell-1$ and
$\mathtt{int2.length}=\ell-1$, then
$(\mathtt{yuanEvenVerts}\;s\;p\;\mathtt{int1}\;\mathtt{int2}).\mathtt{length}=2\ell$.
\leanname{yuan_even_size}{Yuan}
\end{theorem}

\begin{theorem}[Odd candidate list length]\label{thm:yuan-odd}
If $\ell\ge 1$, $\mathtt{intP.length}=\ell-1$ and
$\mathtt{intQ.length}=\ell-1$, then
$(\mathtt{yuanOddVerts}\;s\;p\;q\;\mathtt{intP}\;\mathtt{intQ}).\mathtt{length}=2\ell+1$.
\leanname{yuan_odd_size}{Yuan}
\end{theorem}

\begin{paperprop}[Algorithm~\mref{alg:rings} emits every primitive ring once]\label{prop:alg1}
For $3\le|R|\le D$, Algorithm~\mref{alg:rings} emits every primitive ring of $G$
exactly once, from its lowest vertex, and emits nothing else.
\end{paperprop}

\begin{proof}
\emph{Soundness.} An emitted cycle has distinct vertices (disjoint
interiors, and $p\ne q$ in the odd case) and consecutive vertices
adjacent, so it is a cycle of $G$. Every pair of its vertices is at
cycle distance at most $\ell$, so a strict shortcut has length at
most $\ell-1\le\ell'$ and is visible in the rows $d_{\ell'}$. Keep therefore
emits exactly the candidates that satisfy
Eq.~\eqref{eq:si-franzblau}.

\emph{Completeness.} We take a primitive ring $R$ with $|R|=m\le D$ and
lowest vertex $s$. Every vertex of $R$ is at least $s$, so no vertex
of $R$ is dropped. If $m=2t$, let $p$ be the vertex of $R$ opposite
$s$. The two arcs between $s$ and $p$ have length $t=d_R(s,p)=d_G(s,p)$,
because a primitive ring has no shortcut. Both arcs are therefore
shortest $s$--$p$ paths of $G$ through vertices at least $s$, so they
lie in $\Pi(p)$ and $L(p)=t$. They are distinct,
their interiors are disjoint because $R$ is a cycle, so the even
step assembles $R$. If $m=2t+1$, let $\{p,q\}$ be the edge of $R$
opposite $s$, labelled so that $p<q$. The arcs that join $s$ to $p$ and to
$q$ without that edge have length $t=d_G(s,p)=d_G(s,q)$ and lie in
$\Pi(p)$ and $\Pi(q)$ with $L(p)=L(q)=t$. The odd step then
assembles $R$. Keep passes because $R$ is primitive.

\emph{Uniqueness.} A source $s'>s$ never emits $R$, since it drops
$s$, and a source $s'<s$ is not a vertex of $R$. From $s$, the cycle closes over the vertex (even) or edge
(odd) of $R$ opposite $s$, which is unique, so only one $p$, or one
$p<q$, produces $R$. In the even case the pair of arcs is unordered
and is enumerated once. In the odd case the path ending at $p$ and
the path ending at $q$ are distinguished by their endpoints, and
$p<q$ fixes the orientation. Each ring is therefore assembled once,
and no canonical key is needed.
\end{proof}

\section{Hexagonal and double-diamond cages}
\label{sec:orgfd88a9d}

A hexagonal cage (HC) is the unit of ice Ih. d-SEAMS assembles one
from basal and prismatic six-rings \citep{goswami2020dseams}. The cage is two
basal hexagons joined by three prismatic six-rings. A double-diamond
cage (DDC) is the unit of ice Ic. It has seven primitive six-rings on
fourteen vertices, one equatorial ring and six peripheral rings
\citep{goswami2020dseams}. We write
$I(v)=\{R:v\in R\}$ for the inverted index, each row in ascending ring
index. A candidate six-ring $R=(m_0,\ldots,m_5)$ is equatorial when
every vertex of $R$ lies in at least three other six-rings,
\begin{equation}
\bigl|I(m_k)\setminus\{R\}\bigr|\ge 3
\qquad k=0,\ldots,5,
\label{eq:si-ddc1}
\end{equation}
every consecutive triplet lies in another six-ring,
\begin{equation}
T_k=I(m_k)\cap I(m_{k+1})\cap I(m_{k+2})\setminus\{R\}\ne\emptyset
\qquad k=0,\ldots,5
\label{eq:si-ddc2}
\end{equation}
(indices modulo 6). The last condition asks for six distinct peripherals
$R'_k\in T_k$ and an index $j$ such that $R'_0$, $R'_2$, $R'_4$ share
a vertex and $R'_1$, $R'_3$, $R'_5$ share a vertex. For the same choice, $R'_j$ shares at
least three vertices with each of $R'_{j+2}$ and $R'_{j+4}$, and
$R'_{j+1}$ with each of $R'_{j+3}$ and $R'_{j+5}$. With $j=2$ these are
the alternating-intersection conditions of d-SEAMS
\citep{goswami2020dseams}. Write $\varepsilon(R)$ for the
conjunction, and $P(R)$ for the union of $\{R'_0,\ldots,R'_5\}$ over
all passing choices. The same cage criteria classify confined and nanoribbon
ice \citep{goswami2021tum,goswami2020nanoribbon}.

\paragraph{What the Lean structures store.}
The Lean structure \texttt{HCCage} stores two lists \texttt{basal1},
\texttt{basal2}, each of length six, and \texttt{allVerts} as their
concatenation. The structure \texttt{DDCCage} stores the equatorial
list, the list of peripherals and \texttt{allVerts}, together with
fields \texttt{vertexCount}, \texttt{periphCount}, \texttt{cond1} and
\texttt{cond2} that \emph{assume} the corresponding properties.
Theorems~\ref{thm:hc12}--\ref{thm:ddc2} are therefore statements
about stored lists and fields. They do not derive cage cardinalities
or the DDC conditions from adjacency.

\begin{theorem}[HC list length]\label{thm:hc12}
For an \texttt{HCCage}, the stored list \texttt{allVerts}, defined as
\texttt{basal1 ++ basal2} with each basal list of length six, has
list length 12.
\leanname{hc_has_12_vertices}{Topology}
\end{theorem}

\begin{theorem}[DDC list length]\label{thm:ddc14}
For a \texttt{DDCCage}, whose \texttt{vertexCount} field is
$\mathtt{allVerts.length}=14$, the stored list \texttt{allVerts} has
length 14.
\leanname{ddc_has_14_vertices}{Topology}
\end{theorem}

\begin{corollary}[DDC ring count]\label{thm:ddc7}
For a \texttt{DDCCage}, whose \texttt{periphCount} field is
$\mathtt{peripheral.length}=6$, $\mathtt{peripheral.length}+1=7$.
\leanname{ddc_seven_rings}{Topology}
\end{corollary}

\begin{theorem}[DDC first condition field]\label{thm:ddc1}
For a \texttt{DDCCage}, the stored field
\[\mathtt{cond1}:\mathtt{ddcCond1}\;\mathtt{equatorial}\;\mathtt{peripheral}\]
yields $\mathtt{ddcCond1}\;\mathtt{equatorial}\;\mathtt{peripheral}$.
\leanname{ddc_equatorial_vertex_count}{Topology}
\end{theorem}

\begin{theorem}[DDC second condition field]\label{thm:ddc2}
For a \texttt{DDCCage}, the stored field
\[\mathtt{cond2}:\mathtt{ddcCond2}\;\mathtt{equatorial}\;\mathtt{peripheral}\]
yields $\mathtt{ddcCond2}\;\mathtt{equatorial}\;\mathtt{peripheral}$.
\leanname{ddc_consecutive_triplet}{Topology}
\end{theorem}

\section{Affiliation and the ice score}
\label{sec:org36e36f2}

A ring is HC-affiliated when it is a basal ring of an ordered pair
that passes the basal test $\beta$, or a prismatic ring of such a
pair. We restrict equatorial candidacy to rings that are not
HC-affiliated, as the engine and the baseline commit of the
main-text benchmarks do; the 2020
release tested every ring as an equator. A ring is DDC-affiliated
when it is such an equator or a ring of $P(R_e)$ for one,
\begin{equation}
\begin{aligned}
\alpha_{\mathrm{DDC}}(R)\iff{}&
\bigl[\neg\alpha_{\mathrm{HC}}(R)\wedge\varepsilon(R)\bigr]\\
&\lor\ \exists\,R_e:\ \neg\alpha_{\mathrm{HC}}(R_e)\wedge\varepsilon(R_e)\wedge R\in P(R_e).
\end{aligned}
\label{eq:si-ddc-aff}
\end{equation}
When a peripheral is itself HC-affiliated, that ring carries both flags,
which is the stacking-fault case. Molecule $v$ is ice Ic when a six-ring through $v$ is
DDC-affiliated and ice Ih when a six-ring through $v$ is
HC-affiliated,
\begin{equation}
\mathrm{Ic}(v)\iff \exists\, R\ni v,\;\alpha_{\mathrm{DDC}}(R),
\qquad
\mathrm{Ih}(v)\iff \exists\, R\ni v,\;\alpha_{\mathrm{HC}}(R).
\label{eq:si-ice}
\end{equation}
A molecule with neither label is unassigned.

\paragraph{The Lean predicates.}
Lean's \texttt{alphaHC} is a simplified basal-pair fragment. The predicate
\texttt{alphaHC adj rings R} holds when some stored $R'\ne R$ has an
atom adjacent to an atom of \texttt{R.take 2}. It omits the triplet
conditions of $\beta$ and the prismatic rings. Lean's
\texttt{alphaDDC R P} is \texttt{hexagonal R} together with the
abstract conditions \texttt{ddcCond1} and \texttt{ddcCond2} on a given
peripheral list $P$; it does not compute $P$ from adjacency and has no
HC guard. Theorems~\ref{thm:ice-ic}--\ref{thm:aff-local} are about
those predicates, with the stored rings and the peripheral map
\texttt{periphOf} fixed.

\begin{theorem}[Ic witness]\label{thm:ice-ic}
For the Lean predicate
\[\mathtt{alphaDDC}\;R\;P:=\mathtt{hexagonal}\;R\wedge\mathtt{ddcCond1}\;R\;P\wedge\mathtt{ddcCond2}\;R\;P,\]
a stored ring $R$ containing $v$ and satisfying
$\mathtt{alphaDDC}\;R\;(\mathtt{periphOf}\;R)$ witnesses
$\mathtt{iceIc}\;\mathtt{adj}\;\mathtt{rings}\;\mathtt{periphOf}\;v$.
\leanname{ice_ic_of_ring}{IceScore}
\end{theorem}

\begin{theorem}[Ih witness]\label{thm:ice-ih}
For the Lean predicate \texttt{alphaHC}, defined by existence of a
stored $R'\ne R$ with a true adjacency between an atom of
\texttt{R.take 2} and an atom of $R'$, such a stored ring $R$
containing $v$ witnesses $\mathtt{iceIh}\;\mathtt{adj}\;\mathtt{rings}\;v$.
\leanname{ice_ih_of_ring}{IceScore}
\end{theorem}

\begin{theorem}[Equal adjacency, equal score]\label{thm:ice-invariant}
For fixed \texttt{rings} and \texttt{periphOf}, equal adjacency
functions give equivalent \texttt{iceIc} and \texttt{iceIh}
predicates.
\leanname{ice_score_graph_invariant}{IceScore}
\end{theorem}

\begin{theorem}[Locality of the simplified HC predicate]\label{thm:aff-local}
For the simplified Lean predicate \texttt{alphaHC}, agreement of
\texttt{adj} and \texttt{adj'} on $\mathtt{R.take\ 2}\times R'$ for
every stored $R'$ makes $\mathtt{alphaHC}\;\mathtt{adj}\;\mathtt{rings}\;R$
and $\mathtt{alphaHC}\;\mathtt{adj'}\;\mathtt{rings}\;R$ equivalent.
\leanname{affiliation_locality}{IceScore}
\end{theorem}

\section{Locality and the incremental update}
\label{sec:org7096e5e}

\begin{figure}[htbp]
\centering
\resizebox{0.92\textwidth}{!}{\input{figures/hopbound}}
\caption{Hop-bound update at $D=6$, where $r=0$ and $T$ is the set
of endpoints of $\Delta E$. A ring through an atom of $T$ is dropped
and enumerated again on the next frame. A ring that avoids $T$ keeps
its stored entry. The output equals a full rebuild.}
\label{fig:hopbound}
\end{figure}

Write $G\cup G'$ for the union of the neighbour graphs of two
consecutive frames, $\Delta E=E\mathbin{\triangle}E'$ for the
changed edges, $X$ for their endpoints, $\Delta$ for the largest
degree of $G\cup G'$, and
\begin{equation}
B_r(s)=\bigl\{u:d_{G\cup G'}(s,u)\le r\bigr\}.
\label{eq:si-ball}
\end{equation}
We take distances to $X$ on $G\cup G'$.
The source bound at $s$ is
\begin{equation}
d_{G\cup G'}(s,X)>2\ell-1.
\label{eq:si-hopbound}
\end{equation}

\paragraph{Ball agreement.}
Suppose $uv\in\Delta E$. Then~\eqref{eq:si-hopbound} gives
$u,v\notin B_{2\ell-1}(s)$. Contrapositively, $G$ and $G'$ agree on
every edge with both endpoints in $B_{2\ell-1}(s)$.

\begin{paperprop}[A source reads a bounded ball]\label{prop:local}
The set of rings that Algorithm~\mref{alg:rings} emits from $s$ is a function of the
subgraph induced on $B_{2\ell-1}(s)$. Hence, if $G$ and $G'$ agree on
that subgraph, Algorithm~\mref{alg:rings} emits the same rings from $s$ on both.
\end{paperprop}

\begin{proof}
The sweep from $s$ reads only edges of $G$ with an endpoint at level
at most $\ell-1$, so both endpoints lie within $\ell$ of $s$. Every
candidate vertex lies within $\ell$ of $s$. Keep asks, for a pair
$c_i,c_j$ of candidate vertices, whether $d_G(c_i,c_j)<d_C(c_i,c_j)$.
Since $d_C\le\ell$, that question is the existence of a path of length
at most $\ell-1$ from $c_i$; every vertex of such a path lies within
$\ell+(\ell-1)=2\ell-1$ of $s$, and so does every edge it uses. The
same argument applies on $G'$. Distances on $G$ and on $G'$ are at
least those on $G\cup G'$, so every vertex the computation reaches on
either frame lies in $B_{2\ell-1}(s)$, where the two frames have the
same edges.
\end{proof}

\paragraph{Distance rows.}
Algorithm~\mref{alg:rings} reads the rows $d_{\ell'}(u,\cdot)$ of radius
$\ell'=\max(\ell-1,1)$. Algorithm~\mref{alg:hopbound} stores the rows of $G$ and
recomputes, on $G'$, the rows of
\begin{equation}
U=\bigl\{u:d_{G\cup G'}(u,X)\le\ell'\bigr\},
\label{eq:si-rows}
\end{equation}
and carries every other row.

\begin{paperprop}[Carried rows are exact]\label{prop:rows}
For $u\notin U$, $d_{\ell'}^{G}(u,\cdot)=d_{\ell'}^{G'}(u,\cdot)$.
\end{paperprop}

\begin{proof}
We take a shortest path of $G$ of length at most $\ell'$ from $u$. Each
of its edges has an endpoint at $G$-distance at most $\ell'-1$ from
$u$, hence at $(G\cup G')$-distance at most $\ell'-1$. An edge of
$\Delta E$ has both endpoints in $X$, at distance greater than $\ell'$
from $u$ since $u\notin U$. The path therefore uses no changed edge and
is a path of $G'$. The argument is symmetric in $G$ and $G'$, so the
two rows coincide. The proof needs only $d_{G\cup G'}(u,X)\ge\ell'$;
the engine refreshes the larger set $U$.
\end{proof}

\paragraph{The set $T$.}
The engine does not test sources against Eq.~\eqref{eq:si-hopbound}.
It drops and enumerates again the rings that meet
\begin{equation}
T=\bigl\{v:d_{G\cup G'}(v,X)\le r\bigr\},\qquad r=\max\bigl(\lfloor(\ell-2)/2\rfloor,0\bigr),
\label{eq:si-T}
\end{equation}
so $T=X$ for $D\le7$.

\begin{paperprop}[The update equals a rebuild]\label{prop:alg2}
(i) A cycle of at most $D$ vertices that is primitive in exactly one of
$G$ and $G'$ has a vertex in $T$.
(ii) Recompute the rows of $U$, keep the stored rings of $G$ that
avoid $T$, and run Algorithm~\mref{alg:rings} on $G'$ from every
$s\in T$ with only the vertices $t\in T$, $t<s$, dropped. The result is
every primitive ring of $G'$ with at most $D$ vertices, each exactly
once, and nothing else.
\end{paperprop}

\begin{proof}
(i) Let $Z$ be such a cycle. If $Z$ holds an edge of $\Delta E$, its
endpoints are in $X\subseteq T$. Otherwise $Z$ is a cycle of both
graphs, primitive in one and not in the other; call the second $H$.
In $H$ some pair $a,b\in Z$ has a path $p$ with
$|p|<d_Z(a,b)\le\lfloor|Z|/2\rfloor\le\ell$, so $|p|\le\ell-1$. If $p$
held no edge of $\Delta E$, it would be a path of the other graph as
well, and $Z$ would be primitive in neither. So $p$ holds an edge
$xy\in\Delta E$, the $i$-th of its $|p|$ edges. Along $p$, $a$ is
$i-1$ edges from $x$ and $b$ is $|p|-i$ edges from $y$. The smaller of
the two is at most $\lfloor(|p|-1)/2\rfloor\le\lfloor(\ell-2)/2\rfloor\le r$.
Since $p$ is a path of $G\cup G'$ and $x,y\in X$, $a$ or $b$ lies in $T$.

(ii) After the row update every row is that of $G'$ (Proposition~P3 for
$u\notin U$, recomputation for $u\in U$). Every edge of $\Delta E$ has
its endpoints in $T$, so a cycle that avoids $T$ has the same edges in
both graphs. By (i) it is primitive in $G$ if and only if it is
primitive in $G'$. The kept rings are therefore the primitive rings of
$G'$ that avoid $T$. A primitive ring of $G'$ that meets $T$ has one
lowest vertex $s$ in $T$. The proof of Proposition~P1 holds with
``vertices at least $s$'' replaced by ``vertices other than the $t\in T$
with $t<s$''. It uses the vertex order only to choose the source
that emits a ring, and a strict shortcut still has length at most
$\ell-1\le\ell'$ and lies in the rows $d_{\ell'}$ of $G'$. The run from
$s$ therefore emits that ring once, the run from any other vertex of
$T$ does not emit it, and no run emits a cycle that is not a primitive
ring of $G'$. The kept and the emitted rings are disjoint and together
hold every primitive ring of $G'$ with at most $D$ vertices.
\end{proof}

The row update recomputes at most $|U|\le 2|\Delta E|\,b_{\ell'}$
rows at $O(\Delta\,b_{\ell'})$ each, where
$b_r\le 1+\Delta\sum_{i<r}(\Delta-1)^i$ bounds an $r$-ball, and the
enumeration runs from $|T|\le 2|\Delta E|\,b_r$ sources, with $b_0=1$.
The engine computes $U$ and $T$ by multi-source sweeps on $G\cup G'$.

\paragraph{The source bound.}
Proposition~P2 gives a coarser test that the engine does not use.

\begin{papercor}[Source bound]\label{cor:source}
A source $s$ that satisfies Eq.~\eqref{eq:si-hopbound} emits the same
rings on $G$ and on $G'$.
\end{papercor}

\begin{proof}
By ball agreement $G$ and $G'$ agree on the subgraph induced on
$B_{2\ell-1}(s)$, and Proposition~P2 applies.
\end{proof}

\paragraph{The Lean statements.}
Lean proves the ball-agreement step and two conditional statements
about an abstract per-source enumerator
$\mathtt{full}:\mathtt{Graph}\to\mathbb{N}\to\mathtt{RingSet}$.
They concern the source bound of Corollary~\ref{cor:source} and model
neither Algorithm~\mref{alg:rings} nor the set $T$. The locality of the
actual enumerator, Proposition~P2, is a hypothesis of the Lean
statements, and we prove it and Proposition~P4 only here.

\begin{theorem}[Source bound implies ball agreement]\label{thm:hop-ball}
If Eq.~\eqref{eq:si-hopbound} holds at $s$, then $G$ and $G'$ agree on
every edge whose endpoints both lie in $B_{2\ell-1}(s)$.
\leanname{hopbound_implies_ball}{HopBound}
\end{theorem}

\begin{theorem}[Conditional invariance]\label{thm:hop-inv}
If an abstract enumerator \texttt{full} returns equal ring sets from
$s$ on any two graphs that agree on $B_{2\ell-1}(s)$, then a source
that satisfies Eq.~\eqref{eq:si-hopbound} has
$\mathtt{full}\;G\;s=\mathtt{full}\;G'\;s$.
\leanname{hopbound_invariance}{HopBound}
\end{theorem}

\begin{theorem}[Conditional carry-or-rebuild]\label{thm:incr}
Let \texttt{full} be an abstract per-source enumerator and assume
$\mathtt{full}\;G\;s=\mathtt{full}\;G'\;s$ for every source satisfying
Eq.~\eqref{eq:si-hopbound}. For each source, an output equal to
$\mathtt{full}\;G\;s$ in the source-bound branch and to
$\mathtt{full}\;G'\;s$ otherwise equals $\mathtt{full}\;G'\;s$.
\leanname{incremental_eq_rebuild}{HopBound}
\end{theorem}

\section{Seeded hysteresis}
\label{sec:orgadecea3}

We fix one label $X\in\{\mathrm{HC},\mathrm{DDC}\}$ and write
$\alpha_{\mathrm{mut}}(v)$ and $\alpha_{\cup}(v)$ for label $X$ of
Eq.~\eqref{eq:si-ice} on the mutual four-nearest graph
$G_{\mathrm{mut}}$ and on the union graph $G_{\cup}$, each from the
primitive six-rings of its own graph. Affiliation on $G_{\mathrm{mut}}$
supplies seeds, and the candidates are
$A=\{v:\alpha_{\mathrm{mut}}(v)\lor\alpha_{\cup}(v)\}$. An edge of
$G_{\cup}$ absent from $G_{\mathrm{mut}}$ can make a mutual-graph
six-ring non-primitive or change the bonds that $\beta$ reads, so
$\alpha_{\mathrm{mut}}$ does not imply $\alpha_{\cup}$. We accept a candidate
only when its connected component $C(v)$ of $G_{\cup}[A]$
contains a seed,
\begin{equation}
\mathrm{seeded}(v)
  \iff
  \exists\, u\in C(v)
  \text{ with }
  \alpha_{\mathrm{mut}}(u).
\label{eq:si-seeded}
\end{equation}
In Lean \texttt{seeded} is that existential over an abstract
component relation \texttt{inComp}.

\begin{theorem}[No mutual seed, nothing accepted]\label{thm:seeded-none}
If the mutual pass affiliates no molecule, then Eq.~\eqref{eq:si-seeded}
is false for every molecule.
\leanname{seeded_false_of_no_mutual}{Seeded}
\end{theorem}

\begin{lemma}[A seed in the component suffices]\label{thm:seeded-of}
If a mutual-affiliated molecule lies in $C(v)$, then $\mathrm{seeded}(v)$.
\leanname{seeded_of_seed}{Seeded}
\end{lemma}

Theorem~\ref{thm:seeded-none} is conditional. Whether the mutual pass
affiliates any molecule of a given frame is a property of that frame,
and we measure it in the main text.

\section{Ring completion}
\label{sec:orgfd9db0c}

A six-ring of the union graph with every vertex labelled but one
takes the label on its last vertex. The completed set is the least
set that contains the seeded molecules and is closed under that rule,
\begin{equation}
\mathrm{completed}(v)
  \iff
  \mathrm{seeded}(v)
  \ \lor\ 
  \exists\, R\ni v:\ \forall u\in R\setminus\{v\}:\ \mathrm{completed}(u),
\label{eq:si-completion}
\end{equation}
read as an inductive definition, separately for the HC and the DDC
label. A completed molecule need not lie on any affiliated cage; the
completed label is a separate output from cage membership. The Lean
statement is an inductive predicate with a seed constructor and a fill
constructor.

\begin{paperprop}[The completion loop returns the least fixed point]\label{prop:queue}
Algorithm~\ref{alg:seeded} starts with the seeded set $A$ and a
stack $F$ holding every six-ring of $G_\cup$. It pops a ring $R$; when
$R\setminus A=\{u\}$ it adds $u$ to $A$ and pushes every six-ring
through $u$ that is not on $F$. It terminates with the set of
Eq.~\eqref{eq:si-completion} for every pop order.
\end{paperprop}

\begin{proof}
Each molecule is added at most once, and each addition pushes at
most $d_I$ rings, so at most $R_6+n\,d_I$ rings are popped and the
loop terminates. We induct about the sequence of additions. Every added molecule is the only unlabelled vertex
of a ring whose other vertices were already in $A$, so the final set $A$ is contained in the
completed set. For the converse we show that the final $A$ is
closed under the rule. Let $R$ have every vertex but $u$ in the final
$A$. If every vertex of $R\setminus\{u\}$ was seeded, $R$ was on $F$
at the start; otherwise let $w$ be the vertex of $R\setminus\{u\}$
added last, and $R$ was pushed when $w$ was added or was already on
$F$ then. In both cases $R$ is popped after every vertex of
$R\setminus\{u\}$ is in $A$, so at that pop $R\setminus A\subseteq\{u\}$
and $u$ is added if it is not already in $A$. Hence $A$ contains the
seeds and is closed, and contains the least such set. Neither
direction uses the pop order.
\end{proof}

\begin{theorem}[Completion is sound]\label{thm:completion-sound}
A completed molecule is seeded or lies on a ring whose other
molecules are completed.
\leanname{completed_sound}{Completion}
\end{theorem}

\begin{theorem}[Completion is the least closed set]\label{thm:completion-least}
Every predicate that contains the seeded molecules and is closed
under the all-but-one rule contains every completed molecule.
\leanname{completed_least}{Completion}
\end{theorem}

Theorems~\ref{thm:completion-sound} and~\ref{thm:completion-least}
characterize the inductive closure. They do not verify the stack
program; Proposition~P5 does that on paper.

\begin{theorem}[Empty seed completes nothing]\label{thm:completion-empty}
If no molecule is seeded and every ring has two distinct vertices, no
molecule is completed.
\leanname{completed_empty}{Completion}
\end{theorem}

\begin{lemma}[The edge rule is not the completion]\label{thm:completion-edge}
On the triangle $\{0,1,2\}$ and the four-ring $\{0,1,3,4\}$ with the
triangle seeded, the rule that labels every vertex of a ring sharing
two vertices with a fully seeded ring labels molecule $3$, and
Eq.~\eqref{eq:si-completion} does not.
\leanname{edge_rule_not_completion}{Completion}
\end{lemma}

\section{Cage signatures and occupancy}
\label{sec:orgac8eab0}

A signature is a ring-size census. The enumerator grows a connected
set of primitive rings on the ring-adjacency graph and keeps a set
$F$ that is closed. A face set is closed when every edge of every face belongs to exactly two
faces. An empty face list is not closed.

Occupancy is two predicates. The radius assignment sends a guest to
the nearest periodic centroid within $r$. The inside assignment
unwraps the faces and the guest about that centroid and fans each face
from its centroid into triangles. It counts the guest as inside
when it lies on a triangle or a ray from it meets an odd number of
triangles. A grazing hit on a shared edge retries a perturbed ray.

A water-only primitive ring whose vertices all lie in an ion's first
shell is capped. A primitive ring of the water graph plus a virtual
ion bonded to every shell water, that contains the ion, is broken.

\paragraph{The Lean occupancy model.}
Lean represents the inside test by the parity oracle
$\mathtt{insideParity}\;\mathtt{onFace}\;\mathtt{crossings}:=\mathtt{onFace}\lor\mathtt{crossings}\bmod 2=1$.
It contains no guest coordinates, ray, face set or triangle, and no
relation between \texttt{onTriangle} and \texttt{insideParity}.
Theorems~\ref{thm:on-face-inside}--\ref{thm:even-outside} are
statements about that oracle.

\begin{theorem}[Empty faces are not closed]\label{thm:empty-not-closed}
An empty face list is not a closed polyhedron.
\leanname{empty_faces_not_closed}{Occupancy}
\end{theorem}

\begin{theorem}[Radius assignment lies within $r$]\label{thm:radius-within}
A radius-assigned cage of a guest is at distance at most $r$.
\leanname{radius_self}{Occupancy}
\end{theorem}

\begin{theorem}[Parity oracle, on a face]\label{thm:on-face-inside}
For
\[\mathtt{insideParity}\;\mathtt{onFace}\;\mathtt{crossings}:=\mathtt{onFace}\lor\mathtt{crossings}\bmod 2=1,\]
the value $\mathtt{insideParity}\;\mathtt{True}\;\mathtt{crossings}$ holds for every count.
\leanname{on_face_inside}{Occupancy}
\end{theorem}

\begin{theorem}[Parity oracle, odd count]\label{thm:odd-inside}
For every natural $n$,
\[\mathtt{insideParity}\;\mathtt{False}\;(2n+1)\]
holds.
\leanname{odd_crossings_inside}{Occupancy}
\end{theorem}

\begin{theorem}[Parity oracle, even count]\label{thm:even-outside}
For every natural $n$,
\[\neg\,\mathtt{insideParity}\;\mathtt{False}\;(2n).\]
\leanname{even_crossings_outside}{Occupancy}
\end{theorem}

\begin{theorem}[Empty ring is not capped]\label{thm:empty-not-capped}
The empty vertex list is not a capped ring.
\leanname{empty_ring_not_capped}{Occupancy}
\end{theorem}

\section{Inverted index and resident batch}
\label{sec:org48290e8}

Both cage criteria ask which six-rings contain a given atom
\citep{goswami2020dseams}. One pass
\begin{equation}
I(v)=\bigl\{R:v\in R\bigr\}
\label{eq:si-index}
\end{equation}
answers it: $R\in I(v)$ if and only if $R$ is stored and $v\in R$.

The engine probes the free device memory $M$ in bytes and forms the
integer byte budget $A=\lfloor sM\rfloor$ with the safety fraction
$s=0.8$. Its per-frame estimate $W>0$ is the sum of every per-frame
device array, each rounded up to 256 bytes. The arrays take the planned sizes of
16 six-rings per atom, 16 ring identifiers per row of $I$ and 8 basal
pairs per ring. A frame whose counts pass a planned size is rerun with
that array grown, and the frames still pending go to the host path
when the grown arrays cannot be allocated. The batch is
\begin{equation}
F=\min\bigl(F_{\mathrm{req}},\ \lfloor A/W\rfloor\bigr),
\label{eq:si-resident}
\end{equation}
where $F_{\mathrm{req}}$ is the number of frames requested. Rounding
each $F$-frame array up to 256 bytes adds less than rounding each of
the $F$ per-frame arrays separately, so the $F$-frame allocation is
at most $FW$. Integer division writes $A=\lfloor A/W\rfloor W+\rho$
with $0\le\rho<W$, so
\begin{align}
\lfloor A/W\rfloor\,W
  &= A-\rho
    \le A,
\label{eq:si-fits}\\
\bigl(\lfloor A/W\rfloor+1\bigr)W
  &= A-\rho+W
    > A.
\label{eq:si-max}
\end{align}
The second line is maximality under the per-frame estimate $W$ with no
fixed allocation; it applies when $F$ is not capped by
$F_{\mathrm{req}}$. Lean states both lines for natural numbers with
the budget written as a product $sM$ of naturals. Setting $s=1$ and
$M=A$ gives the integer-budget form above; the fractional $s=0.8$ and
its rounding to $A$ are not part of the Lean statement.

The cage-kernel grid is
\begin{equation}
N_{\mathrm{blk}}
  =\bigl\lceil nF\tau/B\bigr\rceil
  =(nF\tau+B-1)/B,
\label{eq:si-blocks}
\end{equation}
for $B>0$. Let $a=nF\tau$ and $k=a+B-1$. Then
$k=qB+\rho'$ with $0\le\rho'<B$, and
$N_{\mathrm{blk}}B=qB=k-\rho'\ge k-(B-1)=a$.

\begin{theorem}[Inverted index membership]\label{thm:index}
$R\in I(v)$ if and only if $R$ is a stored ring and $v\in R$
(Eq.~\eqref{eq:si-index}).
\leanname{inverted_index_mem}{Index}
\end{theorem}

\begin{theorem}[Natural-division quotient fits]\label{thm:resident-fits}
For natural numbers $s,M,W$, define $F=(sM)/W$ using natural-number
division. Then $F\,W\le sM$.
\leanname{resident_fits}{Resident}
\end{theorem}

\begin{theorem}[Its successor does not fit]\label{thm:resident-max}
For natural numbers $s,M,W$ with $W>0$ and $F=(sM)/W$,
$(F+1)\,W>sM$.
\leanname{resident_maximal}{Resident}
\end{theorem}

\begin{theorem}[Grid covers every source]\label{thm:nblocks}
$n F\tau\le \lceil n F\tau/B\rceil\cdot B$ for natural numbers with
$B>0$ (Eq.~\eqref{eq:si-blocks}).
\leanname{nblocks_covers}{Resident}
\end{theorem}

\section{CHILL windows}
\label{sec:org184afb6}

CHILL and CHILL+ classify the four-neighbour star of a molecule by
its staggered and eclipsed bonds
\citep{moore2010chill,nguyen2015identification}. Each $c_{ij}$ reads
$q_{lm}(j)$, so the label reads one hop beyond the star. They are not
Eq.~\eqref{eq:si-ice}. A bond is staggered when $c_{ij}\le-0.8$. Ice
Ic is $S=4$ and $E=0$; ice Ih is $S=3$ and $E=1$. Those assignments
cannot hold together: $S$ cannot be both $4$ and $3$.

\begin{proposition}[CHILL Ic and Ih are exclusive]\label{thm:chill-xor}
The CHILL ice-Ic and ice-Ih count windows cannot hold together.
\leanname{chill_ic_not_ih}{CHILL}
\end{proposition}

\section{Clustering}
\label{sec:org5328df8}

We model the Stoddard circular-list clustering used by the engine
\citep{goswami2020dseams} in Lean as reachability within
a step bound under an arbitrary successor map, and that relation is
reflexive. Following the successor map
zero times returns the start, so every molecule belongs to the
component started at itself.

\begin{lemma}[Cluster reflexivity]\label{thm:cluster-refl}
Every molecule is in the cluster started at itself.
\leanname{same_cluster_refl}{Cluster}
\end{lemma}

\section{Lean development and its model boundary}
\label{sec:lean}

The declarations are the library \texttt{DSeams} in the
\texttt{lean/} folder of the supplementary archive. The library has thirteen files in
\texttt{lean/DSeams/} with 55 theorem declarations, 2 of them private,
and builds on the Lean toolchain \texttt{v4.15.0}
with no other library. Theorems S1--S34 state 34 of its 53 public theorems. Two commands from that folder, \texttt{elan toolchain install leanprover/lean4:v4.15.0} and \texttt{lake build}, check every declaration. No declaration uses \texttt{sorry} or adds
an axiom. We run \texttt{\#print axioms} on all 53 public theorems, and it reports at
most \texttt{propext}, \texttt{Quot.sound} and \texttt{Classical.choice},
the standard axioms of Lean~4. \texttt{lean/README.md} lists each
theorem's supplement number, Lean name and file.
The seams-core repository also carries a separate Coq development in \texttt{proofs/}, present since v2.11.0, and its pixi environment pins Coq~8.12.2.
Its four files also compile with Coq~8.18.0, and it closes 7 of its 12 results and admits the other 5, among them the no-shortcut property that Theorem~\ref{thm:franzblau-noshort} proves in Lean.

The Lean development proves only the abstract declarations listed
above. In Lean, \texttt{Graph.adj} is an unconstrained Boolean relation, with
no symmetry or irreflexivity. Candidate and cage collections are
lists. The HC and DDC affiliation predicates are simplified as stated
in Section~\ref{sec:org36e36f2}. The cage structures store their
cardinalities and DDC conditions as fields, and occupancy is a parity
oracle. The enumerator \texttt{full} is an arbitrary function subject
to an extensional locality hypothesis, and the batch variables are
natural numbers. The development contains no refinement proof
connecting these objects to the C++ data structures or algorithms.
We prove Propositions~P1--P5 on paper in this supplement. We test the C++
implementation separately on the fixtures of the main text.

\section{Label accuracy}
\label{sec:s-accuracy}
\subsection{Cage membership and the receptive field}
\label{sec:cagestar}
Eq.~\mref{eq:ice} is not a star function.
We fix \(v\) and take two configurations that agree on \((G[B_1(v)],\{x_u:u\in B_1(v)\})\).
They have identical positions for \(v\) and its four neighbours, and identical induced edges among those five molecules.
Displacing molecules of the second shell alone leaves that data unchanged and removes the peripheral rings of the cage through \(v\), so \(\mathrm{Ic}(v)\) is true in the first configuration and false in the second.
A star function returns one value on both, and therefore equals Eq.~\mref{eq:ice} on neither.
We construct the pair from a frame of the cubic fixture and check it in the suite (Section~\ref{sec:verification}).

The bound is on the receptive field rather than on locality.
Section~\mref{sec:algorithms} shows that Eq.~\mref{eq:ice} at \(v\) is decided by the ring complex of a ball of bounded radius, which makes an incremental update exact.
A star function lacks that ball.
The published scores read different amounts of the same neighbourhood.
Eq.~\ref{eq:ql} at \(N_b=4\) is a star function.
CHILL and CHILL+ read one hop further, since \(c_{ij}\) of Eq.~\ref{eq:cij} contains \(q_{lm}(j)\).
The average \(\bar q_l\) takes the mean of \(q_{lm}\) over \(B_1\).
Polyhedral template matching (PTM) matches a template of the first and second shells for cubic and hexagonal diamond \citep{larsen2016robust}, which is how it separates the two polymorphs.
DeepIce \citep{fulford2019deepice} is a network on the Cartesian coordinates of \(n\) nearest neighbours, trained to ice versus liquid, and does not distinguish ice Ic from ice Ih.
IceCoder \citep{maity2025icecoder} is a variational autoencoder on SOAP descriptors of the local oxygen neighbourhood; its published check on nucleation trajectories is agreement with CHILL+.
Neither computes \(\alpha_{\mathrm{HC}}\) or \(\alpha_{\mathrm{DDC}}\).
For the star and template scores the difference from Eq.~\mref{eq:ice} is a measurement rather than a theorem.
We measure it in Table~\ref{tab:sota}.

Two frames show the difference in both directions.
On defect silver iodide the first two adsorbed water layers are HC-dominated \citep{prerna2019ice}.
On the complementary AgI frame the first adsorbed layer sits in complete cages while CHILL labels it liquid \citep{goswami2020dseams}.
Stacking-disordered ice I nuclei mix cubic and hexagonal sequences \citep{lupi2017stacking}.
The cage search reports both hexagonal and double-diamond cages that share rings \citep{goswami2020dseams}.
A per-molecule label reports one class per molecule.

Table~\ref{tab:sota} measures lattice-label retention: the fraction of the sites of an ideal lattice that keep their construction label after a Gaussian displacement of standard deviation \(\sigma\) in each Cartesian component, an RMS displacement of \(\sqrt3\,\sigma\).
The lattices are cubic diamond (\(4\times4\times4\) conventional cells, 512 sites) and lonsdaleite (\(5\times3\times3\) orthohexagonal cells, 360 sites) at an O--O distance of \(2.75\,\mathring{\mathrm{A}}\).
We run five realizations at each \(\sigma\), with generator seeds 5 to 9 of \texttt{scripts/sota\_compare.py} (\texttt{-{}-seed-offset 5}).
The earlier campaign records hold seeds 0 to 4, and no record before this sweep contains seeds 5 to 9.
We count a site as correct when its label is the construction polymorph, and a site labelled both Ic and Ih as wrong.
One input without ice follows, and its column gives the fraction of molecules labelled Ic or Ih.
The liquid is a 512-molecule mW configuration equilibrated at \(215\,\mathrm{K}\), the liquid stage of the seeding workflow.
CHILL+ is CHILL+\((r_c)\).
For PTM we run OVITO's \texttt{PolyhedralTemplateMatchingModifier} with \texttt{rmsd\_cutoff=0.1}, the value its manual gives as a reasonable cutoff (\url{https://www.ovito.org/docs/current/reference/pipelines/modifiers/polyhedral_template_matching.html}), and the cubic- and hexagonal-diamond templates enabled next to its default templates, on the periodic orthorhombic cell.
A molecule matched to cubic diamond is Ic, one matched to hexagonal diamond is Ih, and any other result, no match included, is unlabelled.
The Lechner--Dellago row is a nearest-centroid rule in the \((\bar q_4,\bar q_6)\) plane on four nearest neighbours.
The cubic and hexagonal centroids are the means of ideal Ic and ideal Ih.
The water centroid is the mean of the coexistence liquid at \(274.6\,\mathrm{K}\).
The \((\bar q_3,\bar q_{12})\) row is the same rule in the plane of the averaged \(q_3\) and \(q_{12}\) that \citet{zeron2024hydrate} use for hydrates, computed with freud on the same four nearest neighbours, and its water centroid is that same liquid.
A molecule is ice under the \(q_6\) row when \(q_6>\tau\), and that row cannot separate Ic from Ih.
The ice error at a candidate \(\tau\) is the fraction of the coexistence-crystal molecule-frames with \(q_6\le\tau\).
The liquid error is the fraction of the coexistence-liquid molecule-frames with \(q_6>\tau\).
The candidates are the sorted unique values in the two samples, which are the values at which either error changes.
\(\tau\) is the candidate with the smallest absolute difference of the two errors, then the smaller of the two errors, then the smaller candidate.
Both samples are the five production frames at one state point, \(20160\) crystal molecule-frames and \(20480\) liquid molecule-frames.
The ice error is \(4044/20160 = 0.2005952380952381\) and the liquid error is \(4108/20480 = 0.2005859375\).
That threshold is \(\tau=0.5579227209091187\).
The mean \(q_6\) of the hexagonal crystal is \(0.5803724527359009\) and the mean \(q_6\) of the liquid is \(0.4868989884853363\).
The same rule at the cubic melting point gives \(\tau_{\mathrm{Ic}}=0.5592451095581055\), and both errors are \(4019/20480 = 0.196240234375\).
The means there are \(0.5819523930549622\) and \(0.4874526560306549\).
The cross-check threshold is the equal-error point of the two normal distributions \citet{lechner2008accurate} publish for unaveraged \(q_6\) in Tables 1 and 2.
Those tables are the Lennard-Jones FCC crystal, mean \(0.507298\) and standard deviation \(0.043301\), and the undercooled liquid, mean \(0.360012\) and standard deviation \(0.066518\), at \(k_{\mathrm{B}}T=0.92\varepsilon\) and \(P=5.68\,\varepsilon\sigma^{-3}\), with every neighbour inside \(1.4\sigma\).
That paper states no \(q_6\) threshold.
With ice on the high side the equal-error point of the two normals is \(\tau_{\mathrm{LD}}=(\mu_{\mathrm{liq}}\sigma_{\mathrm{FCC}}+\mu_{\mathrm{FCC}}\sigma_{\mathrm{liq}})/(\sigma_{\mathrm{FCC}}+\sigma_{\mathrm{liq}})=0.4492239774173867909924512152\).
The neighbour shell in that paper is not the four nearest neighbours used here, and the substance is a Lennard-Jones fluid.
Table~\ref{tab:sota} and Table~\ref{tab:md} give the \(q_6\) cells at \(\tau\) and, in the following row, at \(\tau_{\mathrm{LD}}\).
The primary state is the ice--liquid coexistence point of hexagonal ice.
\citet{molinero2009mw} write that hexagonal ice (\(T_{\mathrm{m}}=274.6\pm 1\,\mathrm{K}\)) is more stable than cubic ice (\(T_{\mathrm{m}}=271.5\pm 1\,\mathrm{K}\)).
Table 1 of that paper lists the mW liquid density at \(298\,\mathrm{K}\) as \(0.997\,\mathrm{g\,cm^{-3}}\), the ice density at melting as \(0.978\,\mathrm{g\,cm^{-3}}\) and the liquid density at melting as \(1.001\,\mathrm{g\,cm^{-3}}\).
The crystal and the liquid of \(\tau\) are both at \(274.6\,\mathrm{K}\) and \(1\,\mathrm{atm}\).
The crystal is ice Ih, 4032 molecules, velocity seed 97461.
The liquid is 4096 molecules, packing seed 1746 and velocity seed 77461, melted for \(100\,\mathrm{ps}\) at \(350\,\mathrm{K}\) and then held at the target.
The cross-check is the cubic coexistence point, \(271.5\,\mathrm{K}\) and \(1\,\mathrm{atm}\): ice Ic, 4096 molecules, velocity seed 97151, and a liquid with packing seed 1715 and velocity seed 77151.
The scored temperatures are \(200\), \(215\), \(250\) and \(300\,\mathrm{K}\).
The thermostat and the isotropic barostat are Tdamp \(1000\,\mathrm{fs}\) and Pdamp \(10000\,\mathrm{fs}\).
Each crystal is equilibrated for \(300\,\mathrm{ps}\) of NPT.
Each liquid is equilibrated for \(300\,\mathrm{ps}\) of NPT after the melt.
Every trajectory is then settled for \(50\,\mathrm{ps}\) and saved as five frames over \(100\,\mathrm{ps}\).
The thirty Ih equilibration rows average \(273.769\,\mathrm{K}\), \(4.314\,\mathrm{atm}\) and \(0.978608\,\mathrm{g\,cm^{-3}}\).
Its ten production rows average \(275.004\,\mathrm{K}\), \(-17.828\,\mathrm{atm}\) and \(0.978418\,\mathrm{g\,cm^{-3}}\), with a centre-of-mass mean-squared displacement of \(0.336\,\mathring{\mathrm{A}}^2\) at \(100\,\mathrm{ps}\).
The thirty-one liquid equilibration rows from step 21000 at \(274.6\,\mathrm{K}\) average \(274.594\,\mathrm{K}\) and \(-13.109\,\mathrm{atm}\).
Its ten production rows average \(274.610\,\mathrm{K}\), \(93.596\,\mathrm{atm}\) and \(1.001839\,\mathrm{g\,cm^{-3}}\).
The Ic production rows average \(270.312\,\mathrm{K}\) and \(0.978800\,\mathrm{g\,cm^{-3}}\), and the \(271.5\,\mathrm{K}\) liquid production rows average \(271.060\,\mathrm{K}\) and \(1.002387\,\mathrm{g\,cm^{-3}}\).
\texttt{sota\_compare.py -{}-write-calibration} writes the centroids and the thresholds to \texttt{calibration.txt}.
Table~\ref{tab:sota} and Table~\ref{tab:md} were scored from that file at sha256 \texttt{09a15650}.
The deposited copy keeps those decision lines, states that the 2008 paper gives no \(q_6\) threshold, and appends the measured thermostat rows.
Its sha256 is \texttt{dbac499b}.
At \(\tau\) the ideal lattices, with no displacement, are labelled ice in full: \(1.000\) of the 512 Ic sites and \(1.000\) of the 360 Ih sites.
The jittered columns keep \(0.720\) and \(0.726\) at \(0.2\,\mathring{\mathrm{A}}\) and \(0.237\) and \(0.238\) at \(0.4\,\mathring{\mathrm{A}}\).
The held-out crystals keep \(0.925\) and \(0.924\) at \(200\,\mathrm{K}\) and \(0.848\) and \(0.848\) at \(250\,\mathrm{K}\).
The \(215\,\mathrm{K}\) liquid contributes 135 of 512 molecules.
At \(\tau_{\mathrm{LD}}\) the jittered columns stay above \(0.72\) and the held-out crystals stay above \(0.997\), and that liquid contributes 388 of 512.
The primary row is \(\tau\).
\(\tau_{\mathrm{Ic}}\) is the cubic-coexistence cross-check and \(\tau_{\mathrm{LD}}\) is the published-normal cross-check.
They part on the liquid: \(\tau\) marks \(0.264\) of it and \(\tau_{\mathrm{LD}}\) marks \(0.758\).

At \(\sigma=0.2\,\mathring{\mathrm{A}}\) the cutoff, mutual and seeded cage labels keep every cubic site, and all but the cutoff graph (\(0.999\)) keep every hexagonal one. PTM keeps \(0.991\) of the cubic sites and \(0.991\) of the hexagonal sites. CHILL+ keeps \(0.971\) and \(0.950\).
The Lechner--Dellago rule keeps \(0.563\) of the cubic sites at \(0.2\,\mathring{\mathrm{A}}\) and labels \(0.435\) of them hexagonal.
The \((\bar q_3,\bar q_{12})\) rule keeps \(0.240\) of the cubic and \(0.864\) of the hexagonal sites at \(0.2\,\mathring{\mathrm{A}}\), and \(0.000\) and \(0.098\) at \(0.4\,\mathring{\mathrm{A}}\).
Moving the water centroid onto the coexistence liquid at \(274.6\,\mathrm{K}\) moves the hexagonal cells.
The \((\bar q_4,\bar q_6)\) cells go from \(0.968\), \(0.344\), \(0.096\), \(0.031\) to \(0.970\), \(0.349\), \(0.097\), \(0.032\).
The \((\bar q_3,\bar q_{12})\) cells go from \(0.880\), \(0.393\), \(0.224\), \(0.105\) to \(0.864\), \(0.383\), \(0.216\), \(0.098\).
The cubic cells of both rows are unchanged at three decimals.
At \(0.3\,\mathring{\mathrm{A}}\) the cutoff graph loses edges as neighbours cross \(3.5\,\mathring{\mathrm{A}}\) and keeps \(0.936\) and \(0.899\); the mutual graph keeps \(0.990\) and \(0.940\), the seeded rule \(0.999\) and \(0.994\), PTM \(0.089\) and \(0.099\) and CHILL+ \(0.564\) and \(0.567\).
At \(0.4\,\mathring{\mathrm{A}}\) the seeded rule keeps \(0.698\) of Ic (realizations \(0.594\) to \(0.824\)) and \(0.589\) of Ih (\(0.511\) to \(0.678\)).
At the same displacement the cutoff graph keeps \(0.346\) and \(0.233\), PTM \(0.000\) and \(0.000\) and CHILL+ \(0.100\) and \(0.099\).
Ring completion raises the seeded values to \(0.942\) (\(0.883\) to \(0.994\)) and \(0.728\) (\(0.597\) to \(0.844\)) at \(0.4\,\mathring{\mathrm{A}}\).
On the liquid the cutoff graph labels \(0.047\) of the molecules Ic or Ih, the mutual and seeded rules \(0.023\), CHILL+ \(0.008\), the Lechner--Dellago rule \(0.002\) (1 of 512), the \((\bar q_3,\bar q_{12})\) rule \(0.049\) (25 of 512) and the \(q_6\) rule at \(\tau\) labels \(0.264\) (135 of 512).
The \(q_6\) cross-check at \(\tau_{\mathrm{LD}}\) labels \(0.758\) (388 of 512).
PTM labels none of them.
DeepIce and IceCoder are not rows of that table, because their published scopes are ice versus liquid and agreement with CHILL+, respectively.
At \(\sigma=0.3\,\mathring{\mathrm{A}}\) on the cubic lattice, generator seed 5, the completed rule, CHILL+ and PTM are one session on the core of Table~\mref{tab:compare}: five warm-ups and 30 repetitions, \texttt{scripts/ptm\_rows.py}.
The completed rule and CHILL+ are \texttt{tests/bench\_paper.cpp}, a release build.
CHILL+ calls \texttt{neighListO} at \(3.5\,\mathring{\mathrm{A}}\), \texttt{getCorrelPlus} and \texttt{getIceTypePlusNoPrint}, and that neighbour list is inside the timer.
The completed rule is mode \texttt{seeded-adj}.
Its timer is \texttt{std::chrono::steady\_clock} around \texttt{seededLabels}, which builds the neighbour lists inside the timer.
The \(5.0\,\mathring{\mathrm{A}}\) argument of \texttt{kNearestNeighbourList} only sizes its cell list.
The search falls back to brute force for any molecule whose fourth neighbour lies beyond it, so the labels do not depend on it (\texttt{tests/test\_neighbours.cpp}).
The medians are \(7.104\,\mathrm{ms}\) for the completed rule, \(0.350\,\mathrm{ms}\) for CHILL+ and \(16.312\,\mathrm{ms}\) for PTM.
PTM is \texttt{PolyhedralTemplateMatchingModifier}, called in the Python process of \texttt{scripts/ptm\_rows.py}, which starts \texttt{bench\_paper} as a child process pinned with \texttt{taskset}.
The unrounded medians give \(7.103615/0.349822 = 20.306\), so the completed rule takes that many times as long as CHILL+.
The earlier supplement printed \(15.6\,\mathrm{ms}\) for the completed rule.
That line is an earlier, undeposited run of the command that wrote \texttt{results/reference/}\allowbreak\texttt{sota-heldout-2026-10-31/}\allowbreak\texttt{sota\_compare-seeds5-9.txt}.
The deposited copy of that file prints \(15.8\,\mathrm{ms}\) for \texttt{dseams-topo-seeded-adj} on ice Ic at this displacement and seed.
Both are the one-decimal Python timer \texttt{timed()} around that call through PydSEAMSlib, not this C++ session.

The per-molecule column is the comparison the other methods admit.
The cage is the object of the analysis.
The analysis asks which molecules close an HC or a DDC, which six-rings two cages share and which prisms a confined film or a nanotube carries \citep{goswami2020nanoribbon}.
A growing nucleus is described by its hexagonal and double-diamond cages \citep{goswami2020dseams,goswami2020seeding}.
A stacking-disordered ice I nucleus mixes cubic and hexagonal sequences \citep{lupi2017stacking}.
Hydrate precursors are counted as cages \citep{jacobson2010amorphous,zeron2024hydrate}, and the deformation of a cage is a shape match against its template \citep{goswami2021tum}.
A per-molecule score returns local descriptors rather than the extended topology of these objects.
Table~\ref{tab:sota} reports the per-molecule column on ice I and the liquid.

\begin{table}[htbp]
\caption{\label{tab:sota}Lattice-label retention on ice Ic (512 sites) and Ih (360 sites) under a Gaussian displacement of \(\sigma\) per component (\(0.2\), \(0.3\), \(0.35\), \(0.4\,\mathring{\mathrm{A}}\)), means over realizations 5 to 9.
The last column gives the fraction labelled Ic or Ih on the 512-molecule mW liquid at \(215\,\mathrm{K}\).
Labels other than the PTM row are from PydSEAMSlib built on seams-core \coretag, from freud and from OVITO, at the versions that record \texttt{sota-heldout-2026-10-31} prints.
The PTM row is the mean over seeds 5 to 9 at \texttt{rmsd\_cutoff=0.1}, from \texttt{scripts/ptm\_cutoff\_sweep.py}.
Its liquid cell is the false-crystal fraction from \texttt{scripts/sota\_compare.py -{}-liquid-data} on \texttt{data/mW\_liquid\_T215\_512.data} of the reproducibility package, and that fraction is \(0.000\).
Section~\ref{sec:cagestar} defines the rules.
The Lechner--Dellago and \((\bar q_3,\bar q_{12})\) centroids use the ideal lattices and the coexistence liquid at \(274.6\,\mathrm{K}\).
The \(q_6\) threshold uses ice Ih and the liquid at \(274.6\,\mathrm{K}\) and \(1\,\mathrm{atm}\).
The \(q_6\) rows are ice versus not and give the fraction labelled ice.
The first of them uses \(\tau\) and the second uses \(\tau_{\mathrm{LD}}\).
With every seeded site already correct at \(0.2\,\mathring{\mathrm{A}}\), completion has nothing to add there.
DeepIce (\url{https://github.com/mfulford/DeepIce}) returns ice versus liquid and IceCoder (\url{https://github.com/TeamSuman/IceCoder}) is scored by agreement with CHILL+, so neither has an Ic/Ih row.}
\centering
\scriptsize
\resizebox{\textwidth}{!}{%
\begin{tabular}{lccccccccc}
method & Ic .2 & Ic .3 & Ic .35 & Ic .4 & Ih .2 & Ih .3 & Ih .35 & Ih .4 & liquid\\
\hline
cages, cutoff (Eq.~\mref{eq:ice}) & 1.000 & 0.936 & 0.684 & 0.346 & 0.999 & 0.899 & 0.607 & 0.233 & 0.047\\
cages, mutual 4-NN & 1.000 & 0.990 & 0.772 & 0.419 & 1.000 & 0.940 & 0.698 & 0.311 & 0.023\\
cages, seeded (Eq.~\ref{eq:seeded}) & 1.000 & 0.999 & 0.951 & 0.698 & 1.000 & 0.994 & 0.917 & 0.589 & 0.023\\
cages, seeded + ring completion (Eq.~\ref{eq:completion}) & 1.000 & 1.000 & 1.000 & 0.942 & 1.000 & 1.000 & 0.997 & 0.728 & 0.023\\
CHILL+\((r_c)\) & 0.971 & 0.564 & 0.244 & 0.100 & 0.950 & 0.567 & 0.276 & 0.099 & 0.008\\
PTM (OVITO) & 0.991 & 0.089 & 0.003 & 0.000 & 0.991 & 0.099 & 0.003 & 0.000 & 0.000\\
\((\bar q_4,\bar q_6)\) & 0.563 & 0.020 & 0.003 & 0.000 & 0.970 & 0.349 & 0.097 & 0.032 & 0.002\\
\((\bar q_3,\bar q_{12})\) & 0.240 & 0.009 & 0.002 & 0.000 & 0.864 & 0.383 & 0.216 & 0.098 & 0.049\\
\(q_6\) at \(\tau\) & 0.720 & 0.393 & 0.303 & 0.237 & 0.726 & 0.392 & 0.278 & 0.238 & 0.264\\
\(q_6\) at \(\tau_{\mathrm{LD}}\) & 0.983 & 0.867 & 0.788 & 0.745 & 0.986 & 0.859 & 0.790 & 0.727 & 0.758\\
\end{tabular}
}
\end{table}

On the 512-molecule liquid, deposited with the reproducibility package, the seeded rule labels \(12\) molecules and the completed rule \(12\).
The cutoff graph labels \(24\), CHILL+\((r_c)\) \(4\) as Ic or Ih, template matching \(0\), the Lechner--Dellago rule \(1\) and the \(q_6\) rule at \(\tau\) labels \(135\).
The cross-check at \(\tau_{\mathrm{LD}}\) labels \(388\).
The table uses the frozen file \texttt{calibration.txt}.
A liquid deep in the supercooled regime contains transient cages \citep{moore2011structural}.
A cage method reports those cages.
Nucleation studies therefore threshold on the largest connected cluster rather than on the count.

The jitter sweep is a constructed input.
In Table~\ref{tab:md} we score the same rules on physical trajectories that no rule was fitted on (record \texttt{md-heldout-2026-10-31}).
Ice Ic (4096 molecules) and ice Ih (4032) start from the lattices of Table~\ref{tab:sota} and run at \(200\) and \(250\,\mathrm{K}\), below the \(274.6\,\mathrm{K}\) melting point of mW \citep{molinero2009mw}.
The liquid (4096 molecules) is a random packing melted for \(100\,\mathrm{ps}\) at \(350\,\mathrm{K}\) and equilibrated for \(300\,\mathrm{ps}\) of NPT at \(215\), \(250\) or \(300\,\mathrm{K}\), the liquid stage of the seeding campaign.
Each of the seven conditions has three independent runs with their own velocities and, for the liquid, their own packing.
We run LAMMPS with a \(10\,\mathrm{fs}\) step.
Every run settles for \(50\,\mathrm{ps}\) of NPT at \(1\,\mathrm{atm}\) and is sampled five times, \(20\,\mathrm{ps}\) apart.
The Nose--Hoover damping is 100 timesteps in temperature and 1000 timesteps in pressure, the recommendation of the LAMMPS \texttt{fix nh} documentation (\url{https://docs.lammps.org/fix_nh.html}).
At a \(10\,\mathrm{fs}\) step that is \(1000\,\mathrm{fs}\) and \(10000\,\mathrm{fs}\), the damping of the deposited liquid input.
Record \texttt{md-heldout-2026-10-31} reports a mean-squared displacement, centre of mass removed, of \(0.21\) to \(0.29\,\mathring{\mathrm{A}}^2\) on every crystal run and \(63\) to \(401\,\mathring{\mathrm{A}}^2\) on every liquid run.
Those displacements were not remeasured on the regenerated dumps.
The Lechner--Dellago centroids and both \(q_6\) thresholds are those of Table~\ref{tab:sota}, read from the deposited \texttt{calibration.txt}.
No held-out frame enters the fit.
The \(215\,\mathrm{K}\) liquid of Table~\ref{tab:sota} is not the calibration trajectory.
The cage rules, CHILL+ and PTM have no data-fitted parameter.

On the regenerated crystal frames the seeded and completed rules label every molecule with its construction polymorph.
The cutoff and mutual graphs do the same except for one Ih run at \(250\,\mathrm{K}\), whose five-frame mean is \(0.99995\).
The CHILL+ 15-frame means are at least \(0.998\).
PTM keeps at least \(4093\) of \(4096\) Ic molecules and \(4031\) of \(4032\) Ih molecules in every \(200\,\mathrm{K}\) frame.
At \(250\,\mathrm{K}\) it misses up to \(65\) Ic molecules of a frame and \(73\) Ih molecules of a frame.
Those counts are \texttt{ok\_min} in the output of \texttt{scripts/md\_heldout\_score.py}.
The Lechner--Dellago rule keeps \(0.922\) of Ic at \(200\,\mathrm{K}\) and \(0.740\) at \(250\,\mathrm{K}\), and \(0.999\) and \(0.990\) of Ih.
The \(q_6\) rule at \(\tau\) calls \(0.848\) to \(0.925\) of the crystal ice.
The cross-check at \(\tau_{\mathrm{LD}}\) calls \(0.997\) to \(1.000\).
On the liquid frames the cutoff graph labels \(0.011\) to \(0.029\) of the molecules Ic or Ih.
The mutual, seeded and completed rules label \(0.010\) to \(0.012\) at \(215\,\mathrm{K}\) and at most \(0.001\) at \(250\) and \(300\,\mathrm{K}\).
CHILL+ and the Lechner--Dellago rule label at most \(0.005\).
The \(q_6\) rule at \(\tau\) calls \(0.198\) to \(0.273\) of the liquid ice, and the cross-check calls \(0.683\) to \(0.752\).
PTM labels no molecule on those frames.
On these crystals the cage-rule and CHILL+ means stay at least \(0.998\).
The Lechner--Dellago rule loses \(0.260\) of the Ic molecules at \(250\,\mathrm{K}\).
The primary \(q_6\) rule keeps at least \(0.848\) of either held-out crystal.
Table~\ref{tab:sota} is where \(\tau\) and \(\tau_{\mathrm{LD}}\) part on the liquid.
The \(21\) dumps, \(105\) frames, are not in this source tree.
The supplementary data hold them, with their SHA-256, so the labels can be recomputed without LAMMPS.
The runs use the seeds of record \texttt{md-heldout-2026-10-31} and the deposited templates.
The packaged \texttt{scripts/md\_heldout.py} writes a crystal start whose first line is \texttt{mW ic lattice, 4096 molecules} or \texttt{mW ih lattice, 4032 molecules}, and that file matches the record's SHA-256.
The deposited crystal starts use \texttt{mW lattice, 4096 molecules} and \texttt{mW lattice, 4032 molecules}.
The coordinate lines are the same six-decimal atom lines, so the coordinates agree and the bytes do not.
The liquid start in that script is the equilibrated data file written by \texttt{in.liquid.lammps}.
The deposited liquid \texttt{eq.data} files are the record's equilibrated starts and match its SHA-256; the production dumps are a new run from them, and the record's dumps are not deposited.
\texttt{scripts/md\_heldout\_score.py} reads the \texttt{ITEM: ATOMS} header of these dumps, which are \texttt{id type x y z}, and scores them from \texttt{calibration.txt}.

\begin{table}[htbp]
\caption{\label{tab:md}Held-out mW trajectories regenerated with the seeds of record \texttt{md-heldout-2026-10-31} and the deposited templates.
Every row is scored on those dumps.
On ice Ic (4096 molecules) and Ih (4032) at \(200\) and \(250\,\mathrm{K}\) the entries give the fraction of molecules labelled with the construction polymorph.
On the liquid (4096) at \(215\), \(250\) and \(300\,\mathrm{K}\) they give the fraction labelled Ic or Ih.
Each entry is the mean over 15 frames, five from each of three independent runs.
Every run mean lies within \(0.012\) of it.
The \((\bar q_4,\bar q_6)\) and \(q_6\) rules are those fitted for Table~\ref{tab:sota}.
The \(q_6\) rows are ice versus not and give the fraction labelled ice, the first at \(\tau\) and the second at \(\tau_{\mathrm{LD}}\).
The PTM row uses \texttt{rmsd\_cutoff=0.1}.}
\centering
\scriptsize
\resizebox{\textwidth}{!}{%
\begin{tabular}{lccccccc}
method & Ic 200 & Ic 250 & Ih 200 & Ih 250 & liq.\ 215 & liq.\ 250 & liq.\ 300\\
\hline
cages, cutoff (Eq.~\mref{eq:ice}) & 1.000 & 1.000 & 1.000 & 1.000 & 0.029 & 0.016 & 0.011\\
cages, mutual 4-NN & 1.000 & 1.000 & 1.000 & 1.000 & 0.010 & 0.001 & 0.001\\
cages, seeded (Eq.~\ref{eq:seeded}) & 1.000 & 1.000 & 1.000 & 1.000 & 0.012 & 0.001 & 0.001\\
cages, seeded + ring completion (Eq.~\ref{eq:completion}) & 1.000 & 1.000 & 1.000 & 1.000 & 0.012 & 0.001 & 0.001\\
CHILL+\((r_c)\) & 1.000 & 0.999 & 1.000 & 0.998 & 0.005 & 0.001 & 0.001\\
PTM (OVITO) & 1.000 & 0.990 & 1.000 & 0.989 & 0.000 & 0.000 & 0.000\\
\((\bar q_4,\bar q_6)\) & 0.922 & 0.740 & 0.999 & 0.990 & 0.004 & 0.001 & 0.001\\
\(q_6\) at \(\tau\) & 0.925 & 0.848 & 0.924 & 0.848 & 0.273 & 0.213 & 0.198\\
\(q_6\) at \(\tau_{\mathrm{LD}}\) & 1.000 & 0.997 & 1.000 & 0.997 & 0.752 & 0.697 & 0.683\\
\end{tabular}
}
\end{table}

Table~\ref{tab:sota} is a positive test on ice I. Table~\ref{tab:atlas} is the negative test.
We take every lattice of the GenIce library \citep{matsumoto2018genice} that fits an orthorhombic cell, generate it ideal and under Gaussian displacements of \(0.16\) and \(0.31\,\mathring{\mathrm{A}}\), and classify it by the same three methods.
The seeded cage assignment with ring completion returns ice I on ice Ih, Ic and XI at every displacement.
It returns nothing on the ten dense polymorphs, the four empty hydrate frameworks and the six porous ices of the table.
The cutoff graph has no seed and calls ice VI \(98\%\) ice at rest.
Its eight-coordinated oxygens close six-rings that satisfy the cage criteria of Eq.~\mref{eq:ddc1} and Eq.~\mref{eq:ddc2} as graphs.
CHILL+\((r_c)\) keeps clathrate as a label and returns it on the hydrate frameworks, as designed.
At rest it also returns clathrate on ice V, VI and VII and on the porous XVII, FAU and SOD.
At \(0.31\,\mathring{\mathrm{A}}\) it returns clathrate on \(13\%\) to \(43\%\) of every dense polymorph and \(11\%\) to \(76\%\) of every porous one, where the seeded assignment stays at zero.
The continuous random networks are the residual.
A four-connected amorphous network contains hexagonal and double-diamond cages as subgraphs, with \(21\) HC-affiliated and \(28\) DDC-affiliated six-rings on the cutoff graph of the thousand-molecule CRN1.
At rest the cage assignments report the molecules on them, between \(4.8\%\) and \(11.3\%\), as ice-like motifs.
CHILL+\((r_c)\) reports \(3\) to \(5\%\) ice and \(14\) to \(18\%\) interfacial on the same networks at rest.

\begin{table}[htbp]
\caption{\label{tab:atlas}Fraction of molecules labelled ice I by the seeded cage assignment with ring completion (seeded), by the cutoff cage assignment (cutoff) and by CHILL+ (ice), and the fraction CHILL+ labels clathrate.
The inputs are the GenIce library at rest and under Gaussian displacements of \(\sigma = 0.16\) and \(0.31\,\mathring{\mathrm{A}}\).
Ice I structures are the positive set; everything below the rule is a negative.
XVI and sII share one oxygen network and one row.
Ice II, XIII, sH and EMT have non-orthorhombic cells and are not in the table; the GenIce generator of stacking-disordered Isd exited with an error in the run and has no row.}
\centering
\scriptsize
\resizebox{\textwidth}{!}{%
\begin{tabular}{llrcccccccccc}
structure & class & \(N\) & seeded 0 & seeded .16 & seeded .31 & cutoff 0 & cutoff .16 & cutoff .31 & CHILL+ ice 0 & ice .31 & clathrate 0 & clathrate .31\\
\hline
Ih & ice I & 288 & 1.000 & 1.000 & 1.000 & 1.000 & 1.000 & 0.809 & 1.000 & 0.493 & 0.000 & 0.000\\
Ic & ice I & 216 & 1.000 & 1.000 & 1.000 & 1.000 & 1.000 & 0.954 & 1.000 & 0.514 & 0.000 & 0.000\\
XI & ice I & 288 & 1.000 & 1.000 & 1.000 & 1.000 & 1.000 & 0.854 & 1.000 & 0.528 & 0.000 & 0.000\\
\hline
0 & dense & 216 & 0.000 & 0.000 & 0.000 & 0.000 & 0.000 & 0.000 & 0.000 & 0.000 & 0.000 & 0.134\\
III & dense & 324 & 0.000 & 0.000 & 0.000 & 0.000 & 0.000 & 0.000 & 0.000 & 0.000 & 0.000 & 0.151\\
IV & dense & 384 & 0.000 & 0.000 & 0.000 & 0.276 & 0.510 & 0.167 & 0.000 & 0.000 & 0.000 & 0.284\\
V & dense & 504 & 0.000 & 0.000 & 0.000 & 0.000 & 0.024 & 0.000 & 0.000 & 0.000 & 0.286 & 0.232\\
VI & dense & 270 & 0.000 & 0.000 & 0.000 & 0.985 & 0.000 & 0.000 & 0.000 & 0.000 & 1.000 & 0.426\\
VII & dense & 432 & 0.000 & 0.000 & 0.000 & 0.000 & 0.000 & 0.000 & 0.000 & 0.000 & 0.236 & 0.278\\
VIII & dense & 512 & 0.000 & 0.000 & 0.000 & 0.000 & 0.000 & 0.000 & 0.000 & 0.000 & 0.002 & 0.299\\
IX & dense & 324 & 0.000 & 0.000 & 0.000 & 0.000 & 0.000 & 0.000 & 0.000 & 0.000 & 0.000 & 0.145\\
XII & dense & 648 & 0.000 & 0.000 & 0.000 & 0.000 & 0.000 & 0.019 & 0.000 & 0.000 & 0.000 & 0.293\\
XIV & dense & 192 & 0.000 & 0.000 & 0.000 & 0.000 & 0.000 & 0.000 & 0.000 & 0.000 & 0.000 & 0.245\\
XVI (sII) & hydrate & 136 & 0.000 & 0.000 & 0.000 & 0.000 & 0.000 & 0.000 & 0.000 & 0.000 & 1.000 & 0.772\\
sI & hydrate & 368 & 0.000 & 0.000 & 0.000 & 0.000 & 0.000 & 0.000 & 0.000 & 0.000 & 1.000 & 0.701\\
sT & hydrate & 920 & 0.000 & 0.000 & 0.000 & 0.000 & 0.000 & 0.000 & 0.000 & 0.000 & 0.896 & 0.266\\
sIV & hydrate & 320 & 0.000 & 0.000 & 0.000 & 0.000 & 0.000 & 0.000 & 0.000 & 0.000 & 1.000 & 0.731\\
XVII & porous & 216 & 0.000 & 0.000 & 0.000 & 0.000 & 0.000 & 0.000 & 0.000 & 0.000 & 0.167 & 0.403\\
A & porous & 324 & 0.000 & 0.000 & 0.000 & 0.000 & 0.000 & 0.000 & 0.000 & 0.000 & 0.000 & 0.120\\
B & porous & 384 & 0.000 & 0.000 & 0.000 & 0.000 & 0.000 & 0.000 & 0.000 & 0.000 & 0.000 & 0.112\\
FAU & porous & 192 & 0.000 & 0.000 & 0.000 & 0.000 & 0.000 & 0.000 & 0.000 & 0.000 & 0.177 & 0.365\\
RHO & porous & 384 & 0.000 & 0.000 & 0.000 & 0.000 & 0.000 & 0.000 & 0.000 & 0.000 & 0.000 & 0.419\\
SOD & porous & 324 & 0.000 & 0.000 & 0.000 & 0.000 & 0.000 & 0.000 & 0.000 & 0.000 & 1.000 & 0.762\\
CRN1 & amorphous & 1000 & 0.097 & 0.097 & 0.046 & 0.095 & 0.107 & 0.032 & 0.047 & 0.012 & 0.066 & 0.067\\
CRN2 & amorphous & 1000 & 0.113 & 0.113 & 0.038 & 0.107 & 0.095 & 0.028 & 0.047 & 0.013 & 0.063 & 0.075\\
CRN3 & amorphous & 4096 & 0.048 & 0.048 & 0.020 & 0.048 & 0.041 & 0.012 & 0.031 & 0.005 & 0.086 & 0.086\\
\end{tabular}
}
\end{table}

Ions are the remaining negative case.
For Table~\ref{tab:ions} we replace water molecules of the GenIce ice Ih and Ic lattices and of the sI framework by Na and Cl at drawn lattice sites, between one and sixteen pairs per box.
The ions are then removed from the graph, and the water is labelled.
A substitutional ion leaves each of its four neighbours with three water neighbours.
CHILL+ requires four bonded neighbours, so it drops every water in the first shell of every ion.
Its ice fraction on ice Ih is \(0.972\) at two ions in \(288\) molecules and \(0.62\) at thirty-two.
The seeded cage assignment keeps the lattice through eight ion pairs (Ih \(0.999\), Ic \(0.997\)) and falls to \(0.922\) (Ih) and \(0.791\) (Ic) at sixteen.
The first-shell ice fraction is \(1.000\) through four pairs on Ih and falls to \(0.813\) at sixteen.
On Ih the cutoff graph starts at \(1.000\) and falls to \(0.796\), and on Ic it starts at \(1.000\) and falls to \(0.714\).
The vacancy opens rings that the cutoff graph cannot close.
The sI framework with ions stays at zero under all three.
The PLUMED action and the per-frame vectors of PydSEAMSlib take the ion types as an argument.
Per frame, they report how many ions have an all-ice shell, a mixed shell or none, which is the brine-rejection observable of a growing front.

\begin{table}[htbp]
\caption{\label{tab:ions}Fraction of water labelled ice I with Na and Cl at drawn lattice sites, means over three draws, for seeded cages with completion (seeded), cutoff cages (cutoff) and CHILL+. The last two columns give the mean ice fraction of the ions' first water shells under the seeded assignment and under CHILL+.
Ice Ih has \(288\) sites, Ic \(216\), the sI framework \(368\).}
\centering
\scriptsize
\resizebox{\textwidth}{!}{%
\begin{tabular}{lrrccccc}
host & pairs & ion fraction & seeded & cutoff & CHILL+ & shell, seeded & shell, CHILL+\\
\hline
Ih & 1 & 0.007 & 1.000 & 1.000 & 0.972 & 1.000 & 0.000\\
Ih & 2 & 0.014 & 1.000 & 0.997 & 0.944 & 1.000 & 0.000\\
Ih & 4 & 0.028 & 1.000 & 0.980 & 0.888 & 1.000 & 0.000\\
Ih & 8 & 0.056 & 0.999 & 0.955 & 0.792 & 0.990 & 0.000\\
Ih & 16 & 0.111 & 0.922 & 0.796 & 0.620 & 0.813 & 0.000\\
Ic & 1 & 0.009 & 1.000 & 1.000 & 0.963 & 1.000 & 0.000\\
Ic & 2 & 0.019 & 1.000 & 1.000 & 0.926 & 1.000 & 0.000\\
Ic & 4 & 0.037 & 0.998 & 0.992 & 0.857 & 0.979 & 0.000\\
Ic & 8 & 0.074 & 0.997 & 0.952 & 0.717 & 0.984 & 0.000\\
Ic & 16 & 0.148 & 0.791 & 0.714 & 0.507 & 0.593 & 0.000\\
sI & 1--16 & 0.005--0.087 & 0.000 & 0.000 & 0.000 & 0.000 & 0.000\\
\end{tabular}
}
\end{table}

The electrolyte models of the Madrid group are the target of that observable.
Section~\ref{sec:applications} reports a direct-coexistence ice/brine campaign on the scaled-charge Madrid-2019 ions of Zeron, Abascal and Vega \citep{zeron2019madrid} with TIP4P/2005 water \citep{sedano2026madridtopology}, and the occupancy of GenIce methane hydrates by ring-size signature.

The ion environment lists each ion's shell and the rings of the water network that pass through it.
A ring wholly in the shell is capped.
A ring that exists only once the ion is restored as a vertex bonded to that shell is broken.

Cutoff \(q_l\) is the Steinhardt path and does not replace Eq.~\mref{eq:ice}.
Voronoi-weighted \(q_l\), SOAP and IRA templates are available in the library; they are not the ice score.

\section{Shape matching}
\label{sec:org2fc88fa}

An RMSD against an ideal HC or DDC is a deformation score, not the ice score.
Forming that RMSD requires a point-to-point correspondence between the candidate cage and the template \citep{goswami2021tum}.
The correspondence is not known a priori for a generic point set.
Polyhedral template matching obtains it from the convex hull of the neighbour star \citep{larsen2016robust} and cannot map a non-convex HC, whose prismatic and basal planes are non-planar \citep{goswami2021tum}.
Exhaustive permutation of \(n\) vertices is \(n!\) and does not scale.
IRA \citep{gunde2021ira} is built for that generic case.
It iteratively proposes rotated atom-centered frames and assignments.
When the atomic assignment is unknown, it minimizes a permutationally invariant set--set distance.

In the TUM formalism the classification of the test block is already known from the topological criterion, and the relative order of its primitive rings is a topological information criterion that reduces the correspondence search \citep{goswami2021tum}.
For an HC, a triplet of alternating nodes in one basal ring is bonded to a corresponding triplet in the other.
Advancing the first basal ring by one vertex advances the second, and clockwise and anticlockwise orderings are equivalent.
We therefore test six of the \(12!\) permutations of a twelve-vertex cage.
A DDC is fixed in the same way by its equatorial ring and six peripherals.
Horn's quaternion algorithm \citep{horn1987closed} then returns the rotation
\begin{equation}
R(q)=(q_0^2-\mathbf{q}\cdot\mathbf{q})I
  + 2\mathbf{q}\mathbf{q}^{\mathsf T}
  + 2q_0[\mathbf{q}]_{\times}
\label{eq:horn}
\end{equation}
that minimizes the RMSD for each feasible mapping, and we retain the mapping of least RMSD.
We inspect a sequence of mapped cages with the same RMSD, plotted as a two-dimensional projection along a path \citep{goswami2026rmsd}.
A nudged-elastic-band path uses a related tangent and climbing-image construction \citep{henkelman1999dimer,henkelman2000neb,goswami2026cineb,goswami2026neb}.
Gaussian-process accelerations of those searches \citep{goswami2025gpr,goswami2025pruning,goswami2026pcau,goswami2025bayes} are not the ice score.

IRA \citep{gunde2021ira} recovers a correspondence the ring complex has already given.
SOFI \citep{gunde2024sofi} reports the point group of those self-matches.
An ideal DDC is tetrahedral.
Thermal motion reduces it to \(C_1\), and Eq.~\mref{eq:ice} is unchanged.
Neither is the ice score.
The cage loop parallelizes over cages.
On the 4096-molecule mW cubic frame the double-diamond cage (DDC) shape matcher matches 1950 cages and no hexagonal cage (HC), in \(32.5\) and \(44.8\,\mathrm{ms}\) at the engine commits of records \texttt{elja-1788072} and \texttt{elja-1798666} (files \texttt{tip-cages.txt}).

\section{Verification}
\label{sec:verification}
We check the compiled library against the published formulae, not against its own previous output.
Two disagreements with those formulae change labels or RMSD values, and we state both in this section.
We keep bounds checks, empty-container guards and similar regressions in the test suite (Section~\ref{sec:code}) and do not report them as results.
Figure~\ref{fig:verification} summarises the pipeline.

\begin{figure}[htbp]
\centering
\input{figures/verification}
\caption{How the compiled library is checked. Symbolic algebra with SymPy
derives the published formulae and compares them against a transcription of
the C++. Those scripts do not link the library. The Catch2 suite and the
ring and hydrogen-bond tests of PydSEAMSlib call the compiled symbols; its
Hypothesis tests of the minimum image and the Horn rotation run on Python
transcriptions.}
\label{fig:verification}
\end{figure}
\subsection{Published identities}
\label{sec:orgdd29cd9}

Eight scripts under \texttt{analysis/} cover stacking, hydrogen bonds, Horn orientation, outlier rejection, periodic distance, \(g(r)\), shoelace area and \(Y_{lm}\).
Each derives the intended formula with computer algebra and compares it term by term with a transcription of the corresponding C++ into Python.
They do not link the library.
The \(Y_{lm}\) script is Eqs.~\ref{eq:qlm}--\ref{eq:cs}; the Horn script is Eq.~\ref{eq:horn}; the periodic-distance script is Eq.~\ref{eq:simd_dist}.
They locate a wrong term in the transcription.
Catch2 cases that call the compiled symbols of Table~\ref{tab:symbols} pin the binary.

The spherical harmonics of Eq.~\ref{eq:qlm} must satisfy the addition theorem
\begin{equation}
\sum_{m=-l}^{l} Y_{lm}(\hat r)\,Y_{lm}^{*}(\hat r')
  = \frac{2l+1}{4\pi}\,P_l(\hat r\cdot\hat r')
\label{eq:addition}
\end{equation}
and the Condon--Shortley relation
\begin{equation}
Y_{l,-m}(\hat r)=(-1)^m Y_{lm}^{*}(\hat r)
\label{eq:cs}
\end{equation}
independently of any transcription.
We check the compiled tables against both identities and against high-precision values.
Emitting the orders in conjugate pairs makes Eq.~\ref{eq:cs} exact rather than to within the rounding of two independently evaluated exponentials.

Five of the nine matrix elements in the compiled map of Eq.~\ref{eq:horn} were written from the wrong identities, so the RMSD returned for a known correspondence was not Horn's minimum.
The transcription now matches Eq.~\ref{eq:horn} term for term.

CHILL and CHILL+ define \(N_b(i)=4\) in Eq.~\ref{eq:qlm} \citep{moore2010chill,nguyen2015identification}.
The compiled correlation walk averaged \(q_{lm}\) over every neighbour inside the distance cutoff.
On ordered ice the two averages coincide.
On the disordered example trajectory they do not.
There \(17.6\%\) of the molecules (44 of 250) have two, three, five or six neighbours inside \(3.5\,\mathring{\mathrm{A}}\) rather than four (\texttt{scripts/example\_neighbour\_count.py}).
Commit \texttt{73b7a41d} moves the CHILL+ average to the four nearest neighbours, and against its parent \texttt{4899a598} it changes the label of 7 of the 250 molecules (\(2.8\%\)).
One clathrate becomes interfacial clathrate, one interfacial clathrate becomes water, and five water become interfacial clathrate.
The four-neighbour \(q_6\) of \texttt{7b8ed96b} changes none of the 250 against its parent \texttt{85f8c0b0} (rule \texttt{reclass} of the rerun of Section~\ref{sec:s-rerun}, file \texttt{results/reclass/reclass.txt}).
Relative to the cutoff list, the mutual four-nearest list changes the CHILL+ label of 6 of the 250 molecules, five interfacial clathrate to water and one water to clathrate (\texttt{scripts/example\_chill\_graphs.py}, file \texttt{repro-records/example\_chill\_graphs.txt}).

The cell-list path and the in-plane \(g(r)\) reduce the cell list's periodic images to a single minimum-image pair (Eq.~\ref{eq:simd_dist}) and never the self image, with a regression across boxes larger and smaller than twice the cutoff.

The shortest-path construction (Eqs.~\mref{eq:franzblau}, \mref{eq:yuan}) agrees with the generate-then-filter path of \citet{goswami2020dseams}.
That path remains available in the engine and produces the same ring set.
The affiliation predicate agrees ring for ring with the assembled HC and DDC searches on mW cubic (every six-ring DDC-affiliated, none HC) and on the isolated hexagonal cage of a Catch2 fixture.
The Catch2 case ``DDC needs six distinct peripheral rings'' builds two ring sets.
In these sets one candidate passes the first two equatorial tests, and it passes the intersections of the third only with a peripheral used twice, as \(R'_1=R'_0\) or as \(R'_2=R'_0\).
The host predicate, the DDC search and the offload test reject that candidate, and \texttt{tests/gpu\_kernels\_host.cpp} requires the same of the CUDA kernel.
All four accept it, with its six peripherals, once one more ring supplies a sixth distinct peripheral.
When we remove the distinctness test from any of the host, offload or CUDA code, the case fails.
A Catch2 case that translates every molecule at graph distance two from a chosen \(v\) asserts the star of \(v\) identical and \(\mathrm{Ic}(v)\) flipped.
\subsection{Property-based tests}
\label{sec:org592d063}

The four \texttt{test\_hypothesis\_*} modules of PydSEAMSlib 2.11.1 hold 32 tests of invariants that the identities above imply but do not exhaust.
Of these, 15 are Hypothesis property tests of the minimum image and the Horn rotation, and 17 are fixed-input tests of hydrogen bonds, rings and the Horn map.
Eighteen run on Python transcriptions of the minimum image and the Horn rotation.
Fifteen of those draw random input.
Ten of them test symmetry, bounds, triangle inequality and periodicity of the minimum-image distance.
The other five test \(\det R(q)=1\), \(R(q)^{\mathsf T}R(q)=I\), the inverse, norm preservation and composition of the Horn rotation.
The remaining three check fixed rotations.
Fourteen call the compiled ring and hydrogen-bond builders and check the cycle validity of every emitted ring and the symmetry of the hydrogen-bond network on the example trajectory.

\subsection{Corrections that change results}
\label{sec:s-corrections}
Table~\ref{tab:corrections} lists the corrections since d-SEAMS 1 that change a label, a count or a distance.
None of them is a new method.

\begin{table}[htbp]
\caption{\label{tab:corrections}Corrections since d-SEAMS 1 that change scientific output, with the seams-core commits and sources that carry them.}
\centering
\scriptsize
\setlength{\tabcolsep}{3pt}
\begin{tabular}{L{4.4cm}L{4.0cm}L{4.6cm}}
correction & where & effect\\
\hline
CHILL and CHILL+ average \(q_{lm}\) over exactly the four nearest neighbours; d-SEAMS 1 looped over the whole cutoff row & \texttt{src/bop.cpp}; \texttt{73b7a41d}, \texttt{7b8ed96b} & the two averages differ on the 44 of 250 molecules of the example trajectory that have other than four neighbours inside \(3.5\,\mathring{\mathrm{A}}\), and \texttt{73b7a41d} changes 7 CHILL+ labels (Section~\ref{sec:orgdd29cd9})\\
full triclinic minimum image replaces independent wrapping of the LAMMPS bound spans & \texttt{generic.hpp}; \texttt{e00a39d1}, \texttt{479baf20}, \texttt{f0eb8c3d} & minimum-image property tests (Section~\ref{sec:org592d063})\\
the in-plane RDF normalises by the dump cell instead of the particle bounding box & \texttt{src/rdf2d.cpp}; \texttt{70da7368} & confined-liquid peak at \(2.825\,\mathring{\mathrm{A}}\) (Table~\ref{tab:deposits})\\
DDC cage objects are assembled in a position-keyed order instead of the ring-enumeration order & \texttt{src/topo\_bulk.cpp}; \texttt{37a71a6c} & 2040 cages in every ring order, against 1955 in the natural order of the baseline (Section~\ref{sec:s-invariance})\\
a corrupted atom ID in the shipped mW cubic frame is repaired, and duplicate IDs warn & input fixture and reader; \texttt{2a8d0051} & 8192 rings on that frame instead of 8180\\
\texttt{projAreaSingleRing} no longer swaps the XZ and YZ projections & \texttt{src/order\_parameter.cpp}; \texttt{0ad6a89f} & symbolic and regression tests\\
five of the nine quaternion rotation-matrix elements and the TUM scoring paths that read them & \texttt{src/absOrientation.cpp}; \texttt{0ad6a89f}, \texttt{e7f1c68e}, \texttt{52aec4e4} & identities of Section~\ref{sec:orgdd29cd9} and the Horn tests\\
\end{tabular}
\end{table}

\section{Code Overview}
\label{sec:code}
\subsection{Architecture}
\label{sec:org988c4f1}


The original d-SEAMS release exposed a single \texttt{yodaStruct} executable that mediated YAML--Lua scripts and the C++ back-end \citep{goswami2020dseams}, packaged with Nix \citep{dolstra2004nix,goswami2022nix,goswami2022phylo}.
The same packaging pattern, a compiled library with a command-line interface and a pinned release, is how Flowy is released as a CPC program paper \citep{sallermann2025flowy} and how the finite-element atomic structure engine is released \citep{certik2024featom}. d-SEAMS 2.0 is the same analysis as a C++20 library, \texttt{libyodaLib}, with a \texttt{seams} command-line interface, the Python package PydSEAMSlib (\texttt{import pydseamslib}) and a Lua library loaded with \texttt{require("dseams")} (Figure~\ref{fig:architecture}).
The live builds are pixi \citep{pixi} with Meson \citep{meson}, or the Nix flake.
The release builds against minimage \citep{goswami2026minimage} and linkcell \citep{goswami2026linkcell}, with vesin \citep{vesin} and Catch2 \citep{catch2}.
Highway \citep{highway} and a CUDA device path are optional, and the pixi lock of the reproducibility package pins the version of every dependency.
The PydSEAMSlib wheels are built with cibuildwheel in its repository.
The closed-form \(Y_{lm}\), Horn orientation and topological cage searches remain the compiled path if optional libraries are absent. sphericart \citep{bigi2023sphericart}, IRA/SOFI \citep{gunde2021ira,gunde2024sofi} and nauty \citep{mckay2014nauty} can be compiled in.
Parsers for the quantum-chemistry side of that stack are Wailord \citep{goswami2022wailord}.
Table~\ref{tab:symbols} maps the compiled identifiers to the equations they implement.

\begin{figure}[htbp]
\centering
\resizebox{\textwidth}{!}{\input{figures/architecture}}
\caption{Architecture of d-SEAMS 2.0. Trajectory data enter the C++20
library \texttt{libyodaLib} through the \texttt{seams} CLI,
PydSEAMSlib (\texttt{import pydseamslib}) in Python, or
\texttt{require("dseams")} in Lua. The library dispatches to CHILL/CHILL+,
hexagonal and double-diamond cage tests, and prism identification
(the confined-film criterion of \citet{goswami2020nanoribbon}). Regression
tests (dashed) exercise the compiled library. Compare with the d-SEAMS 1
architecture in Ref.~\cite{goswami2020dseams}, where a single
\texttt{yodaStruct} executable mediated YAML-Lua scripts and the C++
back-end.}
\label{fig:architecture}
\end{figure}

\begin{table}[htbp]
\caption{\label{tab:symbols}Compiled symbols and the equations they implement.}
\centering
\scriptsize
\begin{tabular}{ll}
symbol & equation / role\\
\hline
\texttt{kNearestNeighbourList} & \(k\)-NN walk, Sec.~\mref{sec:algorithms}\\
\texttt{knearest\_into} & wrap of that walk; in-tree if the wrap is absent\\
\texttt{BatchPeriodicDistSq} & Eq.~\ref{eq:simd_dist}\\
\texttt{cageAffiliation} & Eq.~\mref{eq:ice}, batch\\
\texttt{AffiliationUpdater} & Eq.~\mref{eq:ice}, incremental\\
\texttt{ringNetwork} & Eqs.~\mref{eq:franzblau}, \mref{eq:yuan}\\
\texttt{getCorrel} / \texttt{getCorrelPlus} & Eqs.~\ref{eq:cij}--\ref{eq:eclipsed-plus}\\
\texttt{getIceTypeNoPrint} / \texttt{getIceTypePlusNoPrint} & Eqs.~\ref{eq:chill-ic}--\ref{eq:chill-interclath}\\
\texttt{buildRingSearchIndex} & Eq.~\mref{eq:index}\\
\texttt{findHC} / \texttt{findDDC} & \(\beta\), \(\varepsilon\)\\
\texttt{calcMatrixS} / \texttt{calcMatrixN} & Horn \(S\) and \(N\), Eq.~\ref{eq:horn}\\
\texttt{hornAbsOrientation} & \(R(q)\) and the RMSD, Eq.~\ref{eq:horn}\\
\end{tabular}
\end{table}
\subsection{Installation and a first run}
\label{sec:install}

The engine builds with pixi and Meson, or with the Nix flake, and the library is C++20.
The pixi environment pins the conda-forge compilers, Eigen, Catch2 and Meson.
PydSEAMSlib installs with \mbox{\texttt{pip install pydseamslib}} and loads with \mbox{\texttt{import pydseamslib}}, and the Lua library builds from its repository with Meson.
The Meson option \texttt{with\_gpulite} builds the CUDA path of Section~\ref{sec:resident} when CUDA and NVRTC are found.
The LAMMPS compute is the two sources \texttt{extras/lammps/compute\_dseams.cpp} and \texttt{compute\_dseams.h}, which call only the C header \texttt{seams\_c\_api.h}.
They are copied into the \texttt{src} directory of LAMMPS, which links to \texttt{libyodaLib}.
The PLUMED module builds in its own repository.
The release installs with the Meson commands of Section~\ref{sec:s-repro}, and the LAMMPS and PLUMED build steps are in the README of each repository.
\subsection{Python Interface}
\label{sec:org46c0399}

\begin{lstlisting}[language=Python,numbers=none]
import pydseamslib as ds

frame = ds.read("dump.lammpstrj", frame=100)
print(frame.chill_plus())
print(frame.cages())
\end{lstlisting}

The Python handle builds neighbour lists, hydrogen bonds and rings on demand.
With chemfiles installed it also reads every format that chemfiles reads.
We view classified frames with solvis \citep{goswami2024solvis}.
The same package includes the solvation-shell constructions used for aqueous Fe\(^{3+}\) \citep{goswami2024fe3} and the later viscosity diagnostics of those shells \citep{goswami2026viscosity}.
Cage stills, ice I lattices and the inserted seed (Figures~\mref{fig:cages}, \mref{fig:ice-lattices} and \mref{fig:seed-xyz} of the main text) use that AtomicPlotter path.
The original d-SEAMS release rendered snapshots in OVITO and VMD \citep{goswami2020dseams,stukowski2010visualization,humphrey1996vmd}.
Polyhedral template matching in the comparison of Section~\mref{sec:accuracy} is the Larsen classifier \citep{larsen2016robust}.
The engine CLI is the same path without Python.

\begin{lstlisting}[language=bash,numbers=none]
seams read dump.lammpstrj
seams chill-plus dump.lammpstrj
seams cages dump.lammpstrj
seams cages dump.lammpstrj --graph cutoff
seams fingerprint dump.lammpstrj
seams rdf dump.lammpstrj --types 1,2
seams cn dump.lammpstrj --types 1,2
seams density-z dump.lammpstrj
seams domains dump.lammpstrj --site '1=cationHead,2=anion,3=tail' --subset polar
\end{lstlisting}

\texttt{seams chill-plus} evaluates CHILL+\((r_c)\), and \texttt{-{}-graph knn} gives CHILL+(4NN).
\texttt{seams cages} defaults to the seeded rule with \(k=4\) and no ring completion; \texttt{-{}-complete} adds completion, and \texttt{-{}-graph cutoff} evaluates the HC and DDC predicates on the \(3.5\,\mathring{\mathrm{A}}\) cutoff graph.
With \texttt{SEAMS\_OFFLOAD} set to a nonzero value, the cutoff and \(k\)-nearest cage labels run on the device kernels; the seeded rule and its completion run on the host.

The analysis of Section~\ref{sec:applications}, the cage assignment on every frame of the mW nucleation deposit, is the same three calls on each surface.
The command line walks the frames in parallel and emits one line per frame.

\begin{lstlisting}[language=bash,numbers=none]
seams cages nucleation.lammpstrj --type 1 -k 4 --graph seeded \
    --frame 1 --last 148 --jobs 8
\end{lstlisting}

In Python one \texttt{Trajectory} handle loads each frame in place, and \texttt{seeded\_affiliation()} builds the mutual and union graphs and their rings for that frame.

\begin{lstlisting}[language=Python,numbers=none]
from pydseamslib import Trajectory

traj = Trajectory("nucleation.lammpstrj", frame=1, atom_type=1, cutoff=3.5)
for f in range(1, 149):
    traj.load_frame(f)
    score = traj.seeded_affiliation()
    print(f, score.n_ic, score.n_ih, score.n_water)
\end{lstlisting}

In Lua the same assignment is one helper over a cloud.
It returns one Boolean per molecule for each cage type.

\begin{lstlisting}[language=lua,numbers=none]
local dseams = require("dseams")
local function count(flags)
  local n = 0
  for _, on in ipairs(flags) do
    if on then n = n + 1 end
  end
  return n
end
for f = 1, 148 do
  local cloud = dseams.read("nucleation.lammpstrj", {type = 1, frame = f})
  local cages = dseams.cages(cloud, {type = 1})
  print(f, count(cages.ddc), count(cages.hc))
end
\end{lstlisting}

The three loops evaluate the same seeded rule without completion; \texttt{-{}-complete}, \texttt{ring\_adjacent=True} and \texttt{complete = true} add it.
CHILL+ is \texttt{Trajectory.\allowbreak classify\_chill\_plus()} in Python and \texttt{dseams.\allowbreak chill\_plus(cloud)} in Lua, both CHILL+\((r_c)\) on the \(3.5\,\mathring{\mathrm{A}}\) cutoff list.
The per-frame tables of Section~\ref{sec:applications} are written by the in-tree \texttt{walk\_compare} driver, which at the engine commits of Table~\ref{tab:records-full} evaluates CHILL+\((r_c)\) next to the cage labels.

Two more surfaces put the assignment into the tools a nucleation study already runs.
\texttt{compute dseams} in \texttt{extras/lammps} writes a per-atom column of CHILL+ integers through the C ABI function \texttt{seams\_chill\_plus\_graph}, the integer table of \texttt{seams chill-plus}.
Its optional keywords are \texttt{graph cutoff} or \texttt{graph mutual} and \texttt{cutoff}, and the defaults are the cutoff graph at \(3.5\,\mathring{\mathrm{A}}\), the list of \texttt{seams chill-plus}.
It copies the coordinates of the \texttt{nlocal} atoms of its rank and the orthorhombic box lengths, and it reads no LAMMPS neighbour list, ghost atom or tilt factor (\texttt{extras/lammps/compute\_dseams.cpp}), so it serves single-domain orthorhombic runs.
The example \texttt{extras/lammps/in.dseams} loads it as a plugin and dumps the labels.

\begin{lstlisting}[language=bash,numbers=none]
plugin load ./compute_dseams.so
compute ice all dseams
dump 1 all custom 1 dump.chill id type x y z c_ice
\end{lstlisting}

\texttt{DSEAMS\_CAGES} is a PLUMED action \citep{tribello2014plumed} in a loadable module (\url{https://github.com/d-SEAMS/dseams-plumed}).
Given the oxygen atoms and the cell, it reports at every stride the molecules in a cage, the largest cage cluster, the split into Ic, Ih and mixed, and CHILL+\((r_c)\) on the same frame.
With \texttt{IONS} it reads each ion against the assignment as in ice, at the front or in liquid.
With \texttt{LIBRARY} it counts the molecules whose local topology key (Section~\ref{sec:topokeys}) a reference library names.
Any polymorph or hydrate framework then becomes a collective variable.
The counts are integers without derivatives, so they drive \texttt{PRINT}, \texttt{COMMITTOR} basins and stopping conditions rather than a bias.
The seeding runs of Section~\ref{sec:applications} read \texttt{templates/plumed\_seed.dat} of the reproducibility package, with the placeholders filled for each run.

\begin{lstlisting}[numbers=none]
UNITS LENGTH=A
LOAD FILE=@MODULE@
ice: DSEAMS_CAGES ATOMS=1-@NATOMS@ CUTOFF=3.5 CANDIDATE=5.5 K=4 LENGTH_SCALE=1.0 COMPLETE
PRINT ARG=ice.nice,ice.nmax,ice.nclus,ice.nic,ice.nih,ice.nmixed,ice.chillice,ice.chillmax,ice.chillinterfacial,ice.sixrings STRIDE=@STRIDE@ FILE=ICE FMT=%8.0f
COMMITTOR ARG=ice.nmax STRIDE=@STRIDE@ BASIN_LL1=0 BASIN_UL1=@MELT@ BASIN_LL2=@GROW@ BASIN_UL2=1000000 FILE=COMMITTOR
\end{lstlisting}

The module's \texttt{tests/driver\_check.py} runs \texttt{plumed driver} with \texttt{DSEAMS\_CAGES} over a LAMMPS dump and compares the per-frame cage count, largest cluster, CHILL+ ice count and CHILL+ cluster with a \texttt{walk\_compare} table.
PydSEAMSlib also returns per-frame vectors and per-molecule states, with adapters for deeptime and for PyEMMA (in maintenance mode since 2022), and a one-call discretisation of the largest-cluster series.
\subsection{The carried command line}
\label{sec:s-carry}
The \texttt{seams} command line of \coretag{} holds one ring updater for each graph and ring size and one affiliation updater for each graph, and it passes every frame of a range through them (Section~\mref{sec:carry}, Figure~\ref{fig:carry}).
In v2.11.0 (\texttt{9f963b72}) the command line rebuilt the rings and the cage flags of every frame (\texttt{src/seams\_cli.cpp}, lines 1361--1390 at that tag), and \texttt{RingUpdater} and \texttt{AffiliationUpdater} were library classes that only the tests and benchmarks called.
The Meson test \texttt{seams\_carried\_frames\_match\_rebuild} compares the JSON output of a carried range with \texttt{-{}-rebuild}, \texttt{SEAMS\_CARRY=0} and \texttt{-{}-jobs 2}.
It runs the seeded, cutoff and four-nearest graphs, ring completion, three cage signatures, the TUM layers and the topology keys.
Its inputs are the mW crystal, a second bundled LAMMPS dump and a generated trajectory in which three atoms drift by \(0.6\,\mathring{\mathrm{A}}\) per frame.
Every output is identical.

\begin{figure}[htbp]
\centering
\resizebox{\textwidth}{!}{\input{figures/carry}}
\caption{Carried walk of \texttt{seams} over a frame range. Each frame
builds its neighbour graph, and the ring update and the cage update
read its edge diff against the stored state of the previous frame.
Each updater falls back to the whole-frame pass when most of the frame
changes, and after such a pass it skips the diff for the next 1, 2, 4
and then at most 8 frames. With \texttt{-{}-rebuild}, frame-parallel
workers or MPI ranks every frame takes the whole-frame passes. Every
path returns the output of a rebuild.}
\label{fig:carry}
\end{figure}

\subsection{Further analyses in the library}
\label{sec:s-further}
The library carries analyses that this paper does not benchmark.
\texttt{rodgerF4} returns the \(F_4\) parameter of \citet{rodger1996f4} per oxygen, the mean of \(\cos 3\phi\) over the H--O\(\cdots\)O--H dihedrals with its oxygen neighbours, and it returns no value on a model without hydrogens such as mW.
With TIP3P hydrogens on the ice-rule donation plane and a \(15^{\circ}\) libration about the HOH bisector, the Catch2 cases give a mean of \(-0.501\) on ice Ih and \(0.691\) on filled sI (\texttt{tests/test\_order\_parameter.cpp}, lines 446--452 and 469--479).
With collinear hydrogens on the O--O bonds sI gives about \(0.90\) (lines 277--280 of that file).
\texttt{glassFromDensity} separates ice from low-, medium- and high-density amorphous ice by the local number density, with default windows that end at \(0.028\), \(0.040\) and \(0.070\,\mathring{\mathrm{A}}^{-3}\).
\texttt{findIncompleteBySignature} returns the connected face sets of a cage signature that do not close, the cups and incomplete cages of hydrate nucleation.
\texttt{iceXXILibrary} recognises the 152-molecule ice XXI cell by its molecule count, cell lengths, density, mean coordination, six-ring census and number of local topology classes, so it checks that cell and does not label molecules.

\subsection{Testing}
\label{sec:org4b4bd07}

The Catch2 suite \citep{catch2} is 41 test binaries with 509 test cases, plus \texttt{test\_readcon} (6 test cases) when readcon is built (\texttt{tests/meson.build}, lines 49--51); at \coretag{} all 42 binaries pass.
It covers trajectory I/O, rings, cages, the incremental updaters and CHILL/CHILL+ on mW and TIP4P.
\texttt{meson test} and \texttt{nix build} run it, together with \texttt{check\_device\_host}, \texttt{check\_invariance} and \texttt{gpu\_kernels\_host}.
Property tests of Section~\ref{sec:verification} call the symbols of Table~\ref{tab:symbols}.
SymPy \citep{meurer2017sympy} checks the closed-form \(Y_{lm}\) identities.
Optional SIMD paths use Highway \citep{highway}; dump numbers are parsed by the Lemire \texttt{from\_chars} path \citep{lemire2021number}.
The compiled stack sits next to NumPy and SciPy in the Python bindings \citep{harris2020numpy,virtanen2020scipy} and next to the Fortran scientific libraries as a compiled neighbour \citep{kedward2022fortran}.
Theorems S1--S34 of this supplement are 34 of the 53 public theorems of a Lean~4 development \citep{moura2021lean4} over abstract graphs, ring lists, labels and batch arithmetic.
For an enumerator assumed local, it proves ball agreement under the source bound and the carry-or-rebuild case split.
It gives the list lengths of the even and odd Yuan candidates, and the list lengths and stored fields of the HC and DDC structures.
It proves witness and locality statements for simplified affiliation predicates and two statements about the seeded predicate.
It characterises the completion as the least closed set of its inductive definition and proves natural-number inequalities for the batch size and the grid.
The development is the \texttt{lean/} directory of the supplementary material (\texttt{lean.zip}).
We check it with \texttt{lake build} on Lean 4.15.0, with no other library and no \texttt{sorry}.
The \texttt{lean/} directory of seams-core is a separate development that this paper does not use.
It does not model the C++ code and is not required to run the engine.
The property tests use Hypothesis \citep{maciver2019hypothesis} and run in the PydSEAMSlib repository with \texttt{pytest}.

\section{Further benchmarks}
\label{sec:s-bench}
The ring-and-cage stages on the 4096-molecule mW cubic frame take \(2864.284\,\mathrm{ms}\) in the baseline and \(46.821\,\mathrm{ms}\) in 2.0 on pinned core 3 of the cloud VM (Figure~\mref{fig:pipeline}, record \texttt{bench-2026-10-31}), and every six-ring remains DDC-affiliated.
We time neighbour-list construction, ring detection and cage classification separately below.
Figure~\ref{fig:scaling} and Table~\ref{tab:scaling} give the comparison against system size.

We take seams-core commit \texttt{1c8efa38} (2026-05-28), between v2.0.1 and v2.1.0, as the baseline of every comparison in this section.
It does not contain the searches of Section~\mref{sec:algorithms}.
It keeps the ring enumeration of d-SEAMS 1 (closed walks, then the primitive-ring filter), the pair-scan HC search followed by the DDC search, and the all-pairs index-ordered neighbour builder.
Its cutoff list already takes pairs from vesin.
It then maps every row to an atom ID by a linear search of the ID map.
Columns, curves and bars marked ``baseline'' are that commit.
The 2020 release differs from it in two places.
Its cutoff list is the all-pairs loop, and it runs the DDC search before the HC search.
Table~\ref{tab:records-full} lists the campaign records that this section and Section~\ref{sec:applications} read, with the node, its load and the engine commit each built.

\subsection{Machines and protocol}
\label{sec:s-protocol}
Every record of Table~\ref{tab:records} runs seams-core \coretag{}, source tree \coretree, built from the program-files archive, and each \texttt{run.sh} refuses a tree that differs from that archive.
They all ran on one machine, a KVM guest with 8 vCPUs on one socket and one thread per core.
Its CPU is an Intel Xeon of CPUID family 6 and model 207 with \(320\,\mathrm{MiB}\) of L3, and it has \(47\,\mathrm{GiB}\) of memory.
We time on its pinned core 3 with \texttt{OMP\_NUM\_THREADS=1} and run its eight-thread timings on cores 0--7.
The compiler and library versions are in the \texttt{provenance.txt} of each record.
Both libraries are Meson release builds with g++ 13.3.0.
The engine's drivers are its benchmark sources at \texttt{-O3 -funroll-loops}, and the baseline's drivers are the same sources compiled against it at \texttt{-O2}.
The cage and scaling drivers report the minimum over repetitions, and the trajectory driver reports the mean over frames 2--11.
Record \texttt{bench-2026-10-31} runs the cage, trajectory and jittered-cell drivers in three sessions in turn, and we quote the median over sessions.
The command-line timings of Section~\mref{sec:incremental} are whole processes, and each includes reading the dump and writing the last output line.

\begin{table}[htbp]
\caption{\label{tab:records}Records behind the timings of the main text, under \texttt{repro-records/} of the supplementary data.
Every record runs seams-core \coretag{}, source tree \coretree, on the 8-vCPU Xeon VM.
Load is the one-minute load average when the record started, and includes the tail of the work before it.
Table~\ref{tab:records-full} lists every record of this supplement.}
\centering
\footnotesize
\setlength{\tabcolsep}{4pt}
\begin{tabular}{L{3.0cm}L{3.1cm}L{0.8cm}L{5.4cm}}
record & cores & load & read by\\
\hline
\texttt{bench-2026-10-31} & pinned core 3; cores 0--7 & 0.00 & Figure~\mref{fig:pipeline}, ring scaling, Figure~\ref{fig:incremental}, jittered cell, churn sweep, deposited trajectories\\
\texttt{cli-carry-}\allowbreak\texttt{2026-10-31} & pinned core 3 & 0.77 & Table~\ref{tab:carry}\\
\texttt{h2h-2026-10-31} & pinned core 3 & 1.67 & Table~\mref{tab:compare}\\
\texttt{incremental-md-}\allowbreak\texttt{2026-10-31} & pinned core 3 & 1.00 & MD stride trajectories, nucleation excerpt\\
\texttt{asv-churn-}\allowbreak\texttt{2026-10-31} & pinned core 3 & 0.16 & churn sweep as an asv suite (Section~\ref{sec:s-incremental})\\
\end{tabular}
\end{table}

\begin{table}[htbp]
\caption{\label{tab:records-full}Campaign records behind the timings and labels, as directories or files under \texttt{results/reference/} of \url{https://github.com/d-SEAMS/dseams2_repro} at commit \texttt{\reprobase} with the package patch applied.
The records dated 31 October 2026 are also in \texttt{repro-records/} of the supplementary data.
Every Elja CPU timing record built baseline commit \texttt{1c8efa38} next to the engine commit listed; the GPU job built only the engine.
Load is the one-minute load average when the record started; \texttt{elja-1788072} gives it before the workflow.}
\centering
\scriptsize
\setlength{\tabcolsep}{3pt}
\begin{tabular}{L{2.7cm}L{2.1cm}L{2.9cm}L{1.8cm}L{3.1cm}}
record & processor & allocation & engine & read by\\
\hline
\texttt{elja-1788072} & 2\(\times\) Xeon Platinum 8358 & exclusive \texttt{64cpu\_256mem}, load 0.95 & v2.6.0 sources & Table~\ref{tab:scaling}, Figure~\ref{fig:scaling}, cage stages\\
\texttt{paper\_manifest-}\allowbreak\texttt{elja-1785712.json} & 2\(\times\) Xeon Platinum 8358 & \texttt{64cpu\_256mem}, load 4.96 & \texttt{13c95cec} & combined cage call on Elja\\
\texttt{elja-1798666} & 2\(\times\) AMD EPYC 7713 & exclusive \texttt{128cpu\_256mem}, load 0.24 & \texttt{e7856408} & Table~\ref{tab:second-arch}, \(q_l\) check, walk control\\
\texttt{terra-2026-09-01} & Ryzen 9 9950X & shared workstation, load 4.32 & \texttt{770accf9} & walk and PTM controls\\
\texttt{sota-heldout-}\allowbreak\texttt{2026-10-31} & Xeon, 8-vCPU cloud VM & labels only, no timing & \coretag & Table~\ref{tab:sota}, Figure~\mref{fig:sota}\\
\texttt{md-heldout-}\allowbreak\texttt{2026-10-31} & Xeon, 8-vCPU cloud VM & labels only, no timing & \coretag & Table~\ref{tab:md}\\
\texttt{h2h-2026-10-31} & Xeon, 8-vCPU cloud VM & pinned core 3, load 1.67 & \coretag & Table~\mref{tab:compare}, timings beside Table~\ref{tab:sota}\\
\texttt{bench-2026-10-31} & Xeon, 8-vCPU cloud VM & pinned core 3 and cores 0--7, load 0.00 & \coretag & Figure~\mref{fig:pipeline}, Figure~\ref{fig:incremental}, incremental and churn timings\\
\texttt{cli-carry-}\allowbreak\texttt{2026-10-31} & Xeon, 8-vCPU cloud VM & pinned core 3, load 0.77 & \coretag & Table~\ref{tab:carry}\\
\texttt{asv-churn-}\allowbreak\texttt{2026-10-31} & Xeon, 8-vCPU cloud VM & pinned core 3, load 0.16 & \coretag & churn sweep as an asv suite\\
\texttt{md-strides-}\allowbreak\texttt{2026-10-31} & Xeon, 8-vCPU cloud VM & trajectories only, no timing & none & stride trajectories\\
\texttt{incremental-md-}\allowbreak\texttt{2026-10-31} & Xeon, 8-vCPU cloud VM & pinned core 3, load 1.00 & \coretag & stride and nucleation updates\\
\texttt{hydrate-inside-}\allowbreak\texttt{2026-10-31} & Xeon, 8-vCPU cloud VM & labels only & \coretag & ray-parity hydrate occupancy\\
\texttt{walks-2026-10-31} & Xeon, 8-vCPU cloud VM & labels only & \coretag & label walks, Table~\ref{tab:public-cnt}, Figures~\mref{fig:app-nucleation}, \mref{fig:app-public} and \ref{fig:app-msm}\\
\texttt{terra-2026-09-02} & Ryzen 9 9950X & shared workstation & \texttt{0a7fc546} & Figure~\ref{fig:strong}, Tables~\ref{tab:atlas} and \ref{tab:ions}, hydrate occupancy\\
\texttt{tip-gpu-batch-}\allowbreak\texttt{elja-1787449.txt} & 2\(\times\) A100-PCIE-40GB & \texttt{gpu-2xA100}, load 0.00 & not recorded & Section~\ref{sec:resident}\\
\texttt{elja-production} & 2\(\times\) Xeon Platinum 8358 & exclusive \texttt{64cpu\_256mem}, one job per temperature & v2.9.2 (\texttt{d1983d53}) & seeding and brine campaigns\\
\end{tabular}
\end{table}

We take each C++ benchmark-driver timing in this section as a \texttt{std::chrono::}\allowbreak\texttt{steady\_clock} interval around the library call and the assignment or clearing of its output vectors.
The frame is already in memory, and the engine's cage driver builds the inverted index once before its timed calls.
The cage drivers of the cited release time \texttt{findDDC} alone.
They compute the ring types and cages that \texttt{findHC} leaves once, and restore them outside the timer before each call.
The combined HC and DDC call is timed separately.
On pinned core 3 of the cloud machine the combined call of the engine takes \(43.403\,\mathrm{ms}\) against \(40.798\) and \(2.489\,\mathrm{ms}\) for the two calls timed alone (record \texttt{bench-2026-10-31}).
The drivers of record \texttt{paper\_manifest-}\allowbreak\texttt{elja-1785712.json}, at engine commit \texttt{13c95cec} on an Elja Xeon Platinum 8358 node, timed only the HC call and the combined call, which built its local ring-type and cage vectors inside the timer.
There the two cage searches together take \(48.6\,\mathrm{ms}\) in the engine and \(3.9\,\mathrm{s}\) in the baseline.
With the ring construction of the same drivers, \(13.9\) and \(100.7\,\mathrm{ms}\), the path takes \(62.5\,\mathrm{ms}\) in the engine and \(3979\,\mathrm{ms}\) (\(4.0\,\mathrm{s}\)) in the baseline.
The scaling, cage and strong-scaling drivers report the minimum over repetitions.
The synthetic incremental driver reports the smallest update and the smallest rebuild over its five frames, each taken separately.
The trajectory driver takes one sample per frame and reports the mean over frames 2--11.
The scaling driver takes five repetitions up to 4000 atoms and three above, except ring construction, which runs once at 16000 and 32000 atoms; the cage driver takes five.
On Elja the engine is a Meson release build with the conda-forge GCC toolchain of its pixi environment.
The baseline library is a release build with the same toolchain, and its drivers are the engine's benchmark sources compiled against it with \texttt{g++ -std=c++20 -O2}.
Every timed function of the baseline is defined in its library sources, so that flag reaches only the drivers' timers, output vectors and call sites.
We run the single-core timings of the Elja records with \texttt{OMP\_NUM\_THREADS=1} as one-CPU HyperQueue tasks inside the Slurm allocation.
Every single-core row of those records is one core of a dual-socket Intel Xeon Platinum 8358 (32 cores per socket, \(2.60\,\mathrm{GHz}\)).
Record \texttt{bench-2026-10-31} runs one-thread timings on pinned core 3 and eight-thread trajectory timings on cores 0--7 of the 8-vCPU cloud VM of Table~\mref{tab:compare} (Intel Xeon, family 6, model 207).
Its system GCC 13.3 builds both the engine and the baseline.
We report the cage, mW-trajectory and synthetic-incremental values as medians over three sessions, and the nucleation, Niu and churn values from one session.
The scaling systems are synthetic cubic cells at the mW number density of \(0.0332\,\mathring{\mathrm{A}}^{-3}\) with a deterministic jitter.
The mean coordination number is therefore the same at every size, and the number of bonds grows in proportion to the atom count.
\subsection{Neighbour List Construction}
\label{sec:org062385e}

The cutoff list of the original d-SEAMS release is the all-pairs minimum-image loop. \citet{heinz2004pairlist} review replacing that double loop by a grid of cells.
A linked-cell list scales linearly \citep{allentildesley2017}, but most of the pair distances it still evaluates lie outside \(r_c\) \citep{welling2011efficiency}. \citet{vesin} is that cutoff pair list in the engine commits timed here, with a brute-force fallback.
It does not return the \(k\) nearest neighbours.
The seeded, mutual and union graphs of Section~\mref{sec:algorithms} take those \(k\) nominations from linkcell \citep{goswami2026linkcell}, or from the in-tree fallback above if the wrap is absent.
That walk has no cutoff.
Its shells grow until the \(k\)-heap is exact.

Neighbour lists are keyed by atom ID while coordinates are addressed by index.
The baseline resolves that map by a linear search for every row, so its cutoff list stays quadratic although vesin supplies the pairs.
A table built once per frame makes the map constant-time, and the measured scaling is close to linear (Table~\ref{tab:scaling}).
Between 1000 and 32000 atoms, a thirty-two-fold increase in system size, the 2.0 cutoff list grows by a factor of 30 against 435 for the baseline (Figure~\ref{fig:scaling}).

\begin{figure}[htbp]
\centering
\resizebox{\textwidth}{!}{\input{figures/scaling}}
\caption{Wall time against system size on exclusive Elja
\texttt{64cpu\_256mem} (Xeon Platinum 8358, one core; record
\texttt{elja-1788072}, Table~\ref{tab:scaling}). Left: cutoff cell list and the
index-ordered builder. Right: CHILL+ bond correlations and primitive rings. Grey
dash: $O(n)$ through the $n=1000$ 2.0 point. Coral is baseline commit
\texttt{1c8efa38}. Teal is engine v2.6.0.}
\label{fig:scaling}
\end{figure}

\begin{table}[htbp]
\caption{\label{tab:scaling}Single-core wall time in milliseconds on exclusive Elja \texttt{64cpu\_256mem}, record \texttt{elja-1788072}. ``base'' columns are baseline commit \texttt{1c8efa38}; the others are engine v2.6.0.
Cutoff list and index-ordered builder at \(r_c=3.5\,\mathring{\mathrm{A}}\); CHILL+ bond correlations (\texttt{getCorrelPlus}) on the cutoff list; primitive rings with \(D=6\) on the index-ordered list.}
\centering
\scriptsize
\resizebox{\textwidth}{!}{%
\begin{tabular}{rrrrrrrrr}
\(n\) & cutoff, base & cutoff & index, base & index & CHILL+, base & CHILL+ & rings, base & rings\\
\hline
1000 & 1.975 & 1.461 & 6.794 & 1.443 & 3.446 & 1.404 & 116.282 & 8.064\\
2000 & 5.650 & 2.781 & 22.015 & 2.738 & 7.274 & 2.865 & 406.376 & 19.870\\
4000 & 17.793 & 5.305 & 99.424 & 5.291 & 13.915 & 5.648 & 576.520 & 37.208\\
8000 & 61.986 & 11.030 & 387.840 & 10.796 & 27.182 & 11.401 & 965.459 & 78.731\\
16000 & 225.734 & 22.712 & 1516.346 & 22.345 & 59.033 & 23.363 & 3626.729 & 226.554\\
32000 & 859.556 & 43.663 & 5975.106 & 42.700 & 110.161 & 45.881 & 4486.953 & 528.897\\
\end{tabular}
}
\end{table}

The index-ordered builder supplies the hydrogen-bond network and hence the ring analysis; in the baseline it is the all-pairs loop.
We route it through the same cell list.
At 32000 atoms its cost is \(6.0\,\mathrm{s}\) before that change and \(43\,\mathrm{ms}\) after it.
The two builders together take \(6.8\,\mathrm{s}\) before and \(86\,\mathrm{ms}\) after.
Neighbour-list construction is no longer a leading term at any system size in the table.
The \(k\)-nearest path is the linkcell walk.
\subsection{Cage order and ring routes}
The baseline assembly claims peripheral rings as it accepts cages.
Its cage count therefore depends on the order in which equatorial candidates are visited.
On this frame it returns 1955 cages in the natural ring order, and 2.0 returns 2040 in the position-keyed order, with the same flags on every ring (record \texttt{bench-2026-10-31}).
Under 32 random relabellings of the atoms the baseline returns 1865 to 1891 cages, 22 distinct counts (rule \texttt{partition\_orders} of the rerun of Section~\ref{sec:s-rerun}, file \texttt{results/partition/base-cages-permuted.txt}).
Equatorial candidates are now visited in the position-keyed order of Section~\mref{sec:orgf57a77e}.
The Catch2 case ``DDC cage assembly is independent of ring enumeration order'' requires the same cage count and the same number of DDC rings for the natural, sorted and reverse-sorted ring lists of this frame.
The engine returns 2040 cages.
Walk enumeration generates every closed walk of length below \(D\) and keeps only the primitive ones; the shortest-path construction generates only shortest-path pairs (Algorithm~\mref{alg:rings}).
The Catch2 case ``ringNetwork matches the generate-then-filter route'' requires identical canonical ring sets from the two on 400-atom jittered cells of the same construction.
The cells have uniform jitter of full width 0.05, 0.30 and 0.60 lattice spacings, and \(D\) takes every value in 3--8.
\subsection{Strong scaling}
On the cells of Table~\ref{tab:scaling} the shortest-path construction takes \(8.064\) against \(116.282\,\mathrm{ms}\) at \(n=1000\) and \(528.897\) against \(4486.953\,\mathrm{ms}\) at \(n=32000\).
On pinned core 3 of the cloud VM the same construction takes \(3.909\) against \(89.498\,\mathrm{ms}\) at \(n=1000\) and \(125.308\) against \(3478.454\,\mathrm{ms}\) at \(n=32000\) (record \texttt{bench-2026-10-31}, files \texttt{tip-scaling.txt} and \texttt{base-scaling.txt}, one session).
Each source is independent, so the construction also threads, and so does everything in front of it.
The cutoff neighbour list bins the molecules in fractional coordinates of the recovered triclinic cell and builds each row on one thread.
The index-ordered list converts one row per thread.
Each source's rings are collected separately and flattened in source order, so the ring set is identical at every thread count.
For Figure~\ref{fig:strong} we keep a 65536-molecule cell fixed and time every host stage of the pipeline as the best of seven runs at each thread count.
Each thread runs on its own physical core (\texttt{OMP\_PLACES=cores}, \texttt{OMP\_PROC\_BIND=close}).
The machine is the 16-core, 32-thread AMD Ryzen 9 9950X workstation of record \texttt{terra-2026-09-02}, not reserved for the run.
The engine is commit \texttt{0a7fc546}, built with OpenMP and without vesin, so both neighbour lists are the in-tree cell list.
The pipeline takes \(906\,\mathrm{ms}\) on one thread and \(99\,\mathrm{ms}\) on sixteen, a speed-up of \(9.1\).
On the same two thread counts the rings alone take \(802\) and \(76\,\mathrm{ms}\) (\(10.6\times\)), and the Steinhardt kernel takes \(12.6\) and \(4.8\,\mathrm{ms}\).
The two neighbour lists take \(46\) and \(45\,\mathrm{ms}\) on one thread and \(11\) and \(8\,\mathrm{ms}\) on sixteen (\(4.3\times\) and \(5.8\times\)).
At sixteen threads the pipeline is \(43\,\mathrm{ms}\) above the ideal line; the rings account for \(26\,\mathrm{ms}\) of that and the two neighbour lists for \(13\,\mathrm{ms}\).

\begin{figure}[htbp]
\centering
\resizebox{\textwidth}{!}{\input{figures/strong}}
\caption{OpenMP strong scaling of the host ring pipeline on a
65536-molecule cell. The curves are the sum of every stage, the primitive rings
alone, and the two neighbour lists together. Best of seven runs per
point, one thread per core on a 16-core AMD Ryzen 9 9950X (record
\texttt{terra-2026-09-02}, engine \texttt{0a7fc546}). Grey
dash: ideal $1/T$ through the one-thread pipeline time.}
\label{fig:strong}
\end{figure}
\subsection{Incremental runs in detail}
\label{sec:s-incremental}
Across consecutive frames Proposition~P4 turns the per-frame construction into an exact incremental update.
Only the rings through a vertex of \(T\) are enumerated again, and at \(D=6\) the set \(T\) is the endpoints of the changed bonds.
On a jittered 8000-molecule cell with three perturbed sites per frame, the update matches a full rebuild on each of five frames.
The smallest update time over the five is \(1.477\,\mathrm{ms}\) against a smallest rebuild of \(28.839\,\mathrm{ms}\) (record \texttt{bench-2026-10-31}, median over three sessions).
On the last frame \(T\) holds 10 of the 8000 atoms, and the update refreshes 126 distance rows.
We update the eleven stored frames of the 4096-molecule mW cubic trajectory on the mutual four-nearest graph with a \(2\,\mathring{\mathrm{A}}\) candidate skin (Figure~\ref{fig:incremental}).
These are the ideal lattice at step 0 and ten later frames up to step 1000.
That list nominates the four nearest among the pairs within \(5.5\,\mathring{\mathrm{A}}\) at its last rebuild.
It rebuilds once a molecule has moved \(1\,\mathring{\mathrm{A}}\), so every pair now closer than \(3.5\,\mathring{\mathrm{A}}\) is a candidate.
The graph is then the exact mutual four-nearest graph whenever each molecule's fourth neighbour lies within \(3.5\,\mathring{\mathrm{A}}\).
The record marks a skin rebuild on each of the eleven frames, so every frame's graph is nominated afresh from the pairs within \(5.5\,\mathring{\mathrm{A}}\) of that frame.
We also search all pairs without a cutoff.
That search puts the fourth neighbour of every molecule within \(2.76\,\mathring{\mathrm{A}}\) on all eleven frames and returns the same mutual four-nearest graph, so the updated graphs are exact.
No bond of that graph changes after the first frame, so the edge diff certifies that no ring can change.
The update then returns the previous ring network without enumerating anything.
It takes \(0.170\,\mathrm{ms}\) against \(3.316\,\mathrm{ms}\) for a ring rebuild on the same skin-maintained graph, and its output equals that rebuild (record \texttt{bench-2026-10-31}, frames 2--11 means, median over three sessions; ring sets identical on every frame).
The incremental cage classification behaves the same way on the same frames.
The affiliation updater reclassifies zero rings after the first frame and answers in \(0.273\,\mathrm{ms}\) against \(38.376\,\mathrm{ms}\) for the batch predicate.
Its flags equal those of the batch computation, flag for flag, on every frame.
The Catch2 case ``AffiliationUpdater tracks synthetic topology churn exactly'' keeps it equal to the batch predicate under bond removal.
In that case the updater reclassifies only the locality closure of the cuts.
The test suite compares the updater with a full evaluation across the mW trajectory, under synthetic churn and with every neighbour row changed.
On mW cubic the affiliation flags coincide ring for ring with the greedy classification (every six-ring DDC-affiliated, none HC) and on an isolated hexagonal cage the HC flags match the greedy cage exactly.

\begin{figure}[htbp]
\centering
\resizebox{\textwidth}{!}{\input{figures/incremental}}
\caption{Ring and affiliation updates against a full rebuild on the eleven
stored mW cubic frames, mutual four-nearest graph, seams-core \coretag{}
on one pinned core of the cloud VM, record \texttt{bench-2026-10-31}
(frames 2--11 means, median over three sessions).
Rings take $0.170\,\mathrm{ms}$ against $3.316\,\mathrm{ms}$, and affiliation
$0.273\,\mathrm{ms}$ against $38.376\,\mathrm{ms}$. Outputs are identical
to the rebuild on every frame.}
\label{fig:incremental}
\end{figure}

In the figshare mW nucleation excerpt and Niu trajectory 3, bonds change on each transition from frame 1 through frame 11.
On frames 2--11 of the excerpt (4096 molecules, cutoff graph), 2845 bonds change per frame on average, and 1948 change on frames 2--11 of Niu trajectory 3 (2880 molecules).
Both updaters then take the batch path, and their output equals the rebuild on every frame.
In record \texttt{bench-2026-10-31} the nucleation frames take \(13.491\) against \(12.997\,\mathrm{ms}\) for rings and \(14.925\) against \(14.791\,\mathrm{ms}\) for affiliation on the pinned core (one session).
On eight threads they take \(2.523\) against \(2.280\,\mathrm{ms}\) and \(3.187\) against \(3.173\,\mathrm{ms}\).
On frames 2--11 of Niu trajectory 3 the ring update takes \(9.784\) against \(9.681\,\mathrm{ms}\) and the affiliation update \(8.635\) against \(8.557\,\mathrm{ms}\) on the pinned core, and \(1.773\) against \(1.860\,\mathrm{ms}\) and \(2.122\) against \(2.151\,\mathrm{ms}\) on eight threads.
On one core the two updates therefore run within \(3.8\%\) of the rebuild on both trajectories, where every tested transition takes the batch path.
In the churn sweep of the main text the two updates cost \(0.026\) to \(0.042\) and \(0.10\) to \(0.12\) of a rebuild with 3 to 17 changed bonds.
With no changed bond they return the stored result in \(0.166\) and \(0.271\,\mathrm{ms}\) (record \texttt{bench-2026-10-31}, file \texttt{churn-sweep-1t.txt}).
The reproducibility package runs the same driver as an asv suite under Snakemake (\texttt{asv/benchmarks/churn.py} and \texttt{workflow/asv.smk}, record \texttt{asv-churn-2026-10-31}).
Its eight \texttt{track\_} benchmarks read the columns of one driver run at each of twelve radii, and its setup fails on any radius where an update and the batch computation differ.
On pinned core 3 its changed-bond and changed-atom counts equal those of \texttt{bench-2026-10-31} at every radius, and its timings lie within \(2.5\%\) of them.
With 3 changed bonds the ring update takes \(0.343\) against \(13.048\,\mathrm{ms}\) and the affiliation update \(1.878\) against \(18.637\,\mathrm{ms}\).

Record \texttt{incremental-md-2026-10-31} times both updaters on mW trajectories run for the purpose (cutoff graph, rings up to six members, default switch fractions, pinned core 3, times summed over the frames after the first).
There are 168 of them, eleven frames each, saved every 1, 2, 3, 5, 10, 20, 50 or 100 steps of \(10\,\mathrm{fs}\) (record \texttt{md-strides-2026-10-31}).
They continue three runs each of ice Ic (4096 molecules) and ice Ih (4032) at 200 and \(250\,\mathrm{K}\) and of the liquid (4096) at 215, 250 and \(300\,\mathrm{K}\).
On ice at most 0.59\% of the molecules change their neighbour row per frame and neither updater switches.
The ring update takes \(1.899\) to \(5.726\,\mathrm{ms}\) against \(37.538\) to \(47.553\,\mathrm{ms}\) for a rebuild, and the affiliation update \(3.150\) to \(91.941\,\mathrm{ms}\) against \(414.835\) to \(527.047\,\mathrm{ms}\).
In the liquid 8.8 to 17.4\% of the molecules change their row at a stride of one step, and of two steps at \(215\,\mathrm{K}\), and the ring update takes \(57.471\) to \(112.764\,\mathrm{ms}\) against \(145.664\) to \(181.277\,\mathrm{ms}\).
From 23.4\% the changed atoms exceed a fifth of the frame, the ring updater switches on every frame, and it takes 1.3 to 6.1\% longer than the rebuild.
The affiliation updater switches on every liquid frame and takes between 2.1\% less and 2.8\% more than the batch pass.
On the 147 frames of the nucleation excerpt after the first, 67.15\% of the molecules change their row per frame and both updaters switch on every frame.
The ring update takes \(1997.386\) against \(1946.342\,\mathrm{ms}\), 2.6\% longer, and the affiliation update \(2454.199\) against \(2457.197\,\mathrm{ms}\).
Every output equals the rebuild.

Table~\ref{tab:carry} gives the command-line timings of Section~\mref{sec:incremental} (record \texttt{cli-carry-2026-10-31}).
Without the back-off of Section~\mref{sec:affupdate} the ratio on the two deposited trajectories lies between 0.916 and 1.046.
The seeded walk of the excerpt is then 9\% slower than the rebuild, since each of its frames pays for a diff and a batch pass.
The record holds both builds and the patch between them.

\begin{table}[htbp]
\caption{\label{tab:carry}\texttt{seams cages} over a frame range on pinned core 3 of the cloud machine, \texttt{OMP\_NUM\_THREADS=1}, record \texttt{cli-carry-2026-10-31}.
Milliseconds per frame, whole process, medians over five alternating sessions.
The ratio is rebuild over carry, so a value above 1 means the carried walk is faster.
The last column is the same ratio for the build without the back-off.
The carried and rebuilt outputs are identical in every case.}
\centering
\footnotesize
\begin{tabular}{llrrrrr}
trajectory & graph & frames & rebuild & carried & ratio & no back-off\\
\hline
mW cubic, 4096 & seeded & 11 & 105.08 & 16.47 & 6.382 & 6.402\\
mW cubic, 4096 & cutoff & 11 & 46.26 & 8.56 & 5.405 & 5.237\\
nucleation excerpt, 4096 & seeded & 148 & 47.21 & 48.13 & 0.981 & 0.916\\
nucleation excerpt, 4096 & cutoff & 148 & 32.87 & 33.39 & 0.985 & 0.964\\
Niu trajectory 3, 2880 & seeded & 148 & 30.02 & 28.92 & 1.038 & 1.046\\
Niu trajectory 3, 2880 & cutoff & 148 & 19.41 & 19.73 & 0.984 & 0.981\\
\end{tabular}
\end{table}

In the timed engine commits vesin \citep{vesin} supplies the cutoff pair list; \citet{goswami2026linkcell} supplies the \(k\)-nearest walk, and from v2.11.0 the cutoff pairs as well. sphericart, IRA/SOFI and nauty can be compiled in; they stay off the default classification path.
\subsection{Reproduction on a second architecture}
\label{sec:org2f64c1b}

The reproducibility package (\url{https://github.com/d-SEAMS/dseams2_repro}, \reproid) \citep{dseams2repro} runs every rule of this section, the label walks of Section~\ref{sec:applications} and the sweep of Table~\ref{tab:sota} from one workflow.
It clones the engine, its baseline commit and both front ends at the revisions of \texttt{config/ecosystem-lock.json}.
The engine commit of each record in Table~\ref{tab:records-full} is the one its conditions, machine or source-manifest file names.
For \texttt{elja-1788072} none of these names a commit.
Its strong-scaling ring column first exists at \texttt{bc901bed}, which \texttt{e4a7cf4f} merged to the main branch.
The record was committed to seams-core at \texttt{16a710c0}.
At every commit between \texttt{bc901bed} and \texttt{16a710c0}, both included, \texttt{src/} and \texttt{tests/} are the same as at v2.6.0.
Every one of these commits is an ancestor of \texttt{305fb3c6}, the base of \coretag{} three commits after v2.11.0 (\texttt{9f963b72}), in the seams-core history.
With the patch its lock pins \coretag, so its default rerun builds that release.
Its most recent exclusive run, record \texttt{elja-1798666}, is an AMD EPYC 7713 node (64 cores per socket, two sockets) at engine commit \texttt{e7856408}.
That node differs from the Xeon of the tables above in vendor, microarchitecture and memory system.
Table~\ref{tab:second-arch} gives the same rows there, one core per timing except the thread sweep of the last row.
On that node the baseline and the engine take \(8689.6\) and \(59.3\,\mathrm{ms}\) for the index-ordered builder at 32000 atoms.
The cell list takes \(862.8\) and \(61.0\,\mathrm{ms}\), and the hexagonal-cage search takes \(4926.9\) and \(68.7\,\mathrm{ms}\).
The ring-and-cage path on the cubic frame takes \(5255.8\) and \(93.0\,\mathrm{ms}\).
On the stored trajectory the incremental updaters are 20 and 8.6 times faster than the full rebuild of the same engine for rings and affiliation, with outputs equal to the rebuild on every frame.

\begin{table}[htbp]
\caption{\label{tab:second-arch}The exclusive-node rows of this section re-measured on one core of an AMD EPYC 7713 (Elja \texttt{128cpu\_256mem}, load \(0.2\) at start) by the reproducibility package (\url{https://github.com/d-SEAMS/dseams2_repro}), record \texttt{elja-1798666}.
The 2.0 column is engine commit \texttt{e7856408}.
The reference column is baseline commit \texttt{1c8efa38} in the first six rows, the full rebuild of the same engine in the two incremental rows, and one thread against 32 in the last row.
Milliseconds.}
\centering
\scriptsize
\begin{tabular}{lrrr}
stage & reference & 2.0 & factor\\
\hline
index-ordered neighbours, 32000 & 8689.6 & 59.3 & 147\\
cell-list neighbours, 32000 & 862.8 & 61.0 & 14\\
CHILL+, 32000 & 128.6 & 66.6 & 1.9\\
primitive rings, 32000 & 5111.5 & 752.9 & 6.8\\
hexagonal-cage search, mW 4096 & 4926.9 & 68.7 & 72\\
ring-and-cage path, mW 4096 & 5255.8 & 93.0 & 57\\
incremental rings, steady frame & 20.5 & 1.0 & 20\\
incremental affiliation, steady frame & 97.3 & 11.3 & 8.6\\
ring construction, 65536 atoms, 1 to 32 threads & 1893.9 & 119.2 & 16\\
\end{tabular}
\end{table}

The label tables do not depend on the machine at all.
The per-frame walks of Section~\ref{sec:applications} at engine commit \texttt{e7856408} on that node and at its parent commit \texttt{770accf9} on the workstation of record \texttt{terra-2026-09-01} are identical byte for byte, one line per frame.
The walks that section reads are record \texttt{walks-2026-10-31}, seams-core \coretag{} on the cloud VM.
On every frame of the six walks its CHILL+ columns equal those of \texttt{elja-1798666}, and its cage columns are those of the \coretag{} cage tests.
The EPYC record has no PTM rows, because OVITO did not load on that node (\texttt{libxkbcommon.so.0: cannot open shared object file}, line 10 of its \texttt{sota\_compare.txt}).
The sweep of Table~\ref{tab:sota} and Figure~\mref{fig:sota} of the main text ran with PydSEAMSlib 2.11.1 built on seams-core \coretag.
We also run the same installation on seeds 0--4.
There it returns every CHILL+ label of \texttt{elja-1798666}.
The PTM labels of \texttt{terra-2026-09-01} were scored at \texttt{rmsd\_cutoff=0.15}.
At \(0.1\), \(48\) of those \(90\) PTM cells differ, all at \(\sigma\ge 0.20\,\mathring{\mathrm{A}}\).
The completed rule keeps \(0.913\) of Ic and \(0.896\) of Ih at \(0.4\,\mathring{\mathrm{A}}\) (file \texttt{sota\_compare-seeds0-4.txt}).
Of its 273 cutoff, mutual and seeded cells, 22 differ from \texttt{elja-1798666}, all at \(\sigma\ge 0.3\,\mathring{\mathrm{A}}\) and each higher.
The PyPI build on v2.11.0 (\texttt{9f963b72}) returns all 273 (file \texttt{sota\_compare-seeds0-4-v2.11.0.txt}).
On seeds 5--9 of Table~\ref{tab:sota} that build keeps \(0.931\) of Ic and \(0.728\) of Ih under the completed rule at \(0.4\,\mathring{\mathrm{A}}\), against \(0.942\) and \(0.728\) at \coretag{} (record \texttt{sota-heldout-2026-10-31}, file \texttt{sota\_compare-seeds5-9-v2.11.0.txt}).
The completed labels of \texttt{elja-1798666}, \(0.994\) and \(0.990\), come from its engine commit \texttt{e7856408}.
That commit is the parent of \texttt{ba0634f5}, which replaced edge-sharing completion by the all-but-one rule.
\subsection{Steinhardt parameters across codes}

One cross-code check is already recorded.
We evaluate Steinhardt \(q_4\), \(q_6\) and \(q_8\) with the engine, freud and pyscal3 in two rules of the reproducibility package, \texttt{ql\_compare\_}\allowbreak\texttt{dseams} and \texttt{ql\_compare\_}\allowbreak\texttt{python}.
The test systems are an ideal fcc cell (108 atoms in a \(12\,\mathring{\mathrm{A}}\) box, cutoff \(3.2\,\mathring{\mathrm{A}}\)) and an ideal bcc cell (54 atoms, same box, cutoff \(4.0\,\mathring{\mathrm{A}}\)).
The three codes print the same six decimals.
They are \(0.190941\), \(0.574524\) and \(0.403915\) for fcc and \(0.509175\), \(0.628539\) and \(0.212762\) for bcc (record \texttt{elja-1798666}, files \texttt{ql-dseams.txt} and \texttt{ql-}\allowbreak\texttt{python.txt}).
\subsection{Head-to-head calls}
\label{sec:s-h2h}
Every entry of Table~\mref{tab:compare} of the main text runs on pinned core 3.
Record \texttt{h2h-2026-10-31} prints the version of every code in \texttt{raw/env.txt}.
Outside the RINGS rows each entry times the analysis call with the frame in memory, as the best of seven runs after one warm-up.
The d-SEAMS CHILL+ rows use the \(3.5\,\mathring{\mathrm{A}}\) cutoff list, the default of every front end (Table~\mref{tab:params}).
The CNA rows set the seeded cage label against the identify-diamond modifier of OVITO, which applies common-neighbour analysis to the second shell, and the PTM row sets the completed seeded cage rule against polyhedral template matching.
The RINGS rows give two times on each side, the whole process with its input and the ring stage alone.
The ring stage of RINGS is the Primitive Rings line of its own timing report, which excludes its bond search and its distance matrices.
The ring stage of d-SEAMS is its library ring call with the neighbour search, and its whole process is \texttt{seams fingerprint}.
GenIce2 calls cycless 0.7 for its rings and cages, through its ring formatter and the genice2-cage 2.5.2 cage formatter.
Its own \texttt{analice2} front end, version 2.2.13.3, stops with \texttt{AttributeError: 'AnalIce' object has no attribute 'replica\_vectors'} on both the mW and the sI input (files \texttt{raw/analice2\_crash.txt} and \texttt{raw/analice2\_crash\_si\_cage.txt}), so the GenIce2 rows call cycless directly.
The mW RINGS rows are n=3: three recorded repetitions and no further warmup, run at the same time with one repetition pinned to each of cores 3, 4 and 5.
The d-SEAMS repetitions used those same cores and finished before RINGS started.
The library ring call, cores 3, 4 and 5, took 6.735, 6.279 and 7.251~ms (median 6.735, range 6.279 to 7.251).
The RINGS Primitive Rings line on those cores took 52140, 52190 and 52360~ms (median 52190, range 52140 to 52360).
\texttt{seams fingerprint} took 69.500, 68.521 and 72.330~ms (median 69.500, range 68.521 to 72.330).
The RINGS wall times were 3694790.404, 3664208.870 and 3695483.511~ms (median 3694790.404, range 3664208.870 to 3695483.511).
The RINGS entry on CRN1 is unchanged: the best of seven whole processes after one warm-up.
OVITO's CHILL+ modifier and the d-SEAMS CHILL+\((r_c)\) both run at \(r_c=3.5\,\mathring{\mathrm{A}}\).
The agreement column counts the molecules that receive the same label, with clathrate, interfacial clathrate and water matched to OVITO's hydrate, interfacial hydrate and other.
The freud side computes the plain and the averaged parameter as two objects, each over the shared list, and freud stores coordinates in single precision.
On the mW frame the mutual four-nearest graph of d-SEAMS and the four-nearest rows of freud are the same graph.
The Niu frame is the final frame of the oxygen-only dump of Table~\ref{tab:public-cnt}, and the sI framework is the one of Table~\ref{tab:atlas}.
The d-SEAMS side of each row except the PTM row calls the engine through PydSEAMSlib.
The CHILL+\((r_c)\) row calls \texttt{neighListO}, \texttt{getCorrelPlus} and \texttt{getIceTypePlusNoPrint}, so that both CHILL+ columns use the \(3.5\,\mathring{\mathrm{A}}\) cutoff list.
The \(q_6\) and \(\bar q_6\) rows build the list once with \texttt{kNearestNeighbourList}, outside both timers, and time \texttt{steinhardtQl}.
The GenIce2 ring rows call \texttt{neighListO}, \texttt{neighbourListByIndex} and \texttt{ringNetwork}, and the sI cages add \texttt{findBySignature}.
The CNA rows call \texttt{kNearestNeighbourPair}, then \texttt{neighbourListByIndex} and \texttt{ringNetwork} on the mutual and the union list, and then \texttt{seededCageAffiliation}, the label of \texttt{seams cages -{}-graph seeded}.
The PTM row is \texttt{tests/bench\_paper.cpp} in mode \texttt{seeded-adj}, and Section~\ref{sec:cagestar} records why that time replaced the \(15.6\,\mathrm{ms}\) Python call.
The d-SEAMS whole process is this command line.
\begin{lstlisting}[language=bash,numbers=none]
seams fingerprint crn.lammpstrj --cutoff 3.5 --type 1
\end{lstlisting}
The d-SEAMS ring stage is the library call: the \(3.5\,\mathring{\mathrm{A}}\) neighbour list, then \texttt{ringNetwork} at depth 6.
The genice2-cage formatter run from the command line first replicates the sI unit cell and builds a hydrogen-bond graph.
We therefore run it once in the record, and we do not time that run. It returns the same \(16\) and \(48\) cages.
The other side is the public interface of each code.
\begin{itemize}
\item OVITO CHILL+ is \texttt{ChillPlusModifier(cutoff=3.5)} applied to a \texttt{DataCollection} of the \texttt{ovito} module.
\item OVITO PTM is \texttt{PolyhedralTemplateMatchingModifier} with \texttt{rmsd\_cutoff=0.1}, from the same module.
\item OVITO identify diamond is \texttt{IdentifyDiamondModifier}.
Its adaptive common-neighbour analysis is \texttt{CommonNeighborAnalysisModifier}.
\item freud is \texttt{freud.order.Steinhardt(l=6)} with \texttt{average=False} and with \texttt{average=True}, each over the mutual four-nearest list passed as a \texttt{freud.locality.}\allowbreak\texttt{NeighborList}.
\item For GenIce2 rings, \texttt{pairlist.pairs\_iter} at \(3.5\,\mathring{\mathrm{A}}\) fills a networkx graph, and \texttt{cycles\_iter(graph, 6, pos)} of \texttt{cycless.cycles} runs on it.
\item For GenIce2 cages, \texttt{polyhedra\_iter(rings, 14)} of \texttt{cycless.polyhed} runs on those rings.
      It keeps the cages of 12 and 14 five- and six-rings, as the genice2-cage formatter does.
\item RINGS runs \texttt{rings input} on the same coordinates with one species, periodic boundaries, primitive rings only, search depth 3 (rings of up to six members) and the O--O bond at \(3.5\,\mathring{\mathrm{A}}\).
\end{itemize}
Freud's own four-nearest list differs from the mutual list at 175 of the 2880 Niu molecules, which is why both codes read the mutual list in the \(q_6\) rows.
As a whole process the printed three-decimal medians give 3694790.404/69.500 = 53162.5, so RINGS takes that many times as long as \texttt{seams fingerprint} on the mW frame, about \(5.3\times10^{4}\), and 10000 times as long on CRN1.
The unrounded medians give 3694790.404052/69.499658 = 53162.71, and the ring-stage unrounded medians give 52190/6.73451 = 7749.64.
On the mW frame the Primitive Rings line is 52190.000/3694790.404 of the whole-process median.
RINGS builds a \(5\times5\times5\) superlattice for a periodic cell, and that distance matrix is the rest of the whole process.
The adaptive common-neighbour analysis of OVITO takes \(10.7\), \(9.72\) and \(1.50\,\mathrm{ms}\) on the mW, Niu and jittered frames (best of seven, \texttt{raw/cna.txt} of the record).
It has no diamond structure types and labels every molecule of the three frames other.
The completed cage rule takes \(7.104\,\mathrm{ms}\) and PTM takes \(16.312\,\mathrm{ms}\) on the jittered lattice, the medians of the session in Section~\ref{sec:cagestar}.
On the mutual four-nearest list, the other option of every front end, the Niu CHILL+ labels agree with OVITO at 2611 molecules.

\section{Applications}
\label{sec:applications}
A trajectory yields a time series of cage membership without visual inspection of the hydrogen-bond network \citep{goswami2020dseams}.
The original engine was written for that use.
It was applied to confined ices \citep{goswami2020nanoribbon,goswami2020hpice,suraj2021salt}, to heterogeneous nucleation on AgI \citep{zielke2015agi,prerna2019ice} and to the largest-cluster size of a seeding study under shear \citep{goswami2020seeding,goswami2021shearprl,goswami2021shearedreview}.
We apply the same cage assignment to a perfect cubic crystal of the monatomic water model \citep{molinero2009mw} and to a mixed liquid-like TIP4P/Ice frame \citep{abascal2005tip4pice}.
We also apply it to the public nucleation dumps of \citet{niu2019temperature} and to the published figshare mW excerpt, which is the \(n_{\max}\) record of \citet{goswami2020seeding}.
The last two systems are a Madrid-2019 ice/brine front and occupied GenIce methane hydrates.
\subsection{Cubic ice as a double-diamond cage field}
\label{sec:orgb4c8393}

Ice Ic is a network of double-diamond cages.
Every water is a DDC vertex, and the six-membered rings of the crystal belong to those cages rather than to hexagonal cages.
We apply that topological criterion to the 11-frame 4096-molecule mW cubic trajectory.
Across all frames every six-membered ring is a DDC member and no hexagonal cages appear, the ice-Ic signature of \citet{goswami2020dseams}.

Figure~\mref{fig:cages} of the main text shows the two topological units, which we render with solvis \citep{goswami2024solvis}.
The first cubic mW frame assembles 2040 double-diamond cages.
Part (a) is one of them, with fourteen vertices across seven six-rings.
Part (b) is the twelve-atom hexagonal cage of ice Ih.
Its five six-rings are the two basal hexagons and the connecting prism faces.


Figure~\mref{fig:ice-lattices} of the main text shows the \(2\times 2\times 2\) ice Ic and ice Ih oxygen lattices those cages tile.
They come from the same generators as the classifier sweep of Table~\ref{tab:sota}, and we render them with solvis \citep{goswami2024solvis}.

\subsection{Cage membership and CHILL+}
\label{sec:org4411d8e}

CHILL+ classifies the star of bonds around each molecule, whereas a cage label needs a complete cage.
CHILL+\((r_c)\) reads the \(3.5\,\mathring{\mathrm{A}}\) cutoff list and CHILL+(4NN) the mutual four-nearest graph.
The two assignments agree on the 4096-molecule cubic mW frame.
Every molecule belongs to a double-diamond cage, and CHILL+ reports 4096 cubic ices (\(100\%\)), the four-staggered signature of a DDC vertex.
Neither assignment finds a hexagonal cage.
On the mixed TIP4P example included with the engine (250 waters) they disagree.
CHILL+\((r_c)\) marks 12 molecules as interfacial clathrate (\(4.8\%\)), three eclipsed bonds of four, and 238 as water (\(95.2\%\)), while the cage search finds no complete cage.
On the same frame CHILL+(4NN), the variant of \texttt{seams chill-plus -{}-graph knn}, marks 1 molecule as clathrate, 7 as interfacial clathrate and 242 as water; on the cubic frame it also reports 4096 cubic ices.
We collect those counts in Table~\ref{tab:icescores}.

\begin{table}[htbp]
\caption{\label{tab:icescores}Cage membership and CHILL+ on the cubic mW frame and on the mixed TIP4P example.
Ice Ic is DDC membership; ice Ih is HC membership.
CHILL+ rows are CHILL+\((r_c)\) on the \(3.5\,\mathring{\mathrm{A}}\) cutoff list, as in the Catch2 case of \texttt{tests/test\_bop.cpp}.}
\centering
\begin{tabular}{lrr}
score & cubic mW (4096) & mixed TIP4P (250)\\
\hline
cage Ic (DDC) & 4096 (\(100\%\)) & 0\\
cage Ih (HC) & 0 & 0\\
CHILL+ cubic & 4096 (\(100\%\)) & 0\\
CHILL+ hexagonal & 0 & 0\\
CHILL+ interClathrate & 0 & 12 (\(4.8\%\))\\
CHILL+ water & 0 & 238 (\(95.2\%\))\\
CHILL+ bulk ice (Ic+Ih) & 4096 (\(100\%\)) & 0\\
\end{tabular}
\end{table}

Those twelve interfacial labels remain raw \(3E\) stars.
The hydrate pipeline of \citet{nguyen2015identification} clusters \(3E{+}4E\) molecules at \(3.5\,\mathring{\mathrm{A}}\) and discards liquid when those clusters are small and scattered.
It takes the largest such cluster as the hydrate count, and that step does not reclassify these twelve.
The original CHILL rule labels every oxygen as water.
The same four-neighbour limit is why CHILL and cage membership disagree at an adsorbed layer.
CHILL calls that layer liquid, while its molecules already sit in complete hexagonal or double-diamond cages.
The first two water layers on defect AgI were reported as HC-dominated \citep{prerna2019ice}, and the adsorbed layer already sits in complete cages while CHILL labels it liquid \citep{goswami2020dseams}.

A growing nucleus shows the same disagreement on every frame.
In Figure~\mref{fig:app-nucleation} of the main text we follow the 148 stored frames of the 4096-molecule mW nucleation deposit of the original release \citep{goswami2020dseams}.
Each frame carries both scores on the same cutoff graph and the seeded assignment.
On the last frame CHILL+ labels 497 molecules cubic or hexagonal and places 454 of them in one cluster; the seeded cage assignment labels 995 and places 905 in one cluster (1004 and 900 on the cutoff graph).
Of the 995 seeded molecules 363 carry only a DDC label, 255 only an HC label and 377 both, so 740 lie on a DDC and 632 on an HC (record \texttt{walks-2026-10-31}, file \texttt{walks/figshare-nucleation.txt}).
CHILL+ marks a further 397 molecules interfacial and 319 interfacial clathrate on that frame, and with the interfacial label its ice count is 894.
Across the deposit the largest CHILL+ cluster is \(0.24\) to \(0.50\) of the largest seeded cage cluster.

\subsection{Homogeneous nucleation: the largest ice cluster as a time series}
\label{sec:org85ca619}

The observable of a seeding or committor study is \(n_{\max}(t)\), the largest connected cluster of ice-labelled molecules, followed frame by frame.
Whether it grows or melts decides the fate of a nucleus.
Where it crosses a reference size fixes the critical nucleus.
That number depends on the label.
We ask one question of every public nucleation trajectory.
Does the cage assignment, on either graph, recover the published growth or melting of the nucleus, and what does CHILL+ report on the same frames?
\(n_{\max}\) from cage membership is the cluster size \citet{goswami2020seeding} measured with d-SEAMS in a seeding study of ice nucleation under shear and the observable the shear-nucleation papers take from the engine \citep{goswami2021shearprl,goswami2021shearedreview}.
Seeding itself, as a route to homogeneous rates, is that of \citet{sanz2013seeding}, \citet{espinosa2014rates} and \citet{espinosa2016seeding}.

The inputs are every public dump of \citet{niu2019temperature} (Materials Cloud record 2020.0005/v1; TIP4P/Ice, LAMMPS \citep{thompson2022lammps}) and the figshare mW excerpt of the original release (11448702.v1).
Two of the five dumps are the critical-nucleus committor pair of that paper.
In one trajectory the nucleus grows, and in the other it melts.
Their \(N_c=314\pm 20\) at \(230\,\mathrm{K}\) is a free-energy-surface size on their own order parameter.
We use it as a reference line rather than as the occupancy of any frame.
We score every stored frame with CHILL+ and cage membership on the cutoff graph, and with the seeded assignment.
Figure~\mref{fig:app-public} of the main text shows \(n_{\max}\) for four of the dumps, Table~\ref{tab:public-cnt} the final frame of all six.

\begin{table}[htbp]
\caption{\label{tab:public-cnt}Final-frame \(n_{\max}\) on every public nucleation dump of \citet{niu2019temperature} and the figshare mW excerpt, by cage membership on the seeded graph, by cage membership on the cutoff graph, and by CHILL+ (cubic plus hexagonal).
The 1600-water extract stores only part of the cell on most frames.}
\centering
\scriptsize
\begin{tabular}{llccc}
dump & waters & seeded & cutoff & CHILL+\\
\hline
2880 oxygen-only nucleation & 2880 & 2009 & 2005 & 1369\\
2880 four-site nucleation & 2880 & 2870 & 2866 & 2774\\
1600-water extract & 1573 & 1567 & 1566 & 1422\\
critical nucleus, growing & 2880 & 660 & 647 & 295\\
critical nucleus, melting & 2880 & 18 & 30 & 1\\
figshare mW nucleation excerpt & 4096 & 905 & 900 & 454\\
\end{tabular}
\end{table}

The cage assignment recovers the published fate of every nucleus.
The growing member of the committor pair starts at \(n_{\max}=371\) and ends at 660; the melting member starts at 384 and ends at 18.
Both oxygen-only nucleation runs grow to a cell-spanning cluster, the four-site dump freezes through \(n_{\max}=2870\) of \(2880\) waters, and the 1600-water extract ends at \(n_{\max}=1567\) of \(1573\).
On the 2880-water oxygen-only path the cage cluster first exceeds \(N_c\) at frame 32 on the cutoff graph (\(n_{\max}=352\)) and at frame 34 on the seeded graph (\(n_{\max}=332\)).
On the two remaining nucleation runs both graphs cross on the same frame.

The two graphs agree on the fate of every nucleus.
At the end of five of the six dumps they differ by a few percent (Table~\ref{tab:public-cnt}); on the melting committor member they end at 18 and 30.
We also compare them over every frame of the walks of record \texttt{walks-2026-10-31}.
They differ by up to 249 molecules on the oxygen-only run (frame 126: 188 against 437) and 80 on the mW excerpt (frame 130).
The sweep of Table~\ref{tab:sota} predicts this effect.
A thermally displaced first-shell neighbour leaves the cutoff graph and takes its cage with it, while the four-nearest graph keeps the edge.

CHILL+ reports a different nucleus on the same frames.
We compute the ratio of the largest CHILL+ cubic-plus-hexagonal cluster to the seeded cage cluster on every frame (\texttt{scripts/cluster\_ratio.py} on the walks of record \texttt{walks-2026-10-31}).
The median ratio is \(0.39\) on the mW excerpt (range \(0.24\) to \(0.50\)), \(0.45\) on the growing committor member and \(0.35\) on the oxygen-only run.
On the committor pair CHILL+ places the two starting nuclei at 154 and 151 molecules, where the seeded cage cluster has 371 and 384.
The nucleus size is therefore a property of the coordinate.
On the same frames the two clusters differ by a factor of two to three.
A critical size or committor stated in one coordinate does not carry over to the other.

For that pair PydSEAMSlib writes the per-frame columns \(n_{\mathrm{ice}}\), \(n_{\max}\), \(n_{\mathrm{Ic}}\), \(n_{\mathrm{Ih}}\), \(n_{\mathrm{mixed}}\), cubicity and the CHILL+ counts.
We plot the same columns from the walks in Figure~\ref{fig:app-msm} (\texttt{msm\_features.csv} of record \texttt{walks-2026-10-31}).
On the growing member the seeded \(n_{\max}\) starts at \(371\) and ends at \(660\).
On the melting member it starts at \(384\) and ends at \(18\).
The CHILL+ \(n_{\max}\) starts at \(154\) and ends at \(295\) on the growing member.
On the melting member it starts at \(151\) and ends at \(1\).

\begin{figure}[htbp]
\centering
\input{figures/app_msm}
\caption{Per-frame series of the Niu critical-nucleus pair.
Left: growing member.
Right: melting member.
Solid: seeded cage $n_{\max}$.
Dashed: CHILL+ cubic-plus-hexagonal $n_{\max}$.
Points are every tenth stored frame of the Materials Cloud dump, from the first.}
\label{fig:app-msm}
\end{figure}

Inside the integrator the same labels decide when a seeding run stops.
Figure~\mref{fig:seed-xyz} of the main text shows that insertion.
An ice Ic sphere of \(80\) molecules is cut into a random packing of \(256\) mW sites, and liquid molecules closer than \(2.6\,\mathring{\mathrm{A}}\) are removed.
This is the same construction as the production campaign at \(4096\) molecules.
Teal bonds are the seed.
Gray spheres are the remaining liquid.


The same frames (cages, ice I cells, seed, cubic mW) open in a chemiscope widget (\texttt{viewer/index.html}).
The page loads \texttt{dseams.chemiscope.json} and colours the map by a static kind label (liquid, cubic/DDC, hexagonal/HC).
The classify path is a C ABI (\texttt{seams\_chill\_plus}) so the same labels can run under emscripten when a WASM build is present.

A spherical ice Ic seed of \(N\) molecules sits in a liquid of \(4096\) mW molecules at \(205\), \(215\) and \(225\,\mathrm{K}\).
Each temperature has one liquid, a random packing melted for \(100\,\mathrm{ps}\) at a \(5\,\mathrm{fs}\) step and equilibrated for \(300\,\mathrm{ps}\) of NPT at \(10\,\mathrm{fs}\).
For each of six sizes (\(50\) to \(700\)) we cut one seed of exactly \(N\) sites from the perfect lattice and centre it in the box in a fixed orientation.
We then remove every liquid molecule within \(2.6\,\mathring{\mathrm{A}}\) of the seed, which leaves \(n_{\mathrm{tot}}=N+n_{\mathrm{liq}}\) molecules (Table~\ref{tab:ntot}).
The twelve replicas of a \((T,N)\) cell therefore share the same starting coordinates.
Each replica draws its own Gaussian velocities at \(T\) and starts at the cut configuration, with no minimisation or restrained relaxation.
We run it for up to \(5\,\mathrm{ns}\) in LAMMPS with a \(10\,\mathrm{fs}\) step and \texttt{fix npt} at \(1\,\mathrm{bar}\), with temperature and isotropic pressure damping of \(1\) and \(10\,\mathrm{ps}\).
The grown fraction of a cell is therefore a velocity-shooting estimate for one configuration.
The PLUMED action \texttt{DSEAMS\_CAGES} (\texttt{CUTOFF=3.5}, \texttt{CANDIDATE=5.5}, \texttt{K=4}, \texttt{COMPLETE}) logs the seeded and completed labels and the CHILL+\((r_c)\) cluster every \(5\,\mathrm{ps}\).
With \texttt{COMPLETE} a molecule can carry an HC or DDC label through Eq.~\ref{eq:completion} without lying on an accepted cage.
Every production count below is therefore a count of seeded and completed labels.
Here \(n_{\max}\) is the largest connected cluster of labelled molecules, and the figure legends call it the cage cluster.
\texttt{COMMITTOR} stops a run when labelled \(n_{\max}\) falls to \(20\) or reaches \(2N\).
On the first logged stride the box carries \(0.72N\) to \(1.15N\) Ic labels, and \(103\) to \(142\) molecules carry an Ih label at \(205\,\mathrm{K}\), \(12\) to \(37\) at \(215\) and \(225\,\mathrm{K}\) (\texttt{cubicity.txt}).
We report the grown fraction of that campaign in Figure~\ref{fig:app-committor} and Table~\ref{tab:committor} (record \texttt{elja-production}, engine v2.9.2, file \texttt{committor.txt}).

\begin{figure}[htbp]
\centering
\input{figures/app_committor}
\caption{Grown fraction of a spherical ice Ic seed in $4096$ mW
molecules as a function of inserted seed size $N$, twelve replicas
per point. A replica grows when the labelled cluster reaches $2N$
and melts when it falls to $20$, and the run stops there. Filled squares (legend ``cage''):
seeded and completed $n_{\max}$. Open circles: CHILL+ cubic-plus-hexagonal $n_{\max}$
against the same thresholds, read at the stopping frame.}
\label{fig:app-committor}
\end{figure}

\(N\) is the number of ice-Ic sites inserted, and \(p\) is the fraction of replicas whose labelled cluster reaches \(2N\) within the run.
\(p\) rises above one half between \(100\) and \(200\) molecules at \(205\,\mathrm{K}\) (\(0/12\) to \(11/12\) grown), between \(200\) and \(300\) at \(215\,\mathrm{K}\) (\(0/12\) to \(11/12\)), and between \(300\) and \(450\) at \(225\,\mathrm{K}\) (\(1/12\) to \(8/12\)).
The 95\% Wilson intervals for twelve replicas are \(0\)--\(0.24\) for \(0/12\), \(0.015\)--\(0.35\) for \(1/12\), \(0.39\)--\(0.86\) for \(8/12\), \(0.65\)--\(0.99\) for \(11/12\) and \(0.76\)--\(1\) for \(12/12\).
Twelve of twelve seeds of \(700\) molecules grow at every temperature.
The stopping rule censors the CHILL+ columns.
A run stops when the labelled cluster enters a basin.
On every run where that cluster reached \(2N\), the CHILL+ cluster at that frame lay between the two thresholds (the open fraction of Table~\ref{tab:committor}), as the smaller CHILL+ nucleus of Figure~\mref{fig:app-public} implies.
The slopes in Table~\ref{tab:committor} are positive in every cell with \(p>1/2\).

\begin{table}[htbp]
\caption{\label{tab:committor}Grown fraction \(p\) and open fraction of the seeding campaign, twelve replicas per cell.
A run stops when the seeded and completed \(n_{\max}\) reaches \(2N\) (grown) or falls to \(20\) (melted); the CHILL+ columns apply the same thresholds to the CHILL+ cluster at the stopping frame.
The slope is the mean change of the seeded and completed \(n_{\max}\) over the recorded window, in molecules\,ps\(^{-1}\).}
\centering
\scriptsize
\resizebox{\textwidth}{!}{%
\begin{tabular}{rrrrrrrr}
T (K) & N & runs & \(p\) cage & \(p\) CHILL+ & open cage & open CHILL+ & slope cage\\
\hline
205 & 50 & 12 & 0.083 & 0.000 & 0.000 & 0.000 & -0.70\\
205 & 100 & 12 & 0.000 & 0.000 & 0.000 & 0.000 & -2.32\\
205 & 200 & 12 & 0.917 & 0.000 & 0.000 & 0.917 & 0.63\\
205 & 300 & 12 & 1.000 & 0.000 & 0.000 & 1.000 & 1.26\\
205 & 450 & 12 & 1.000 & 0.000 & 0.000 & 1.000 & 2.65\\
205 & 700 & 12 & 1.000 & 0.000 & 0.000 & 1.000 & 4.02\\
215 & 50 & 12 & 0.000 & 0.000 & 0.000 & 0.000 & -3.76\\
215 & 100 & 12 & 0.000 & 0.000 & 0.000 & 0.000 & -11.97\\
215 & 200 & 12 & 0.000 & 0.000 & 0.000 & 0.000 & -5.94\\
215 & 300 & 12 & 0.917 & 0.000 & 0.000 & 0.917 & 1.29\\
215 & 450 & 12 & 1.000 & 0.000 & 0.000 & 1.000 & 2.83\\
215 & 700 & 12 & 1.000 & 0.000 & 0.000 & 1.000 & 5.74\\
225 & 50 & 12 & 0.000 & 0.000 & 0.000 & 0.000 & -4.28\\
225 & 100 & 12 & 0.000 & 0.000 & 0.000 & 0.000 & -11.87\\
225 & 200 & 12 & 0.000 & 0.000 & 0.000 & 0.000 & -10.51\\
225 & 300 & 12 & 0.083 & 0.000 & 0.000 & 0.083 & -2.51\\
225 & 450 & 12 & 0.667 & 0.000 & 0.000 & 0.667 & 0.60\\
225 & 700 & 12 & 1.000 & 0.000 & 0.000 & 1.000 & 4.23\\
\end{tabular}
}
\end{table}

Table~\ref{tab:ntot} gives \(n_{\mathrm{tot}}\) from rule \texttt{seed\_counts} of the rerun of Section~\ref{sec:s-rerun}, which repeats the packing, the liquid stage and the cut of every cell.
The packing and the seeds are the same in every run.
The liquid is not, because LAMMPS runs the \(50\,000\)-step liquid stage on four OpenMP threads.
We therefore quote the counts of that one run with its record, next to the range over three runs of the rule on the same machine with the same inputs, seeds and thread count.
The range is 2 to 21 molecules, at most \(0.52\%\) of \(n_{\mathrm{tot}}\), and no cell gives the same count in all three runs.
The liquids of the campaign itself were not archived, so its totals cannot be regenerated.

\begin{table}[htbp]
\caption{\label{tab:ntot}Molecules after the cut, \(n_{\mathrm{tot}}=N+n_{\mathrm{liq}}\), in each \((T,N)\) cell of the seeding campaign.
\(n_{\mathrm{liq}}\) and \(n_{\mathrm{tot}}\) are one run of rule \texttt{seed\_counts} (rerun of Section~\ref{sec:s-rerun}, file \texttt{results/production/seed\_counts.txt}).
The last column is the range of \(n_{\mathrm{tot}}\) over three runs of the rule on that machine, with the same inputs and seeds and four LAMMPS threads.}
\centering
\scriptsize
\begin{tabular}{rrrrr}
T (K) & N & \(n_{\mathrm{liq}}\) & \(n_{\mathrm{tot}}\) & \(n_{\mathrm{tot}}\), three runs\\
\hline
205 & 50 & 4018 & 4068 & 4068--4081\\
205 & 100 & 3955 & 4055 & 4055--4061\\
205 & 200 & 3830 & 4030 & 4030--4042\\
205 & 300 & 3722 & 4022 & 4017--4022\\
205 & 450 & 3538 & 3988 & 3988--3991\\
205 & 700 & 3237 & 3937 & 3937--3947\\
215 & 50 & 4018 & 4068 & 4068--4074\\
215 & 100 & 3956 & 4056 & 4054--4056\\
215 & 200 & 3838 & 4038 & 4026--4038\\
215 & 300 & 3721 & 4021 & 4012--4021\\
215 & 450 & 3541 & 3991 & 3985--3991\\
215 & 700 & 3229 & 3929 & 3927--3942\\
225 & 50 & 4023 & 4073 & 4068--4073\\
225 & 100 & 3951 & 4051 & 4049--4058\\
225 & 200 & 3830 & 4030 & 4028--4047\\
225 & 300 & 3714 & 4014 & 4006--4027\\
225 & 450 & 3524 & 3974 & 3974--3982\\
225 & 700 & 3232 & 3932 & 3926--3935\\
\end{tabular}
\end{table}

The same ICE files split the labelled molecules into Ic (DDC label only), Ih (HC label only) and mixed (both labels).
Figure~\ref{fig:app-cubicity} plots the last-frame mean of \(n_{\mathrm{Ic}}\), \(n_{\mathrm{Ih}}\) and \(n_{\mathrm{mixed}}\) against \(N\).
We define cubicity here as the DDC-containing label fraction \((n_{\mathrm{Ic}}+n_{\mathrm{mixed}})/n_{\mathrm{ice}}\).
This molecule count puts every mixed molecule on the cubic side, and it is not the stacking-sequence cubicity of a layer analysis.
Cubicity is \(0.73\) to \(0.89\) whenever the grown fraction exceeds one half, and \(0.14\) to \(0.30\) on cells with \(p=0\).
The last-frame means pool the grown and melted replicas of a cell.

\begin{figure}[htbp]
\centering
\input{figures/app_cubicity}
\caption{Last-frame split of the seeded and completed labels of the seeding campaign, twelve replicas per point, grown and melted replicas pooled.
Filled squares: $n_{\mathrm{Ic}}$.
Triangles: $n_{\mathrm{Ih}}$.
Open circles: $n_{\mathrm{mixed}}$.
The same ICE files as Figure~\ref{fig:app-committor}.}
\label{fig:app-cubicity}
\end{figure}

\begin{table}[htbp]
\caption{\label{tab:cubicity}Last-frame means of the seeded and completed label counts \(n_{\mathrm{Ic}}\) (DDC only), \(n_{\mathrm{Ih}}\) (HC only), \(n_{\mathrm{mixed}}\) (both) and \(n_{\mathrm{ice}}\), and the DDC-containing fraction \((n_{\mathrm{Ic}}+n_{\mathrm{mixed}})/n_{\mathrm{ice}}\) (cubicity), on the seeding campaign, twelve replicas per cell, grown and melted replicas pooled.}
\centering
\scriptsize
\resizebox{\textwidth}{!}{%
\begin{tabular}{rrrrrrrr}
T (K) & N & n & \(n_{\mathrm{Ic}}\) & \(n_{\mathrm{Ih}}\) & \(n_{\mathrm{mixed}}\) & \(n_{\mathrm{ice}}\) & cubicity\\
\hline
205 & 50 & 12 & 20.8 & 54.9 & 0.9 & 76.6 & 0.283\\
205 & 100 & 12 & 21.2 & 49.0 & 0.0 & 70.2 & 0.302\\
205 & 200 & 12 & 222.2 & 126.6 & 124.0 & 472.8 & 0.732\\
205 & 300 & 12 & 397.0 & 154.3 & 114.8 & 666.1 & 0.768\\
205 & 450 & 12 & 638.1 & 208.6 & 141.3 & 988.0 & 0.789\\
205 & 700 & 12 & 936.1 & 226.8 & 282.3 & 1445.2 & 0.843\\
215 & 50 & 12 & 7.6 & 24.2 & 0.0 & 31.8 & 0.238\\
215 & 100 & 12 & 5.1 & 20.9 & 0.0 & 26.0 & 0.196\\
215 & 200 & 12 & 5.8 & 31.3 & 0.0 & 37.1 & 0.155\\
215 & 300 & 12 & 305.1 & 135.6 & 155.7 & 596.3 & 0.773\\
215 & 450 & 12 & 517.2 & 165.2 & 266.8 & 949.3 & 0.826\\
215 & 700 & 12 & 1050.1 & 161.9 & 227.1 & 1439.1 & 0.887\\
225 & 50 & 12 & 3.6 & 22.1 & 0.0 & 25.7 & 0.140\\
225 & 100 & 12 & 6.4 & 17.1 & 0.0 & 23.5 & 0.273\\
225 & 200 & 12 & 4.7 & 19.2 & 0.0 & 23.8 & 0.196\\
225 & 300 & 12 & 36.1 & 19.1 & 16.7 & 71.8 & 0.734\\
225 & 450 & 12 & 288.7 & 152.9 & 197.3 & 638.9 & 0.761\\
225 & 700 & 12 & 934.8 & 206.0 & 291.5 & 1432.2 & 0.856\\
\end{tabular}
}
\end{table}

The production workflow also collects those PRINT rows in \texttt{ml\_\allowbreak labels.csv}, one row per replica and stride.
Its columns are temperature, seed size, replica, time, \(n_{\mathrm{Ic}}\), \(n_{\mathrm{Ih}}\), \(n_{\mathrm{mixed}}\), \(n_{\max}\), CHILL+ ice and CHILL+ \(n_{\max}\).
Per-molecule states (water, Ic, Ih, mixed) come from PydSEAMSlib on the dump.

The workflow rescores any dump placed under \texttt{data/shear}, such as the shear trajectories of \citet{goswami2021shearprl}, with the same cage walk as Figure~\mref{fig:app-public}.
It scores water after dropping every non-water type, so a substrate or protein atom never enters the neighbour graph.
\subsection{Ice/brine front}
\label{sec:org8c68582}

The scaled-charge Madrid-2019 ions of Zeron, Abascal and Vega \citep{zeron2019madrid} with TIP4P/2005 water reproduce the freezing-point depression of NaCl solutions \citep{lamas2022freezing}, and their GROMACS topologies are deposited \citep{sedano2026madridtopology}. We split a GenIce ice Ih block of \(1536\) TIP4P/2005 molecules (\(4\times4\times6\) cells) at half height along \(z\).
Then \(0\), \(10\), \(20\) or \(40\) NaCl pairs take the oxygen sites of twice as many waters drawn at random from the upper half, with a different draw per replica.
At \(40\) pairs \(1456\) waters remain, which is \(1.52\,\mathrm{mol\,kg^{-1}}\) over the box.
We hold the lower half fixed, melt the upper half for \(100\,\mathrm{ps}\) of NVT at \(400\,\mathrm{K}\) and quench it for \(100\,\mathrm{ps}\) of NVT at the run temperature.
The whole box then gets fresh Gaussian velocities.
The front runs for \(5\,\mathrm{ns}\) at \(230\), \(240\) and \(245\,\mathrm{K}\), below the \(250\,\mathrm{K}\) melting point of the water model \citep{lamas2022freezing}.
It runs under \texttt{fix npt} at \(1\,\mathrm{bar}\) with the box free along \(z\) only (temperature and pressure damping \(0.2\) and \(2\,\mathrm{ps}\)).
Every stage uses a \(2\,\mathrm{fs}\) step with SHAKE, an \(8.5\,\mathring{\mathrm{A}}\) cutoff for the Lennard-Jones and real-space Coulomb terms with no tail correction, and PPPM at a relative accuracy of \(10^{-5}\).
The PLUMED action (\texttt{ION\_CUTOFF=3.5}, \texttt{COMPLETE}) logs the seeded and completed label count, the largest labelled cluster and the three ion classes every picosecond.
One \(3.5\,\mathring{\mathrm{A}}\) shell radius applies to both ion species.
The ions enter the action as a second atom list and never enter the graph.
The Python package applies the same exclusion and reads GROMACS frames through ASE with the ions kept as their atomic numbers.
The LAMMPS input carries the Madrid-2019 parameters in kcal/mol next to their kJ/mol source values.
Table~\ref{tab:brine} and Figure~\ref{fig:app-brine} report that two-replica demonstration (record \texttt{elja-production}, engine v2.9.2, file \texttt{brine\_summary.txt}).
It has twenty-four runs, two at each of four pair counts and three temperatures, and each run has \(5001\) strides of \(1\,\mathrm{ps}\).
Figure~\mref{fig:app-brine-box} of the main text is the starting box of one 40-pair replica.


\begin{figure}[htbp]
\centering
\input{figures/app_brine}
\caption{Labelled $n_{\max}$ and liquid ions.
Last-stride means, two replicas.
Left: largest cluster of seeded and completed labels.
Right: ions whose shell is liquid.
Filled squares: $230\,\mathrm{K}$.
Filled circles: $240\,\mathrm{K}$.
Open circles: $245\,\mathrm{K}$.}
\label{fig:app-brine}
\end{figure}

Table~\ref{tab:brine} reports last-stride means over two replicas.
With no salt the labelled cluster grows at \(230\), \(240\) and \(245\,\mathrm{K}\).
It starts at \(930\), \(939.5\) and \(927\) molecules and ends at \(1203.5\), \(1275\) and \(1319.5\).
Over the same \(5\,\mathrm{ns}\) the mean \(n_{\max}\) changes by \(+229\), \(+176.5\) and \(+30\) molecules with ten pairs, by \(+86.5\), \(-33.5\) and \(-74.5\) with twenty, and by \(-6.5\), \(-32.5\) and \(-135.5\) with forty, at \(230\), \(240\) and \(245\,\mathrm{K}\).
No run ends with more than two ions in an all-ice shell.
The other ions sit at the front or in the liquid (\(8.0\), \(27.5\) and \(59.5\) liquid ions at \(230\,\mathrm{K}\) for \(10\), \(20\) and \(40\) pairs; \(15.0\), \(36.5\) and \(72.5\) at \(245\,\mathrm{K}\)).
On ions substituted into ideal ice lattices, Table~\ref{tab:ions} gives a CHILL+ first-shell ice fraction of \(0.000\), so CHILL+ alone cannot place an ion in ice.

\begin{table}[htbp]
\caption{\label{tab:brine}Ice/brine front after \(5\,\mathrm{ns}\), means of two replicas. \(n_{\max}(0)\) and \(n_{\max}(5\,\mathrm{ns})\) are the largest cluster of seeded and completed labels; the slope is \((n_{\max}(5\,\mathrm{ns})-n_{\max}(0))/5000\,\mathrm{ps}\); \(n_{\mathrm{ion,ice}}\), \(n_{\mathrm{ion,front}}\) and \(n_{\mathrm{ion,liq}}\) are the last-stride ion classes.}
\centering
\scriptsize
\resizebox{\textwidth}{!}{%
\begin{tabular}{rrrrrrrrr}
T (K) & pairs & n & \(n_{\max}(0)\) & \(n_{\max}(5\,\mathrm{ns})\) & slope & \(n_{\mathrm{ion,ice}}\) & \(n_{\mathrm{ion,front}}\) & \(n_{\mathrm{ion,liq}}\)\\
\hline
230 & 0 & 2 & 930.0 & 1203.5 & 0.055 & 0.0 & 0.0 & 0.0\\
230 & 10 & 2 & 899.0 & 1128.0 & 0.046 & 1.5 & 10.5 & 8.0\\
230 & 20 & 2 & 898.0 & 984.5 & 0.017 & 0.5 & 12.0 & 27.5\\
230 & 40 & 2 & 893.0 & 886.5 & -0.001 & 0.5 & 20.0 & 59.5\\
240 & 0 & 2 & 939.5 & 1275.0 & 0.067 & 0.0 & 0.0 & 0.0\\
240 & 10 & 2 & 910.5 & 1087.0 & 0.035 & 0.0 & 8.5 & 11.5\\
240 & 20 & 2 & 922.5 & 889.0 & -0.007 & 0.5 & 7.0 & 32.5\\
240 & 40 & 2 & 876.5 & 844.0 & -0.007 & 0.0 & 19.0 & 61.0\\
245 & 0 & 2 & 927.0 & 1319.5 & 0.079 & 0.0 & 0.0 & 0.0\\
245 & 10 & 2 & 911.0 & 941.0 & 0.006 & 0.0 & 5.0 & 15.0\\
245 & 20 & 2 & 891.0 & 816.5 & -0.015 & 0.0 & 3.5 & 36.5\\
245 & 40 & 2 & 857.5 & 722.0 & -0.027 & 0.0 & 7.5 & 72.5\\
\end{tabular}
}
\end{table}
\subsection{Occupied methane hydrates}
\label{sec:org0197cec}

A clathrate is characterised by which of its cages are filled.
The engine enumerates cages by their ring-size signature on the water graph (Section~\mref{sec:algorithms}) and assigns each guest by the nearest periodic centroid within a radius or by ray-parity inside the triangulated faces.
On the ideal GenIce methane hydrates the signature search returns exactly the \(16\) \(5^{12}\) and \(48\) \(5^{12}6^{2}\) cages of the \(2\times2\times2\) sI cell and the \(16\) \(5^{12}\) and \(8\) \(5^{12}6^{4}\) cages of the sII unit cell.
We set the centroid radii to \(4.0\), \(4.4\) and \(4.8\,\mathring{\mathrm{A}}\) for the \(5^{12}\), \(5^{12}6^{2}\) and \(5^{12}6^{4}\) cages.
When all cages of a structure are scored together, the radius is the largest among its cages, \(4.4\,\mathring{\mathrm{A}}\) for sI and \(4.8\,\mathring{\mathrm{A}}\) for sII (\texttt{scripts/hydrate\_occupancy.py}).
Under the centroid predicate every cage contains one methane, no cage contains two and no methane is free.
The largest methane-to-centroid distance is \(0.031\,\mathring{\mathrm{A}}\) in sI and \(1.225\,\mathring{\mathrm{A}}\) in sII (record \texttt{terra-2026-09-02}, \texttt{centre\_max}).
The ray-parity predicate of \texttt{seams cages -{}-inside} gives the same counts on both hydrates (record \texttt{hydrate-inside-2026-10-31}, \texttt{scripts/hydrate\_inside\_cli.py}).
The filled-cage count is the observable a hydrate nucleation run commits on, and it is a component of the PLUMED action next to the ion classes.
Figure~\ref{fig:app-hydrate} plots that count on the GenIce methane hydrates (record \texttt{terra-2026-09-02}, file \texttt{hydrate\_occupancy.txt}).
Every enumerated cage is occupied once.
Figure~\mref{fig:app-hydrate-box} of the main text shows one unit cell of each framework.


\begin{figure}[htbp]
\centering
\input{figures/app_hydrate}
\caption{Methane occupancy.
GenIce sI and sII.
Teal: occupied cages.
Coral: enumerated cages of that ring-size signature.
The two coincide.
The $2\times 2\times 2$ sI cell has $16$ $5^{12}$ and $48$ $5^{12}6^{2}$ cages, and sII has $16$ $5^{12}$ and $8$ $5^{12}6^{4}$ cages.}
\label{fig:app-hydrate}
\end{figure}
\subsection{The original application deposits}
\label{sec:orgfdf0af5}

The 2020 release ran five deposited trajectories \citep{goswami2020dseams}.
They are bulk cubic ice, homogeneous nucleation, a square ice nanotube, a confined monolayer and its two-dimensional radial distribution.
The reproducibility package (\url{https://github.com/d-SEAMS/dseams2_repro}) fetches those exact deposits and checksums them against figshare.
It reruns each analysis through the Python package, and the ring and cage analyses through the Lua library as well (Table~\ref{tab:deposits}).
The in-plane radial distribution is the two-dimensional \texttt{rdf\_2d} of \texttt{notebooks/05\_rdf2d.py}, with a bin of \(0.05\,\mathring{\mathrm{A}}\).
Three of the five ask questions a per-molecule ice label cannot answer.
We identify the nanotube by its four-membered prisms, the monolayer by its ring-size distribution and height coverage, and the confined liquid by its in-plane radial distribution.
Those are the confined-ice criteria of \citet{goswami2020nanoribbon}, and they run on the same primitive-ring network as the cage tests.

\begin{table}[htbp]
\caption{\label{tab:deposits}The five figshare deposits of the original release, rerun through d-SEAMS 2.0 (PydSEAMSlib; the first four also through \texttt{require("dseams")}).
The cubic ice, nanotube and monolayer rows are record \texttt{elja-1798666}, file \texttt{figshare-demos/figshare-demos.json}.
The nucleation and RDF rows are the rerun of Section~\ref{sec:s-rerun}, files \texttt{results/figshare-demos/figshare-demos.json}, \texttt{results/walks/figshare-nucleation.txt} and \texttt{results/notebooks/05\_rdf2d.ipynb}.
Frame is the one the original analysis used.}
\centering
\scriptsize
\setlength{\tabcolsep}{3pt}
\begin{tabular}{llrll}
deposit & trajectory & atoms & frame & result\\
\hline
11448720 & mW cubic ice & 4096 & 1 & 4096 DDC vertices, 0 HC; CHILL+ 4096 cubic\\
11448702 & mW nucleation & 4096 & 148 & 740 DDC, 632 HC molecules; \(n_{\max}\) 905\\
11448768 & square ice nanotube & 250 & 1 & 427 primitive rings; four-membered prisms only\\
11448741 & confined monolayer & 314 & 1 & 119 primitive rings; four-membered rings dominate\\
11448711 & confined liquid, RDF & 314 & 1 & 2D \(g(r)\) peak at \(2.825\,\mathring{\mathrm{A}}\)\\
\end{tabular}
\end{table}

\section{Reproduction}
\label{sec:s-repro}
The reproducibility package \citep{dseams2repro} is \url{https://github.com/d-SEAMS/dseams2_repro} at commit \texttt{\reprobase} (tag v2.7.0).
The records dated 31 October 2026 ran on that commit with the first revision of the package patch (sha256 \texttt{0e53bc8e}, tree \texttt{5dde84a5}).
The revision cited here (sha256 \texttt{276ec3f2}, tree \texttt{4a2f6b0e}) adds the PTM cutoff \(0.1\) and \texttt{scripts/ptm\_cutoff\_sweep.py}, fits \(q_6\) at the mW coexistence point, reads the dump column header, and writes the freud means with \texttt{scripts/sota\_freud.py}.
It rewrites the Lechner--Dellago curves of Figure~\mref{fig:sota} with \texttt{scripts/regen\_ld\_figure.py}.
The 0.1 PTM rows ran on it.
The patch adds the program-files archive, those records and the scripts that write the tables and figures from them.
The supplementary data hold the patch in \texttt{patches/} and the records in \texttt{repro-records/}.
With the patch, the pixi lock of the package pins OVITO, freud, GenIce2 and cycless next to the engine.
genice2-cage is installed with pip, and RINGS is built from its SourceForge archive.
Its source lock builds the engine as the program files do, from commit \texttt{305fb3c6} and the ten release patches, and refuses a patched tree other than \coretree.
Its \texttt{wire} step points the PydSEAMSlib and yodaStruct builds at that engine checkout.
Each record holds its \texttt{run.sh}, its raw output and a \texttt{provenance.txt} with the archive checksum and the machine.
Each \texttt{run.sh} refuses a source tree that differs from the archive.
The four timing records of the main text rerun with
\begin{lstlisting}[language=bash]
git clone https://github.com/d-SEAMS/dseams2_repro
cd dseams2_repro
git checkout f625883238bf0fb3b5b56b61a48994e4cb12e43e
# D: the supplementary data, with patches/ and repro-records/
git apply $D/patches/dseams2_repro-v2.7.0-fixes.patch
cp -r $D/repro-records/. results/reference/
A=$PWD/program-files/seams-core-v2.12.0.tar.gz
tar -xzf $A && S=$PWD/seams-core-v2.12.0
meson setup $S/build $S --buildtype=release \
  -Dwith_cli=true -Dwith_tests=true -Dwith_gpulite=disabled
meson compile -C $S/build
pip download --no-deps --no-binary pydseamslib \
  pydseamslib==2.11.1
tar -xzf pydseamslib-2.11.1.tar.gz && P=pydseamslib-2.11.1
rm -rf $P/subprojects/seams-core && tar -xzf $A -C $P/subprojects
mv $P/subprojects/seams-core-v2.12.0 $P/subprojects/seams-core
pip install --target $PWD/pyd ./$P
python scripts/figshare_demos.py fetch data/figshare
python scripts/public_ice.py fetch data/public-ice
R=results/reference
# baseline: git archive of d-SEAMS 1 commit 1c8efa38 in $B,
# built in $B/build
bash $R/bench-2026-10-31/run.sh $S $S/build $B $B/build 3 $A
BEFORE=1 bash $R/cli-carry-2026-10-31/run.sh $S $S/build 3 $A 5
export H2H_SEAMS=$S/build/seams PYDSEAMS=$PWD/pyd \
  H2H_SCRIPTS=$PWD/scripts ARCHIVE=$A
bash $R/h2h-2026-10-31/run.sh && \
  bash $R/h2h-2026-10-31/rings_mw.sh
pip install --target $PWD/lmp lammps==2025.7.22.4.0
export SEAMS_LMP=$PWD/lmp/lammps/lmp PYTHONPATH=$PWD/pyd
python scripts/md_heldout.py --out data/md-heldout \
  --record $R/md-heldout-2026-10-31
python scripts/md_strides.py --heldout data/md-heldout \
  --out data/md-strides --record $R/md-strides-2026-10-31
bash $R/incremental-md-2026-10-31/run.sh $S $S/build 3 $A
\end{lstlisting}
Every command after the \texttt{cp} runs inside \texttt{pixi shell}, which provides the compilers, Meson, Python and Snakemake.
Each timing script pins core 3 and \texttt{OMP\_NUM\_THREADS=1}.
The head-to-head scripts take OVITO, freud, NumPy and SciPy from \texttt{\$H2H/sotaenv}, GenIce2 and cycless from \texttt{\$H2H/env/gi} and RINGS from \texttt{\$H2H/rings}.
They read the Niu trajectory from \texttt{\$H2H/data} and the mW crystal from \texttt{input/traj} of the source tree, write under \texttt{\$H2H/out}, and list the versions in \texttt{raw/env.txt} of the record (Section~\mref{sec:repro}).
The \texttt{summary.py} of the bench and carry records, the \texttt{results.md} of the head-to-head record and the \texttt{summary.txt} of the incremental-md record give the numbers quoted here.
The asv suite of the churn sweep reruns from the same shell with
\begin{lstlisting}[language=bash]
pip install asv==0.6.4 asv_runner==0.3.2
SEAMS_NO_HQ=1 snakemake -s workflow/Snakefile \
  --configfile config/config.yaml --cores 1 \
  --forcerun asv_churn --config \
  churn_driver=$S/build/tests/bench_churn_sweep churn_cpu=3 \
  asv_machine=xeon-8vcpu-core3 -- \
  results/asv-churn/churn_table.tex \
  results/asv-churn/churn_summary.txt
\end{lstlisting}
which writes the table of the sweep and its summary under \texttt{results/\allowbreak asv-churn/}.
The documentation of seams-core carries the version pins and the benchmark tables at \url{https://docs.dseams.info/paper/index.html}.

\subsection{Rerun of the workflow}
\label{sec:s-rerun}
The values that cite this section come from one run of the whole workflow of the package.
It started from a fresh clone of commit \texttt{f625883} with the package patch of sha256 \texttt{87ecb706} applied, and ran
\begin{lstlisting}[language=bash]
pixi run test
pixi run reproduce
\end{lstlisting}
in a login environment that sets only \texttt{HOME}, \texttt{USER}, \texttt{LANG}, \texttt{TERM} and \texttt{PATH}.
Both exited 0, between 15:03 and 18:41~UTC on 31 October 2026, and inside the run the engine passed its 87 tests and its 5 MPI tests.
The machine is a KVM guest with 8 vCPUs, an Intel Xeon with one socket, eight cores and one thread per core, and it exposes no cpufreq governor (\texttt{results/conditions.txt} and \texttt{results/hardware.json} of the run).
\texttt{scripts/reproduce.sh} waited 90~s for the one-minute load to fall below 0.5, to 0.38.
One HyperQueue worker (v0.19.0) owned CPUs 0--7 and bound every task to its allotment with \texttt{taskset}, and the measurement DAG ran one job at a time.
Processes outside the run used \(2.31\%\) of the busy CPU seconds of the machine.
Over the measurement DAG no two consecutive samples of the once-a-minute load sampler lie more than 90~s apart, so the VM did not pause.
The production DAG ran at its smoke matrix (\texttt{PRODUCTION=smoke}, the default), and of its outputs we quote only the seeding totals.
We quote no timing of this run; every timing in this paper comes from the records of Table~\ref{tab:records-full}.

The lock of that run builds seams-core \texttt{305fb3c6} with four patches.
The first (sha256 \texttt{9b218dc1}) holds the ring, cage-affiliation, update, device and command-line changes in an earlier form than the release patches above.
The other three fix the parser of the strong-scaling columns, give the engine's benchmark script its Catch2 wrap, driver flags and a load settle (\texttt{cffbc780}), and set the version.
The run reproduces the 14 files of record \texttt{walks-2026-10-31} with no line differing.
Two earlier runs of rule \texttt{seed\_counts} on the same machine, with the same inputs, seeds and thread count, give the other two counts of each cell in Table~\ref{tab:ntot}.

\bibliographystyle{elsarticle-num-names}
\bibliography{dseams2}
\end{document}
